\RequirePackage{etex}
\documentclass[preprint,authoryear]{elsarticle}
\usepackage{times}
\usepackage[T1]{fontenc}
% T1 fonts will be used to generate the final print and online PDFs,
% so please use T1 fonts in your manuscript whenever possible.
% Other font encondings may result in incorrect characters.
%
\usepackage{graphicx}
\usepackage{etex}
% RECOMMENDED %%%%%%%%%%%%%%%%%%%%%%%%%%%%%%%%%%%%%%%%%%%%%%%%%%%
%

% to typeset URLs, URIs, and DOIs
\usepackage{url}
\def\UrlFont{\rmfamily}

\usepackage{amsmath}
\usepackage{amssymb}
\usepackage{amsthm}
\usepackage{amsfonts}
\usepackage{graphicx}
\usepackage{multirow}
\usepackage[ruled,linesnumbered,vlined]{algorithm2e}
\usepackage{booktabs}
\usepackage[switch]{lineno}
\usepackage{comment}
\usepackage{times}
\usepackage{xcolor}
\usepackage{soul}
\usepackage[utf8]{inputenc}
\usepackage{todonotes}
\newcommand{\carmine}[1]{\todo[color=blue!20]{{\bf Carmine:} #1}}
\newcommand{\group}{\mathit{group}}
\newcommand{\lits}{\mathit{lits}}
\newcommand{\mwh}{\mathit{mwh}}
\newcommand{\mwl}{\mathit{mwl}}
\newcommand{\smwh}{\mathit{smwh}}
\newcommand{\wh}{\mathit{wh}}
\newcommand{\ml}{\mathit{ml}}
\newcommand{\inc}{\mathit{inc}}
\newcommand{\mps}{\mathit{mps}}
\newcommand{\blw}{\mathit{blw}}
\newcommand{\ext}{\xrightarrow{\text{ext}}}
\newcommand{\G}{\mathbb{G}}
\newcommand{\atoms}{\mathit{atoms}}
\newcommand{\parts}{\mathit{parts}}
\newcommand{\p}{\mathit{part}}
\newcommand{\bnd}{\mathit{bnd}}
\newcommand{\rsn}{\mathit{rsn}}
\newtheorem{lemma}{Lemma}
\newcommand{\cupdot}{\,\dot{\cup}\,}
\newtheorem{corollary}{Corollary}
\newcommand{\wasp}{\textsc{wasp}\xspace}
\newcommand{\clingo}{\textsc{clingo}\xspace}
\newcommand{\amowaspij}{\textsc{amowasp-base-py}\xspace}
\newcommand{\amoclingoijpy}{\textsc{amoclingo-base-py}\xspace}
\newcommand{\amoclingoijc}{\textsc{amoclingo-base-c}\xspace}
\newcommand{\amoclingoinf}{\textsc{amoclingo-inf}\xspace}
\newcommand{\amoclingoinfr}{\textsc{amoclingo-inf-r}\xspace}
\newcommand{\amoclingoinfrhybrid}{\textsc{amoclingo-inf-r-h}\xspace}
\newcommand{\amoclingoinfrlazy}{\textsc{amoclingo-inf-r-l}\xspace}
\newcommand{\amoclingoinfmr}{\textsc{amoclingo-inf-mr}\xspace}
\newcommand{\amoclingoinfmrlazy}{\textsc{amoclingo-inf-mr-l}\xspace}
\newcommand{\amoclingoinfmrhybrid}{\textsc{amoclingo-inf-mr-h}\xspace}

\newcommand{\eowaspij}{\textsc{eowasp-base-py}\xspace}
\newcommand{\eoclingoijpy}{\textsc{eoclingo-base-py}\xspace}
\newcommand{\eoclingoijc}{\textsc{eoclingo-base-c}\xspace}
\newcommand{\eoclingoinf}{\textsc{eoclingo-inf}\xspace}
\newcommand{\eoclingoinfr}{\textsc{eoclingo-inf-r}\xspace}
\newcommand{\eoclingoinfrlazy}{\textsc{eoclingo-inf-r-l}\xspace}
\newcommand{\eoclingoinfrhybrid}{\textsc{eoclingo-inf-r-h}\xspace}
\newcommand{\eoclingoinfmr}{\textsc{eoclingo-inf-mr}\xspace}
\newcommand{\eoclingoinfmrlazy}{\textsc{eoclingo-inf-mr-l}\xspace}
\newcommand{\eoclingoinfmrhybrid}{\textsc{eoclingo-inf-mr-h}\xspace}


\definecolor{clingo}{HTML}{1E90FF}      % dark blue
\definecolor{amoclingo}{HTML}{000000}   % light blue

\definecolor{wasp}{HTML}{ff6347}        % dark green
\definecolor{amowasp}{HTML}{83ff83}     % light green


% the following package is optional:
\usepackage{latexsym}
\usepackage{xspace}
\usepackage{url}
% \usepackage{soul}
\usepackage{subcaption}
\usepackage{caption}


\usepackage{tikz}
\usepackage{pgfplots}
\usetikzlibrary{plotmarks}
\usepgfplotslibrary{groupplots}
\pgfplotsset{
    compat=1.18 % version for compiler
	, filter discard warning=false % remove notes while compiling
	, legend cell align=left
	, minor grid style={loosely dotted, lightgray}
	, major grid style={loosely dashed, lightgray}
}

\def\naf{\ensuremath{\raise.17ex\hbox{\ensuremath{\scriptstyle\mathtt{\sim}}}}\xspace}
\newtheorem{theorem}{Theorem}[section]
\newtheorem{proposition}{Proposition}[section]
\theoremstyle{definition}
\newtheorem{definition}{Definition}[section]
\newtheorem{example}{Example}[section]

\def\naf{\ensuremath{\raise.17ex\hbox{\ensuremath{\scriptstyle\mathtt{\sim}}}}\xspace}


\def\amosum{\ensuremath{\textsc{amosum}}\xspace}
\def\eosum{\ensuremath{\textsc{eosum}}\xspace}

\usepackage{listings}
\makeatletter
\lst@Key{countblanklines}{true}[t]{\lstKV@SetIf{#1}\lst@ifcountblanklines}
\lst@AddToHook{OnEmptyLine}{%
    \vspace{-0.7\baselineskip}%
    \lst@ifnumberblanklines\else%
    \lst@ifcountblanklines\else%
    \advance\c@lstnumber-\@ne\relax%
    \fi%
    \fi}
\makeatother
\lstset{
    mathescape,
    columns=fullflexible,
    basicstyle=\small\ttfamily,
    breaklines=true,
    keepspaces=true,
    %frame=lines,
    escapechar=\% \_,
    %numbers=left,
    %firstnumber=last,
    numberblanklines=false,
    countblanklines=false,
    numbersep=5pt,
    upquote=true,
    literate={~} {$\sim$}{1},
    commentstyle=\colour{gray},
    escapeinside={{/*}{*/}},
}


\definecolor{darkgreen}{rgb}{0.0, 0.2, 0.13}
\definecolor{darkorange}{rgb}{1.0, 0.55, 0.0}
\definecolor{darkviolet}{rgb}{0.58, 0.0, 0.83}




\journal{Artificial Intelligence}
\begin{document}
\begin{frontmatter}
\title{AMOSUM Constraints: Enhancing ASP Propagation with Aggregates and At-Most-One Constraints}
%
\author[1]{Mario Alviano}
\ead{mario.alviano@unical.it}
\author[1]{Carmine Dodaro\corref{cor1}}
\ead{carmine.dodaro@unical.it}
\author[2]{Thomas Eiter}
\ead{thomas.eiter@tuwien.ac.at}
\author[1]{Salvatore Fiorentino}
\ead{salvatore.fiorentino@unical.it}
\author[2]{Tobias Geibinger}
\ead{tobias.geibinger@tuwien.ac.at}
\author[1]{Marco Maratea}
\ead{marco.maratea@unical.it}

\cortext[cor1]{Corresponding author}
\affiliation[1]{
    organization={Department of Mathematics and Computer Science, University of Calabria},
    addressline={Via P. Bucci, Cubo 30B},
    postcode={87036},
    city={Rende (CS)},
    country={Italy}
}
\affiliation[2]{
    organization={Knowledge-based Systems Group, Institute of Logic and Computation, TU Wien},
    addressline={Favoritenstraße 9-11},
    postcode={1040},
    city={Vienna},
    country={Austria}
}

\begin{abstract}
% Aggregates such as sum and count are among the most frequently used linguistic extensions of Answer Set Programming (ASP).
% At-most-one (AMO) constraints are a specific form of aggregates that excludes the simultaneous truth of multiple elements in a set.
% This article unleashes a powerful propagation strategy in case groups of elements in an aggregate are also involved in AMO constraints.
% In fact, the combined knowledge given by aggregates and AMO constraints significantly increases the effectiveness of search space pruning, resulting in sensible performance gains.
% Answer Set Programming (ASP) is a popular problem-solving formalism with a rule-based language. While plain ASP is already very expressive, several language extensions have been developed to ease modeling or to improve the underlying solving.
% %
% Aggregates such as \textsc{sum} and \textsc{count} are among the most frequently used linguistic extensions of ASP.
% At-most-one (\textsc{amo}) constraints are a specific form of aggregate that excludes the simultaneous truth of multiple elements in a set.
% This article introduces \textsc{amosum} constraints, a novel construct that compactly combines \textsc{sum} and \textsc{amo} constraints and enables enhanced propagation in ASP solvers.
% By leveraging the combined inference potential of aggregates and exclusivity constraints, \textsc{amosum} allows the solver to deduce deterministic consequences that are otherwise unreachable when the constraints are treated separately.
% %
% We also provide a theoretical foundation for \textsc{amosum} constraints, introducing the notion of maximum possible sum to generalize propagation beyond previously studied cases.
% We formalize several inference rules, including propagation for false literals, and prove their soundness via a detailed analysis of reason structure and minimization. 
% In particular, we develop efficient algorithms to construct minimal reasons, improving clause learning and search space pruning.
% %
% Finally, we extend the ASP language with a native syntax for \textsc{amosum} and a complementary \textsc{eosum} construct for exactly-one constraints.
% At the ground level, these constructs correspond to structured pseudo-Boolean (PB) constraints, making the proposed approach relevant to native PB solving as well.
% The implementation is integrated into the ASP solvers \textsc{wasp} and \textsc{clingo}, with both Python- and C-based propagators, and supports lazy variants of propagation for improved efficiency.
% An extensive empirical evaluation on synthetic and combinatorial benchmarks demonstrates substantial performance gains, in some cases solving instances orders of magnitude faster than state-of-the-art solvers.
{\color{blue}
Answer Set Programming (ASP) is a widely adopted paradigm for knowledge representation and declarative problem solving, whose language has been extended with constructs that simplify modeling and improve solving performance. Among them, aggregates such as \textsc{sum} and \textsc{count} are extensively used, with at-most-one (\textsc{amo}) constraints representing a special case of \textsc{count} aggregates that allows at most one literal in a given set to be true.

This paper introduces \textsc{amosum}, a novel construct that combines \textsc{sum} and \textsc{amo} constraints within a unified propagator. By exploiting their interaction, \textsc{amosum} enables deterministic inferences that cannot be derived when the two constraints are handled independently. We establish its theoretical foundations by introducing the notion of \emph{maximum possible sum}, formalizing propagation rules, including rules for false literals, and developing efficient algorithms for constructing subset-minimal reasons.

We also extend the ASP language with a native syntax for \textsc{amosum} and the complementary \textsc{eosum} construct for exactly-one constraints. At the ground level, both constructs correspond to structured pseudo-Boolean constraints. The proposed approach is implemented in \textsc{wasp} and \textsc{clingo} and includes lazy propagation variants. An extensive empirical evaluation on synthetic and combinatorial benchmarks demonstrates substantial performance improvements over state-of-the-art solvers.
}
\end{abstract}

\begin{keyword}
Answer Set Programming\sep 
Aggregate Constraints\sep
At-Most-One Constraints\sep
Reason Minimization\sep
Lazy Propagation
\end{keyword}
\end{frontmatter}

%\clearpage
% \tableofcontents

% \clearpage
% New things:
% \begin{itemize}
%     \item Salvatore: New way to compute a reason: whenever the sum of other groups contributes to N and the bound is B, in this group the reason is all false that are above or equal B-N.
%     \item Technique: Full propagator, propagation cost in terms of steps
%     \item Technique: Inferring true literal (proof of correctness, computational cost or empirical experiments)
%     \item Computation of the reason: Improvement of reason's computation
%     \item Minimization of the reason (subset minimal)
%     \item Construct: KR linguistic extension
%     \item Implementation: Lazy
%     \item Implementation: Clingo C++ and Python API    
% \end{itemize}

% \clearpage

\section{Introduction}

Answer Set Programming (ASP) is a widely adopted formalism for knowledge representation and automated reasoning in the field of artificial intelligence \citep{MarekT99,Niemela99}.
In ASP, combinatorial problems are expressed using logic rules enriched with constructs designed to concisely capture complex constraints.
Solutions to these problems are provided as \emph{stable models} \citep{GelfondL90}, i.e., classical models satisfying an additional stability condition:
once the negative literals are fixed, producing the \emph{program reduct}, any deviation from the model (i.e., falsifying a true atom) violates at least one rule in the reduct.

In their simplest form, ASP programs comprise normal rules, each normal rule having a single head atom and a body being a conjunction of literals.
Intuitively, a head atom is required to be true whenever the associated body is true, so that stable models are classical models of the propositional knowledge base obtained by mapping normal rules to implications (body implies head).
The stability condition also enforces other properties of stable models, such as being supported models (every true atom is the head of some rule with true body)~\citep{DBLP:journals/mlcs/Fages94} and unfounded-free models (the support is acyclic)~\citep{DBLP:journals/tcs/Dung92}.
All such properties are used to achieve efficient computation by combining a \emph{conflict-driven clause learning} (CDCL) algorithm \citep{DBLP:journals/ia/Dodaro24} with \emph{propagators}.
CDCL is a modern form of non-chronological backtracking based on the pattern \textit{choose}-\textit{propagate}-\textit{learn}~\citep{DBLP:series/faia/0001LM21}:
a branching literal is heuristically \textit{chosen} and \emph{propagated} to infer deterministic consequences, where possible, and each inferred literal is associated with a reason being a clause (i.e., a set of literals, one of which is required to be true) among those associated with a propagator.
Reasons are used to \emph{learn} new clauses in case of conflict, so to effectively prune the search space.

To enhance expressiveness, ASP programs often incorporate \emph{aggregate constraints} \citep{DBLP:conf/aaaiss/BartholomewLM11,DBLP:journals/ai/FaberPL11,DBLP:journals/tocl/Ferraris11,DBLP:journals/tplp/GelfondZ14,DBLP:journals/ai/LiuPST10,DBLP:journals/ai/SimonsNS02}, particularly \textsc{sum} and \textcolor{blue}{\textsc{count} aggregates in rule bodies.}
In a \textsc{sum}, Boolean literals are associated with weights, and the sum of the weights of true literals is required to satisfy a given (in)equality.
\textcolor{blue}{A \textsc{count} aggregate is the special case in which all literals have unit weight, and therefore constrains the number of literals that are true.}

\begin{example}[Emergency Response Team Selection]\label{ex:emergency}
A city control room must dispatch support units to assist firefighters at a chemical incident.
The options are 
\emph{Drone team} ($x$), with operator $o_1$, to give situational awareness (contribution 1),
\emph{HazMat unit} ($y$), with operator $o_1$, to handle chemicals (contribution 2), and
\emph{Medical triage team} ($z$), with operator $o_2$, to handle casualties (contribution 2).
The incident is considered adequately covered if the dispatched support provides coverage score 3 or higher, that is, if the inequality
\begin{align}\label{eq:intro-example}
    1 \cdot x + 2 \cdot y + 2 \cdot z \geq 3
\end{align}  
is satisfied;
for instance, if $z = 1$ and either $x = 1$ or $y = 1$, or both $x = y = 1$.
In terms of ASP syntax, this scenario can be expressed by the following program:
\begin{align}
    \mathit{unit}(x,1,o_1). &\qquad \mathit{unit}(y,2,o_1).\qquad \mathit{unit}(z,2,o_2). \label{eq:emergency:facts}  \\
    \{\mathit{dispatch}(U)\}\ & \text{:- } \mathit{unit}(U,C,O). \label{eq:emergency:choice} \\
    & \text{:- } \#\mathit{sum}\{C,U : \mathit{dispatch}(U),\  \mathit{unit}(U,C,O)\} < 3. \label{eq:emergency:constraint}
\end{align}
Above, \eqref{eq:emergency:facts} are \emph{facts} representing the available units and their contribution,
\eqref{eq:emergency:choice} is a \emph{choice rule} selecting a set of units to dispatch, and
\eqref{eq:emergency:constraint} is a \emph{constraint} enforcing inequality \eqref{eq:intro-example}.
\end{example}
%
Constraints such as the one shown in Example~\ref{ex:emergency} are common in real-world problems and are handled by dedicated propagators in ASP solvers.
In this work, we assume \textsc{sum}s to be used as constraints, to have nonnegative weights and lower-bound enforcement, and target the computation of classical models satisfying them, making our techniques applicable across various stable model semantics \citep{DBLP:journals/tplp/AlvianoFG23}.

\textcolor{blue}{A particularly relevant special case of \textsc{count}, and hence also of \textsc{sum}, is the \emph{at-most-one} (\textsc{amo}) constraint~\citep{DBLP:conf/ictai/Surynek20}, which requires that at most one literal in a given set must be true.
In practice, both general \textsc{sum} constraints and \textsc{amo} constraints are handled by a specialized propagator~\citep{DBLP:conf/iclp/GebserKKS09,DBLP:journals/fuin/FaberLMR11}.}

%
\begin{example}[Continuing Example~\ref{ex:emergency}]\label{ex:emergency:amo}
If $z$ is assigned $0$ (i.e., the Medical triage team is not dispatched), then $x = y = 1$ is necessary to satisfy the constraint (i.e., both Drone team and HazMat unit must be dispatched).
\textcolor{blue}{However, the city has only one certified operator, who can either control the drone system or lead the HazMat deployment, but cannot perform both tasks simultaneously.
Accordingly, $x$ and $y$ are subject to an \textsc{amo} constraint, which can be expressed in ASP as follows:}
\begin{align*}
    & \text{:- } \#\mathit{sum}\{1,x : \mathit{dispatch}(x);\ 1,y : \mathit{dispatch}(y)\} > 1.
\end{align*}
Given the \textsc{amo} constraint, $x$ and $y$ cannot both be true, yielding a conflict.
From this, a solver might eventually deduce that $z$ must be true.
Importantly, this inference arises from the \emph{combined effect} of \textsc{sum} and \textsc{amo} constraints, yet typical propagators treat them in isolation, missing such deterministic consequences.
\end{example}

% This article introduces \emph{\textsc{amosum} constraints}, a novel constraint construct that tightly integrates \textsc{sum} and \textsc{amo} constraints into a unified formalism.
% Our key idea is that literals belonging to mutually exclusive \textsc{amo} groups contribute at most the \emph{maximum weight in their group} to a \textsc{sum}.
% This leads to stronger propagation than what is possible when the constraints are handled independently.
% Although this article focuses on ASP, the proposed propagation principles also apply to native pseudo-Boolean (PB) solving, since a ground \textsc{amosum} constraint corresponds to a weighted PB inequality equipped with a family of AMO relations.

This article introduces \emph{\textsc{amosum} constraints}, a novel constraint construct that tightly integrates \textsc{sum} and \textsc{amo} constraints into a unified formalism. 
Our key idea is that literals belonging to mutually exclusive \textsc{amo} groups contribute at most the \emph{maximum weight in their group} to a \textsc{sum}.
This leads to stronger propagation than what is possible when the constraints are handled independently.
\textcolor{blue}{Although our focus is on ASP, the same propagation principles apply to native pseudo-Boolean (PB) solving, since a ground \textsc{amosum} constraint corresponds to a structured PB system comprising a weighted inequality and a family of \textsc{amo} relations.}
%
An early version of \textsc{amosum} constraints was presented at the \emph{33rd International Joint Conference on Artificial Intelligence} (IJCAI; \citeauthor{DBLP:conf/ijcai/AlvianoDFM24}~\citeyear{DBLP:conf/ijcai/AlvianoDFM24}), where the focus was on a single inference rule and was implemented as a Python-based external propagator for \textsc{wasp}~\citep{DBLP:journals/tplp/DodaroR20}.  
The present article offers a broader and more mature framework: it refines the inference mechanism, introduces static and on-the-fly reason minimization, provides a dedicated ASP language construct for expressing \textsc{amosum} constraints declaratively, and adopts a more flexible integration strategy (in the form of a lazy propagator).  
As for the implementation, in addition to the one in \textsc{wasp}, the propagator has now been ported to \textsc{clingo}~\citep{DBLP:conf/iclp/GebserKKOSW16}, both through its C and Python external interfaces.  

Specifically, the present article addresses the following research questions:

\paragraph{Q1} \emph{How can inference and learning be improved in the solver?}
% We propose novel propagation rules for AMOSUM constraints based on the \emph{maximum possible sum} under current assignments.
% We further introduce mechanisms to infer both true and false literals and formally prove the soundness of their associated reasons.

\paragraph{Q2} \emph{Can such combinations be expressed as a new language construct?}
% We introduce the \texttt{\#AMOSUM} construct, extending the ASP language to directly encode SUM+AMO combinations with a clear and declarative syntax.

\paragraph{Q3} \emph{Which combinations of \textsc{sum} and \textsc{amo} constraints can be effectively propagated jointly?}
% We formally identify a class of SUM constraints partitioned into groups, each subject to an AMO constraint.
% This design avoids overestimation of the sum and enables stronger propagation.

\paragraph{Q4} \emph{Can reasons be minimized to improve learning and efficiency?}
% We define and analyze a polynomial-time algorithm for minimizing reasons based on a notion of \emph{incrementality}, and show how minimization improves learning and solver performance.

\paragraph{Q5} \emph{How does the approach generalize to exactly-one constraints and upper bounds?}
% We extend the language with a dual construct, \texttt{\#EOSUM}, that enforces \emph{exactly-one} semantics on each part, and generalize \textsc{amosum} to support upper-bound (\texttt{$\leq$}) variants.

\paragraph{Q6} \emph{How effective is the approach in practice?}
% We implement our propagator in both \textsc{wasp} and \textsc{clingo}, with Python and C interfaces, and propose \emph{lazy} variants of propagation that delay reasoning until strategically beneficial.
% Extensive experiments on synthetic and combinatorial benchmarks show that \textsc{amosum} achieves substantial performance gains, often solving instances that are otherwise intractable using state-of-the-art solvers.
\medskip

In summary, this article contributes both new language constructs and formal reasoning techniques that significantly enhance the propagation power of ASP solvers when dealing with aggregates and mutual exclusion.
The main contributions of this work can be summarized as follows:
\begin{itemize}
\item
We introduce \emph{\textsc{amosum} constraints}, a new ASP construct that jointly encodes \textsc{sum} and at-most-one (\textsc{amo}) constraints, enabling stronger and earlier propagation than traditional approaches.

\item 
We extend the ASP language with \texttt{\#AMOSUM} and \texttt{\#EOSUM} constructs, supporting both at-most-one and exactly-one variants, with both $\geq$ and $\leq$ semantics.
For example, the constraint  \eqref{eq:emergency:constraint} in Example~\ref{ex:emergency} can be encoded as:
\begin{align*}
%    \label{ex:emergency:amosum}
    \amosum\{C : \textcolor{blue}{dispatch(X), \;unit(X,C,O) [O]}\} \;\geq\; 3.
\end{align*}

In this case, the annotation $[O]$ specifies that two units sharing the same operator cannot simultaneously be true in an answer set, thereby enforcing operator uniqueness.

Now suppose that instead of requiring each operator to be assigned to \textit{at most one} unit, we impose the stricter condition that each operator must be assigned to \textit{exactly one} unit.  
In this case, the constraint \eqref{eq:emergency:constraint} can be formulated as:
\begin{align*}
%    \label{ex:emergency:eosum}
    \eosum\{C : \textcolor{blue}{dispatch(X), \;unit(X,C,O) [O]}\} \;\geq\; 3.
\end{align*}

\item 
We define \emph{new inference rules} for \textsc{amosum} constraints based on the \emph{maximum possible sum}, allowing inference of both true and false literals beyond classical \textsc{sum} and \textsc{amo} propagators. For each rule, we also establish the soundness of its corresponding reason.
\item 
We design a \emph{reason minimization algorithm} and formally prove its correctness. The algorithm removes unessential literals in polynomial time and produces subset-minimal reasons, improving clause learning and solver performance. 
\item 
We implement the proposed techniques in the solver \textsc{clingo}, and we introduce \emph{lazy propagation strategies} that dynamically defer inference when advantageous.
\textcolor{blue}{Furthermore, we also implement the \textsc{eosum} counterpart with the $\ge$ constraint, and for both \textsc{amosum} and \textsc{eosum} we implement the $\le$ variant.}
\item 
Finally, we present an extensive \emph{experimental evaluation} of the proposed methods. The results show that they lead to substantial performance gains across multiple benchmarks, including synthetic and real-world combinatorial problems.
\end{itemize}

\noindent{\em Organization}. The remainder of this paper is organized as follows. 
Section~\ref{sec:prel} introduces the necessary preliminaries, including the syntax and semantics of the underlying formalism, a treatment of clauses and at-most-one (\textsc{amo}) constraints as special cases, and a discussion of stable model search with propagators.
Section~\ref{sec:amosum} defines the proposed \textsc{amosum} constraints, elaborating on their syntax, semantics, propagation mechanisms, and key properties.
In Section~\ref{sec:minimization}, we address the problem of reason minimization and introduce a novel on-the-fly minimization technique.
Section~\ref{sec:language-extension} discusses the extension of the ASP language to support the new constructs, along with implementation details such as meta-encoding strategies, propagation control, and lazy evaluation. 
Section~\ref{sec:exp} presents our experimental evaluation, describing benchmark settings and reporting on the performance of the proposed approach.
Section~\ref{sec:rw} reviews related work, and Section~\ref{sec:conclusion} concludes the paper with final remarks and directions for future research.







% \#amo\_sum {
%           W,C: col(node1,C), colour\_weight(C,W) [node1]; 
%           W,C: col(node2,C), colour\_weight(C,W) [node2];
%           W,C: col(node3,C), colour\_weight(C,W) [node3]
%           } >= B, bound(B).




\section{Preliminaries}
\label{sec:prel}
This section reports syntax and semantics of normal programs extended with \textsc{sum} constraints (Section~\ref{sec:syntax}).
Clauses and \textsc{amo} constraints are then given as special cases of \textsc{sum} (Section~\ref{sec:amo}).
After that, the stable model search procedure implemented by modern ASP solvers is illustrated with a focus on the concept of propagator (Section~\ref{sec:propagators}).

\subsection{Syntax and Semantics}\label{sec:syntax}
Let $\mathcal{A}$ be a set of propositional \emph{atoms}.
A \emph{literal} is an atom possibly preceded by the \emph{default negation} symbol \naf.
Given a literal $\ell$, $\overline{\ell}$ denotes the \emph{complement} of $\ell$, i.e., $\overline{p} = \naf p$ and $\overline{\naf p} = p$ for every $p \in \mathcal{A}$;
for two sets $L,X$ of literals, we let $\overline{L}=\{\overline{\ell} \mid \ell \in L\}$ and
$L \cupdot X = (L \setminus \overline{X}) \cup X$.
% ; abusing of notation if 
%  $X= \{\ell\}$ then $L\cup \dot{\ell}$ is written in place of $L \cup \dot{\{\ell\}}$.

A \emph{(normal) rule} has the form
\begin{align}\label{eq:rule}
    p \leftarrow \ell_1, \ldots, \ell_n
\end{align}
where $n \geq 0$, $p$ is an atom, and $\ell_1,\ldots,\ell_n$ are literals.
A \emph{\textsc{sum} constraint} (or simply {\em constraint}) has the form
\begin{align}\label{eq:sum}
    \textsc{sum}\{w_1 : \ell_1;\ \cdots;\ w_n : \ell_n\} \geq b
\end{align}
where $n \geq 0$, $\ell_1,\ldots,\ell_n$ are distinct literals such that $\ell_i \neq \overline{\ell_j}$ (for all $1 \leq i < j \leq n$),
and $b,w_1,\ldots,w_n$ are nonnegative integers.
We remark  that in the ASP-Core-2 standard~\citep{DBLP:journals/tplp/CalimeriFGIKKLM20}, the constraint (\ref{eq:sum}) is written as the headless rule \linebreak \lstinline|:- #sum{$w_1,\ell_1$ : $\ell_1$; ...; $w_n,\ell_n$ : $\ell_n$} < $b$|. 
A \emph{program} is a set of rules and constraints.

For convenience, we introduce the following piece of syntactic sugar, which we refer to as a choice rule:
\begin{align}
\label{eq:choice-rule}
\{p\} \leftarrow \ell_1, \ldots, \ell_n 
\end{align}
This notation stands for the following set of rules:
\begin{align}
\label{eq:choice-rule-rewriting}
& \textcolor{blue}{p} \leftarrow \ell_1, \ldots, \ell_n, \textcolor{blue}{\naf p'} \\
& \textcolor{blue}{p'} \leftarrow \ell_1, \ldots, \ell_n, \textcolor{blue}{\naf p}
\end{align}
where $p'$ is a fresh atom.

\begin{example}[Running example]\label{ex:running}
Let $\Pi_{\mathit{run}}$ be the following:
\begin{align*}
\begin{array}{rll}
    r_{\alpha\phantom{'}}: & \alpha\phantom{'} \leftarrow \naf \alpha' & \alpha \in \{x,y,z\}\\
    r_{\alpha'}: & \alpha' \leftarrow \naf \alpha & \alpha \in \{x,y,z\}\\  % MALVI: fixed \alpha' after review
    \sigma_{1}: & \textsc{sum}\{
        1 : \textcolor{blue}{\naf x};\ 1 : \textcolor{blue}{\naf y}\phantom{;\ 2 : z}
    \} \geq 1 \\
    \sigma_{2}: & \textsc{sum}\{
        1 : \phantom{\naf}x;\ 2 : \phantom{\naf}y;\ 2 : z
    \} \geq 3
    % r_{1,i}: & a_i \leftarrow \naf b_i & i = 1..3\\
    % r_{2,i}: & b_i \leftarrow \naf a_i & i = 1..3\\
    % \sigma_{1,i}: & \textsc{sum}\{
    %     1 : \naf a_i;\ 1 : \naf b_i
    % \} \geq 1 & i = 1..3 \\
    % \sigma_1: & \textsc{sum}\{
    %     1 : a_1;\ 3 : b_1; \\
    %   & \phantom{\textsc{sum}\{}
    %     1 : a_2;\ 3 : b_2; \\
    %   & \phantom{\textsc{sum}\{}
    %     1 : a_3;\ 3 : b_3
    % \} \geq 7
\end{array}
\end{align*}
Note that there are six atoms ($x,y,z,x',y',z'$), six rules and two constraints.
\end{example}

We adopt the following notation.
%
Atoms, rules, and constraints occurring in a program $\Pi$ are denoted as $\mathit{atoms}(\Pi)$, $\mathit{rules}(\Pi)$, and $\mathit{constraints}(\Pi)$, respectively.
%
For a rule $r$ of the form \eqref{eq:rule}, $H(r) := p$ (the head atom of $r$), 
$B(r) := \{\ell_1, \ldots, \ell_n\}$ (the body literals of $r$),
$B^+(r) := B(r) \cap \mathcal{A}$ (the atoms occurring in the positive body of $r$), and
$B^-(r) := \overline{B(r)  \setminus \mathcal{A}}$ (the atoms occurring in the negative body of $r$).
%
For a constraint $\sigma$ of the form \eqref{eq:sum}, let $\mathit{bnd}_\sigma := b$ (the bound of $\sigma$ is $b$),
$\mathit{wh}_\sigma(\ell_i) := w_i$ (the weight of $\ell_i$ in $\sigma$ is $w_i$; for all $i \in [1..n]$), and
$\mathit{lits}_\sigma := \{\ell_1, \ldots, \ell_n\}$ (the literals occurring in $\sigma$ are $\ell_1, \ldots, \ell_n$).
If $\sigma$ is clear from the context, 
%it is possibly omitted 
we may omit it from the above notation.
Finally, we define 
the relation ${}\in{}$ as $(w_i : \ell_i) \in \sigma$ for $i \in [1..n]$, to be read as $(w_i : \ell_i)$ is an element in (the aggregation set of) $\sigma$.


\begin{example}[Continuing Example~\ref{ex:running}]\label{ex:notation}
Rule $r_x$ is such that $H(r_x) = x$, $B^+(r_x) = \emptyset$ and $B^-(r_x) = \{x'\}$.
Regarding $\sigma_2$, 
$\mathit{bnd}_{\sigma_2} = 3$,
$\mathit{wh}(y) = 2$,
$y \in \mathit{lits}_{\sigma_2}$, and
$(2 : y) \in \sigma_2$.
\end{example}

An \emph{assignment} $I$ is an ordered set (or list) of literals such that $I \cap \overline{I} = \emptyset$;  % MALVI: added (or list) after review
literals in $I$ are true, literals in $\overline{I}$ are false, and all other literals are undefined.
By abuse of notation, let $I(\ell) := 1$ if $\ell \in I$ and $0$ otherwise (true literals assigned to 1);
let $I^\uparrow(\ell) := 1$ if $\ell \notin \overline{I}$ and $0$ otherwise (true and undefined literals assigned to 1).
Moreover, in the subsequent sections, we will iteratively construct an assignment by using a list and occasionally consider its \emph{prefixes} (i.e., assignments obtained in previous iterations).
%
An assignment $I$ is a \emph{(total) interpretation} of a program $\Pi$ if $I \cup \overline{I} = \mathit{atoms}(\Pi) \cup \overline{\mathit{atoms}(\Pi)}$.
%
For total interpretations, the satisfaction relation $\models$ (\emph{is model of}) is inductively defined as follows:
\begin{itemize}
    \item for a literal $\ell$, $I \models \ell$ if $\ell \in I$;
    \item for a rule $r$ of the form (\ref{eq:rule}), $I \models B(r)$ if $I \models \ell$ for all $\ell \in B(r)$;
    \item for a rule $r$ of the form (\ref{eq:rule}), $I \models r$ if $I \models H(r)$ whenever $I \models B(r)$;
    \item for a constraint $\sigma$ of the form (\ref{eq:sum}), $I \models \sigma$ if 
$\sum_{i = 1}^{n}{w_i \cdot I(\ell_i)} \geq \mathit{bnd}_\sigma$; and 
    \item for a program $\Pi$, $I \models \Pi$ if $I \models r$ for all $r \in \mathit{rules}(\Pi)$ and $I \models \sigma$ for all $\sigma \in \mathit{constraints}(\Pi)$.
\end{itemize}
% for a literal $\ell$, $I \models \ell$ if $\ell \in I$;
% for a rule $r$ of the form (\ref{eq:rule}), $I \models \mathit{imp}(R,\ell)$ if $I \models \ell$ for all $\ell \in \mathit{imp}(R,\ell)$, and $I \models r$ if $I \models H(r)$ whenever $I \models \mathit{imp}(R,\ell)$;
% for a constraint $\sigma$ of the form (\ref{eq:sum}), $I \models \sigma$ if 
% $\sum_{i = 1}^{n}{w_i \cdot I(\ell_i)} \geq \mathit{bnd}_\sigma$;
% for a program $\Pi$, $I \models \Pi$ if $I \models r$ for all $r \in \mathit{rules}(\Pi)$ and $I \models \sigma$ for all $\sigma \in \mathit{constraints}(\Pi)$.
%
The \emph{reduct} $\Pi^I$ of a program $\Pi$ with respect to an interpretation $I$ is $\{H(r) \leftarrow B^+(r) \mid r \in \mathit{rules}(\Pi),$ $I \models \textcolor{blue}{B(r)} \}$.
(Note that $\mathit{constraints}(\Pi^I) = \emptyset$.)
%
An interpretation $I$ is a \emph{stable model} of a program $\Pi$ if $I \models \Pi$ and there is no interpretation $J$ such that \textcolor{blue}{$J \cap \atoms(\Pi^I) \subset I \cap \atoms(\Pi^I)$} and $J \models \Pi^I$ \citep{DBLP:journals/ai/FaberPL11}.
Let $\mathit{SM}(\Pi)$ denote the set of stable models of $\Pi$.

\begin{example}[Continuing Example~\ref{ex:running}]\label{ex:semantics}
We first observe that, in every stable model of $\Pi_\mathit{run}$, $\alpha'$ is true whenever $\alpha$ is false, and vice versa ($\forall \alpha \in \{x,y,z\}$).
$\mathit{SM}(\Pi_\mathit{run})$ contains 
\textcolor{blue}{$\{x, z\}$},
\textcolor{blue}{$\{y, z\}$} (negative literals and \emph{prime} atoms are omitted for simplicity).
\end{example}


\subsection{Clauses and AMO as Special Cases}\label{sec:amo}

Given a set $\{\ell_1, \ldots, \ell_n\}$ of $n \geq 0$ literals, an \emph{at-least-one (\textsc{alo})} constraint enforces truth of at least one literal in the set.
It is also commonly referred to as \emph{clause}, and can be expressed as the following \textsc{sum} constraint:
\begin{align}\label{eq:clause}
    \textsc{sum}\{1 : \ell_1;\ \cdots;\ 1 : \ell_n\} \geq 1
\end{align}
Clause \eqref{eq:clause} is also written (if not ambiguous) as
\begin{align}
    \label{eq:clauses}
    \{\ell_1, \ldots, \ell_n\}
\end{align}

In the following, given a program $\Pi$ and a clause $C$, we write $\Pi \models C$ if for each interpretation $I$, $I \models \Pi$ implies $I \models C$.
Note that for each clause $C \in \Pi$, $\Pi \models C$.

A clause $C \in \Pi$ is unessential \textcolor{blue}{with respect to $\Pi$} if and only if it is entailed \textcolor{blue}{by the remaining program, that is, if $\Pi\setminus\{C\}\models C$.}
% In other words, a clause is unessential if $SM(\Pi) = SM(\Pi \cup \{C\})$.

Given two clauses $C,D$ of the form \eqref{eq:clauses} and a literal 
$\ell$ such that $\ell \in C$ and $\overline{\ell} \in D$ then the \textit{resolution}
\citep{DBLP:journals/jacm/DavisP60,DBLP:journals/jacm/Robinson65}
step of $C$ and $D$ upon $\ell$, written $C \otimes_{\ell} D$, corresponds to 
$(C \setminus \{\ell\}) \cup (D \setminus \{\overline{\ell}\})$.
%Intuitively, if $C \in \Pi$ and $D \in \Pi$ then since an interpretation cannot simultaneously satisfy $\ell$ and $\overline{\ell}$ then $C \setminus \{\ell\}$ or $D \setminus \{ \overline{\ell}\}$ have to be satisfied, thus the following clause  $(C \setminus \{\ell\}) \cup (D \setminus \overline{\ell})$ has to be satisfied.
Note that if an interpretation $I$ fulfills both $I \models C$ and $I \models D$, then $I \models C \otimes_{\ell} D$ holds.

Given a set $\{\ell_1, \ldots, \ell_n\}$ of $n \geq 1$ literals, an \textsc{AMO} constraint inhibits truth of pairs of literals in the set.
It can be expressed as the following \textsc{sum} constraint:
\begin{align}\label{eq:amo}
    \textsc{sum}\{1 : \overline{\ell_1};\ \cdots;\ 1 : \overline{\ell_n}\} \geq n - 1
\end{align}
Essentially, it equivalently asks for falsifying all but possibly one literal in the set.
In fact, given an interpretation $I$, $I \models \eqref{eq:amo}$ if 
$\sum_{i = 1}^{n}{I(\overline{\ell_i})} \geq n - 1$, or equivalently
$\sum_{i = 1}^{n}{I(\ell_i)} \leq 1$.
The \textsc{AMO} constraint \eqref{eq:amo} is compactly written $\textsc{amo}\{\ell_1,\ldots,\ell_n\}$.

\begin{example}[Continuing Example~\ref{ex:running}]
Note that $\sigma_1$ is the clause $\{\overline{x}, \overline{y}\}$, or also the AMO constraint $\textsc{amo}\{x, y\}$.
\end{example}


\subsection{Stable Model Search and Propagators}\label{sec:propagators}
Stable model search is implemented in modern ASP solvers using a \emph{conflict-driven clause learning} (CDCL) algorithm \citep{DBLP:journals/ai/GebserKS12}.
This algorithm relies on the pattern \textit{choose}-\textit{propagate}-\textit{learn}.
In a nutshell, the approach involves incrementally constructing a stable model, starting with an empty list $I$ (representing the empty assignment).
%
During each computational step, a branching literal is heuristically \textit{chosen} for addition to $I$.
Subsequently, it is \emph{propagated} to introduce new deterministic consequences to $I$, where possible. 
Each deterministic consequence $\ell$ incorporated into $I$ is associated with a \emph{reason} being the clause comprising $\ell$ and the false literals leading to the inclusion of $\ell$ in $I$.
Moreover, for each reason $R$, $\Pi \models R$.

% After a literal $\ell$ is inferred, thanks to some inference rule, under interpretation I, then a \textit{reason}, written $\mathit{reason}(z)$, 
% is specified. $\mathit{reason}(\ell)$ defines \textit{why} $\ell$ has been propagated.\\
For a literal $\ell$ and a reason $R$, the set $\mathit{imp}(R, \ell) = \overline{R\setminus\{\ell\}}$ is called implicant of $\ell$, and $\ell$ is implied by $R$.
% Intuitively, whenever $I \models \Pi$, then the interpretation $I$ must satisfy $\bigwedge_{\ell' \in \mathit{imp}(R,\ell)} \ell' \rightarrow \ell$.
\textcolor{blue}{Intuitively, if all literals in $\mathit{imp}(R,\ell)$ are true in a model of $\Pi$, then $\ell$ must also be true.}
% Equivalently, if all literals in $\mathit{imp}(R,\ell)$ are true in a model of $\Pi$, then $\ell$ must also be true.
% that is, $\Pi \models \bigwedge_{\ell' \in \mathit{imp}(R,\ell)} \ell' \rightarrow \ell$.

%
A \emph{conflict} arises if the complement of a deterministic consequence is already in $I$, and its analysis leads to \emph{learn} new clauses via \emph{(backward) resolution} starting from the reasons of the conflicting literals. 
(Recall that resolution combines two clauses 
$\{p\} \cup L$ and $\{\overline{p}\} \cup L'$, written $(\{p\} \cup L)\otimes_p(\{\overline{p}\} \cup L'),$ to obtain a new clause $L \cup L'$, where $p$ is an atom, and $L,L'$ are sets of literals.)
During backward resolution, previously added literals are removed from $I$ until the learned clause results in the inclusion of \textcolor{blue}{only one literal, which is either the last branching literal or one of its deterministic consequences. At this point}, a new deterministic consequence is \textcolor{blue}{inferred}, driving the search into a different branch.
%
This iterative process continues until either $I$ represents a stable model or the empty clause is learned. The latter means that the input program does not admit stable models.

A \emph{propagator} is a module designed to compute deterministic consequences of a given assignment.
The most basic among these modules is \emph{unit propagation}:
a literal $\ell$ is added to $I$ if there is a clause $\{\ell\} \cup L$ such that $L \subseteq \overline{I}$, that is, if $\ell$ belongs to a clause that can be satisfied only by extending $I$ with $\ell$.
In this case, $\mathit{reason}(\ell)$ is defined as $\{\ell\} \cup L$.
%
Modern ASP solvers enrich the input program with clauses enforcing $I$ to be a model of rules of the form \eqref{eq:rule} (i.e., $\{p, \overline{\ell_1}, \ldots, \overline{\ell_n}\}$), as well as clauses enforcing other required properties of stable models (which are out of the scope of this article).


\begin{example}\label{ex:unit-propagation}
Let $\Pi$ have, among others, the rules
\begin{align*}
    \begin{array}{ccc}
        x \leftarrow \naf z \qquad &
        y \leftarrow \naf z \qquad &
        w \leftarrow x, y
    \end{array}
\end{align*}
Hence, a modern ASP solver materializes the clauses
\begin{align*}
    \begin{array}{ccc}
        \{x,z\} \qquad &
        \{y,z\} \qquad &
        \{w,\overline{x},\overline{y}\}
    \end{array}
\end{align*}
If the current assignment $I$ is the list $[\overline{w}]$ and $\overline{z}$ is selected as the branching literal, $I$ is updated to $[\overline{w}, \overline{z}]$; unit propagation infers $x,y$ from the first two clauses and then $\overline{y}$ (a conflict) from the third clause.
Hence, we have $I = [\overline{w}, \overline{z}, x, y, \overline{y}]$,
$\mathit{reason}(x) = \{x,z\}$,
$\mathit{reason}(y) = \{y,z\}$, and
$\mathit{reason}(\overline{y}) = \{w,\overline{x},\overline{y}\}$.
By resolving $\mathit{reason}(y)$ and $\mathit{reason}(\overline{y})$, we obtain $\{w,\overline{x},z\}$, which still contains more than one atom \textcolor{blue}{corresponding either to the last branching literal or one of its deterministic consequences (namely, $x$ and $z$).}
By resolving $\{w,\overline{x},z\}$ and $\mathit{reason}(x)$, we obtain $\{w,z\}$, which contains only one atom inferred from the branching point (namely, $z$).
CDCL therefore continues with $I = [\overline{w}, z]$ and $\mathit{reason}(z) = \{w,z\}$.
% Let $r_1$ be the rule $a_1 \leftarrow b_1, \naf c_1$, let $r_2$ be the rule $d_1 \leftarrow b_1, \naf c_1, \naf a_1, \naf d_1$, and assume $I=\{b_1, \overline{c_1}, \overline{d_1}\}$.
% Then, unit propagation on $r_1$ adds $a_1$ to $I$, and $\mathit{reason}(a_1) = \{a_1, \overline{b_1}, c_1\}$.
% Moreover, unit propagation on $r_2$ adds $\overline{a_1}$ to $I$, and $\mathit{reason}(\overline{a_1}) = \{\overline{a_1}, \overline{b_1}, c_1, \overline{d_1}\}$.
% Thus, since $I$ contains both $a_1$ and $\overline{a_1}$, the learning step applies backward resolution on their reasons and learn the clause $\{\overline{b_1}, c_1, \overline{d_1}\}$.
\end{example}

For a \textsc{sum} constraint $\sigma$ of form \eqref{eq:sum}, solvers typically employ a specific \textit{aggregate propagator}~\citep{DBLP:conf/iclp/GebserKKS09,DBLP:journals/fuin/FaberLMR11},
%, reported here as Algorithm~\ref{alg:sumpropagation}.
which essentially adds to $I$ the literal $\ell_i$ ($i \in [1..n]$) if $\ell_i$ is required to (possibly) reach the bound $b$, i.e., if
\begin{align}\label{eq:sum:propagation}
    \sum_{j \in [1..n], j \neq i}{w_j \cdot J^\uparrow(\ell_j)} < b
\end{align}
where $J$ is the first prefix of $I$ meeting the condition (\ref{eq:sum:propagation}).
In this case, $\mathit{reason}(\ell_i)$ is 
\begin{align}
\label{eq:sum:reason}
    \{\ell_i\} \cup (\mathit{lits}_\sigma \cap \overline{J})
\end{align} that is, the false literals occurring in $\sigma$ are enforcing truth of $\ell_i$.
%
In the special case of \eqref{eq:amo}, that is, if $\sigma$ is $\textsc{amo}\{\ell_1, \ldots, \ell_n\}$, the literal $\overline{\ell_i}$ ($i \in [1..n]$) is added to $I$ if there is $\ell_j$, with $j \neq i$, such that $\ell_j \in I$.
In this case, $\mathit{reason}(\overline{\ell_i})$ is $\{\overline{\ell_i}, \overline{\ell_j}\}$.
%
\textcolor{blue}{The propagation rules can be informally understood as a consistency-preserving mechanism: a literal is inferred when it cannot be falsified without violating the constraint, given the current partial assignment and AMO partitioning.}

\begin{example}[Continuing Example~\ref{ex:running}] \label{ex:sumpropagator}
If $I$ is empty, no literal can be inferred from $\sigma_1$ and $\sigma_2$.
If $I$ is $[\overline{z}]$, then the application of \eqref{eq:sum:propagation} to the literals of $\sigma_2$ gives
\begin{align*}
    2 \cdot [\overline{z}]^\uparrow(y) + 2 \cdot [\overline{z}]^\uparrow(z) = 2 \cdot 1 + 2 \cdot 0 = 2 < 3\\
    1 \cdot [\overline{z}]^\uparrow(x) + 2 \cdot [\overline{z}]^\uparrow(z) = 1\cdot 1 + 2 \cdot 0 = 1 < 3\\
    1 \cdot [\overline{z}]^\uparrow(x) + 2 \cdot [\overline{z}]^\uparrow(y) = 1 \cdot 1 + 2 \cdot 1 = 3 \not< 3
\end{align*}
Hence, $x$ and $y$ are inferred with
$\mathit{reason}(x) = \{x,z\}$ and
$\mathit{reason}(y) = \{y,z\}$.
Note that, once $I = [\overline{z},x,y]$, the application of \eqref{eq:sum:propagation} to $\sigma_1$ gives
\begin{align*}
    1 \cdot [\overline{z},x,y]^\uparrow(\overline{y}) = 1 \cdot 0 = 0 < 1\\
    1 \cdot [\overline{z},x,y]^\uparrow(\overline{x}) = 1 \cdot 0 = 0 < 1
\end{align*}
For example, a conflict is raised when $\overline{y}$ (or, symmetrically, $\overline{x}$) is added to $I$ with $\mathit{reason}(\overline{y}) = \{\overline{x}, \overline{y}\}$.

% Assume Algorithm~\ref{alg:sumpropagation} is called with $r=r_4$, $I=\{a_2,\naf b_2\}$, and $\ell=b_2$.
% As first step, $\mathit{mps}$ is updated (we recall it was set to 12) by removing the weight associated to $b_2$, that is $\mathit{mps}(r_4)$ is 9.
% Then, the algorithm iterates over the undefined literals in $\mathit{as}(r)$, that are $\{a_1,b_1,a_3,b_3\}$.
% The first iteration considers $a_1$, $\mathit{mps}(r_4)-\mathit{wh}(r_4, a_1)\geq \mathit{bnd}(r_4)$, therefore no literals are added to $S$.
% The second iteration considers $b_1$, $\mathit{mps}(r_4)-\mathit{wh}(r_4, b_1) < \mathit{bnd}(r_4)$, therefore $S$ is extended by adding ($b_1$, \{$\naf b_2$\}).
% The third iteration considers $a_3$, $\mathit{mps}(r_4)-\mathit{wh}(r_4, a_3)\geq \mathit{bnd}(r_4)$, therefore no literals are added to $S$.
% The second iteration considers $b_3$, $\mathit{mps}(r_4)-\mathit{wh}(r_4, b_3) < \mathit{bnd}(r_4)$, therefore $S$ is extended by adding ($b_3$, \{$\naf b_2$\}).
% The algorithm then terminates returning $S$.
\end{example}




\section{{\sc amosum} Constraints}

\label{sec:amosum}

In this section, \textsc{sum} constraints of the form \eqref{eq:sum} are replaced by the more general \textsc{amosum} constraints (Section~\ref{sec:amo:syntax}).
An \textsc{amosum} constraint combines one \textsc{sum} constraint with a set of \textsc{amo} constraints, in a convenient syntax.
We provide inference rules and their reasons (Section~\ref{sec:amo:propagator}), yielding a more effective propagator, as formally shown in Section~\ref{sec:amo:properties}. \textcolor{blue}{We then prove the correctness of the reasons in Section~\ref{sec:correctness}.}  


\subsection{Syntax and Semantics}\label{sec:amo:syntax}

An \emph{\textsc{amosum} constraint} (or simply constraint from this point) $\sigma$ has the form
\begin{align}\label{eq:amosum}
    \amosum\{w_1 : \ell_1\ [s_1];\ \cdots;\ w_n : \ell_n\ [s_n]\} \geq b
\end{align}
where $n \geq 0$, $\ell_1,\ldots,\ell_n$ are distinct literals such that $\ell_i \neq \overline{\ell_j}$ for all $1 \leq i < j \leq n$,
and $b,w_1,\ldots,w_n,s_1,\ldots,s_n$ are nonnegative integers.
%
The notation introduced in Section~\ref{sec:syntax} is further extended as follows:
$\mathit{parts}_\sigma := \{s_1, \ldots, s_n\}$ (the literals in $\sigma$ are partitioned into at most $n$ parts $s_1,\ldots,s_n$),
$\mathit{part}_\sigma(\ell_i) := s_i$ (the part of $\ell_i$ in $\sigma$ is $s_i$, for all $i \in [1..n]$), with the symmetric convention $\mathit{part}_\sigma(\overline{\ell_i})=\mathit{part}_\sigma(\ell_i)$.
Moreover, relation ${}\in{}$ is now defined as $(w_i : \ell_i\ [s_i]) \in \sigma$ for all $i \in [1..n]$ and \textcolor{blue}{$\mathit{lits}_\sigma|_s := \{\ell \mid (w:\ell [s]) \in \sigma\}$} (the literals in $\sigma$ belonging to the part $s$).
%
Regarding the relation \emph{is model of}, for $\sigma$ of the form \eqref{eq:amosum}, $I \models \sigma$ if 
$\sum_{i = 1}^{n}{w_i \cdot I(\ell_i)} \geq \mathit{bnd}_\sigma$, and 
%$\sum_{(w : \ell\ [s]) \in \sigma}{w \cdot I(\ell)} \geq \mathit{bnd}_\sigma$, and 
$\sum_{\ell \in \mathit{lits}_\sigma|_s}{I(\ell)} \leq 1$ for all $s \in \mathit{parts}_\sigma$.
Essentially, the definition for \eqref{eq:sum} is extended by enforcing an AMO constraint on each part of $\sigma$.

\begin{example}[Continuing Example~\ref{ex:running}]\label{ex:running-revised}
$\Pi_{\mathit{run}}$ is rewritten by replacing $\sigma_1$ and $\sigma_2$ with
\begin{align*}
\begin{array}{rl}
    \sigma_{3}: & \amosum\{
        1 : x\ [1];\ 2 : y\ [1];\ 2 : z\ [2]
    \} \geq 3.
\end{array}
\end{align*}
Note that $\mathit{parts}_{\sigma_3}=\{1, 2\}$, 
$\mathit{part}_{\sigma_{3}}(x) = \mathit{part}_{\sigma_{3}}(y) = 1$, 
$\mathit{part}_{\sigma_{3}}(z) = 2$, 
$\mathit{lits}_{\sigma_{3}}|_1 = \{x, y\}$, and
$\mathit{lits}_{\sigma_{3}}|_2 = \{z\}$.

\end{example}

\subsection{Propagation}\label{sec:amo:propagator}
A constraint $\sigma$ of the form \eqref{eq:amosum} is associated with three different inference rules; for each inference rule we assume to start with an assignment $I$ that is consistent with the \textit{at most one} constraint, that is, $\sum_{\ell \in \mathit{lits}_\sigma|_s}{I(\ell)} \leq 1$ for all $s \in \mathit{parts}_{\sigma}$.

\paragraph{AMO Rule}
First of all, the \textsc{amo} inference given by the parts of $\sigma$:
\begin{itemize}
    \item[\textbf{(RA)}] This inference rule adds the literal $\overline{\ell}$ to $I$ if there is $\ell' \in \mathit{lits}_\sigma|_{\mathit{part}_\sigma(\ell)}$ such that $\ell' \in I$.
In this case, $\mathit{reason}(\overline{\ell})$ is $\{\overline{\ell}, \overline{\ell'}\}$.

\end{itemize}
%

\paragraph{SUM Rules}
To describe the \textsc{sum} related inference rules we introduce the notion of max possible sum.
Given an assignment $I$ and a constraint $\sigma$, its \emph{maximum possible sum} $\mathit{mps}_\sigma(I)$
is the maximum possible value of 
\begin{align*}
    \sum_{i = 1}^{n}{w_i \cdot I'(\ell_i)}
\end{align*} 
such that 
\begin{align*}
    \sum_{\ell \in \mathit{lits}_\sigma|_s}{I'(\ell)} \leq 1 \quad \forall  s \in \mathit{parts}_{\sigma} 
\end{align*}
and $I'$ is a total interpretation such that $I' \supseteq I$. 
In other words, it is the maximum possible sum reachable starting from $I$ that satisfies all the at most one constraints for $\sigma$.
Given a literal $\ell$, the maximum possible sum reachable assuming $\ell$ true is
\begin{align}
    \label{eq:mps}
    mps_\sigma(I, \ell) := \sum_{s \in \parts_\sigma} \mathit{mwh}_\sigma(I \cupdot \{ \ell \},s)
\end{align}
where 
\begin{align*}
    \mathit{mwh}_\sigma(I,s) := \left\{
        \begin{array}{ll}
            \wh_\sigma(\ell') & \text{if } \lits_\sigma|_s \cap I = \{\ell'\} \\
            \max\{w \cdot I^\uparrow(\ell') \mid (w : \ell'\ [s]) \in \sigma\} & \text{otherwise.}
        \end{array}
    \right\}.
\end{align*}
Above, the intuitive reading of the second case is ``the maximum weight that part $s$ can contribute to the overall sum.''
Note that $\mps_\sigma(I)$ can be equivalently defined as follows:
\begin{align}
    \label{eq:mps2}
    mps_\sigma(I) := \sum_{s \in \parts_\sigma} \mathit{mwh}_\sigma(I,s).
\end{align}
Now we can proceed to explain the two inference rules:
\begin{itemize}
    \item[\textcolor{blue}{\textbf{(RT)}}] This inference rule is analogous to the one provided by the aggregate propagator, i.e., a literal is inferred true if it is needed to reach the bound: 
    
    \textcolor{blue}{An undefined literal $\ell$ is added to $I$ if \textcolor{blue}{$\ell \in \lits_{\sigma}$} and  $J$ is the first prefix of $I$ such that
    \begin{align}\label{eq:amosum:propagation:true}
        \mps_{\sigma}(J,\overline{\ell}) < \bnd_\sigma\,.
    \end{align}
    In this case, $\mathit{reason}(\ell)$ is
    \begin{align}\label{eq:amosum:reason:rt}
    \{\ell\} \cup \bigcup_{s \in \mathit{parts}_\sigma}{
        \mathit{rsn}_\sigma(J,\ell,s)
    }
    \end{align}}
    where 
    $\mathit{rsn}_\sigma(J,\ell,s) := \{\overline{\ell'}\}$ if $\mathit{lits}_\sigma|_s \cap J = \{\ell'\}$ (i.e., a true literal in $\lits_\sigma|_s$), and
    $ \{\ell' \in \lits_{\sigma}|_s \cap \overline{J} \mid \wh(\ell') \textcolor{blue}{>} \mwh(J \cupdot \{ {\overline{\ell}} \}, s)\}$ otherwise (i.e., the false literals in the part of $\ell$ that could have increased the overall sum).
    \item[\textcolor{blue}{\textbf{(RF)}}] This inference rule has no counterpart in \textsc{amo} or \textsc{sum} constraints, and \emph{enforces falsity of literals in a part} whose weight is guaranteed to be insufficient to reach the bound: 
    
    The literal $\overline{\ell}$ is added to $I$ if \textcolor{blue}{$\ell \in \lits|_\sigma$, $\lits_{\sigma}|_s \cap J = \emptyset$}, \textcolor{blue}{where $s=\mathit{part}_\sigma(\ell)$}, and $J$ is the first prefix of $I$ such that
    \textcolor{blue}{\begin{align}\label{eq:amosum:propagation:false}
        \mps_{\sigma}(J,\ell) < \bnd_\sigma\, .
    \end{align}}
    %a \textcolor{red}{
    % the literal $\overline{\ell}$ is added to $I$ if $\ell \in \lits_{\sigma}$ and
    % \begin{align}\label{eq:amosum:propagation:false:new}
    %     \mps(I,\ell) < \bnd_\sigma
    % \end{align}
    % }
    In this case, $\mathit{reason}(\overline{\ell})$ is
    \begin{align}\label{eq:amosum:reason:rf}
    \{\overline{\ell}\} \cup \bigcup_{s \in \mathit{parts}_\sigma \setminus \{\mathit{part}_\sigma(\ell)\}}{
        \mathit{rsn}_\sigma(J,\ell,s)
    }.
    \end{align}
\end{itemize}


% As outlined in the introduction, ASP programs can contain multiple rules of the form (\ref{eq:aggr}). Within these rules, there are cases when literals in the aggregates are inherently subject to an additional condition. Specifically, it becomes possible to devise a partition of the literals in the aggregate set, where literals within the same partition are constrained by an at-most-one (AMO) condition.
% In instances where such a condition is present, the aggregate propagators detailed in the previous section exhibit inherent inefficiencies.
% The objective of this section is to introduce a dedicated AMO-aware aggregate for such cases.

% For a rule $r$ of the form (\ref{eq:aggr}), let $\mathit{partitions}(r)$ be a partition of $\mathit{as}(r)$ subject to the AMO condition; for a literal $\ell$ and a rule $r$, let $\mathit{partition}(\ell, r)=\{P \mid P \in \mathit{partitions}(r), \ell \in P\}$ be the partition of $r$ containing $\ell$.


\begin{example}\label{ex:amosum:propagation:more-2}
Let us consider the following constraint:
\begin{align*}
\begin{array}{rl}
    \sigma_{4}: & \amosum\{
        1 : x\ [1]; 2 : y\ [1]; 2 : z\ [2]; 3 : w\ [2]
    \} \geq 3.
\end{array}
\end{align*}
If $I = [\overline{w}]$, the second inference rule \textcolor{blue}{(RT)} associated with $\sigma_4$ infers $z$ with $\mathit{reason}(z) = \{z, w\}$; the classical sum inference rule~\eqref{eq:sum:propagation} would not propagate in this case.
%
The same holds if $I = [\overline{x}, \overline{w}]$, with the addition of $y$ with $\mathit{reason}(y) = \{y, x, w\}$;
note that $w$ is included in $\mathit{reason}(y)$ because it could increase the overall sum.
%
On the other hand, if $I = [\overline{y}, \overline{w}]$, then $x$ and $z$ are inferred with 
$\mathit{reason}(x) = \{x, y, w\}$ and
$\mathit{reason}(z) = \{z, w, y\}$;
note that $y$ is included in $\mathit{reason}(z)$ because it could increase the overall sum (in this specific case $y$ could be ignored, in Section~\ref{sec:minimization} we see how to detect such case).
Finally, if $I = [x, \overline{y}, \overline{w}]$, then $z$ is inferred with $\mathit{reason}(z) = \{z, w, \overline{x}\}$;
in this case, $\overline{x}$ is included because it enforces a value for part 1 (again, in this specific case $\overline{x}$ could be ignored).
\end{example}

\begin{example}\label{ex:amosum:propagation:3}
Let us consider the following constraint:
\begin{align*}
\begin{array}{rl}
    \sigma_{5}: & \amosum\{
        1 : x\ [1];\, 2 : y\ [1];\, 2 : z\ [2];\, 2 : w\ [2]
    \} \geq 4.
\end{array}
\end{align*}
\textcolor{blue}{For $I = [\,]$, the inference rule (RF) of $\sigma_5$ infers $\overline{x}$ thanks to the application of \eqref{eq:amosum:propagation:false}:}
\begin{align*}
    \mps([\,], x) =  1 + \max\{2 \cdot [\,]^\uparrow(z), 2 \cdot [\,]^\uparrow(w)\} = 1 + 2 = 3 < 4.
\end{align*}
In this case, $\mathit{reason}(\overline{x}) = \{\overline{x}\}$.
\end{example}

\subsection{Properties}\label{sec:amo:properties}

Semantically, it is clear that \textsc{amosum} constraints do not extend the capabilities of the language.
% In fact, an \textsc{amosum} $\sigma$ constraint of the form \eqref{eq:amosum} can be expressed in terms of (no more than $n+1$) \textsc{sum} constraints (by design) as follows:
% \begin{align*}
%     \textsc{sum}\{w_1 : \ell_1;\ \cdots;\ w_n : \ell_n\} \geq b\\
%     \textsc{amo}\{\mathit{lits}_\sigma|_{s_i}\}\phantom{{}\geq b} & \phantom{xxx} \forall i \in [1..n]
% \end{align*}

\begin{proposition}
Let $\sigma$ be an \textsc{amosum} constraint of the form
\begin{equation}
\label{eq:amosum:prop}
    \textsc{amosum}\{w_1 : \ell_1\ [s_1];\ \dots;\ w_n : \ell_n\ [s_n]\} \;\geq\; b.
\end{equation}
Then $\sigma$ can be equivalently expressed using at most $n+1$ standard 
\textsc{sum}  constraints as follows:
\begin{align*}
    \textsc{sum}\{w_1 : \ell_1;\ \cdots;\ w_n : \ell_n\} \geq b\\
    \textsc{amo}\{\mathit{lits}_\sigma|_{s_i}\}\phantom{{}\geq b} & \phantom{xxx} \forall i \in [1..n]
\end{align*}

\end{proposition}


On the other hand, it is also true that \textsc{amosum} constraints immediately subsume \textsc{sum} (and \textsc{amo}) constraints.

\begin{theorem}
Every \textsc{sum} constraint can be expressed by an \textsc{amosum} constraint.
\end{theorem}
\begin{proof}
Let $\sigma$ be a \textsc{sum} constraint of the form \eqref{eq:sum}.
Construct the following \textsc{amosum} constraint $\sigma'$:
\begin{align*}
    \amosum\{w_1 : \ell_1\ [1];\ \cdots;\ w_n : \ell_n\ [n]\} \geq b.
\end{align*}
Note that each part $s \in \mathit{parts}_{\sigma'}$ is a singleton, and therefore trivially satisfies \begin{align*}
\sum_{\ell \in \mathit{lits}_{\sigma'}|_s}{I(\ell)} \leq 1.
\end{align*}
Hence, $I \models \sigma$ if and only if $I \models \sigma'$, for every interpretation $I$.
\end{proof}

The real advantage gained by extending the language with \textsc{amosum} constraints is in the associated inference rules, in particular the rules (RT) and (RF), while (RA) does not add anything to what can be inferred by \textsc{amo} constraints.
%It can be shown that the second inference rule 
We next show that (RT) alone is sufficient to capture all deterministic consequences that can be identified by \textsc{sum} constraints, and actually can result in a larger set of inferred literals.

\begin{theorem}\label{thm:inferences}
Let $\sigma$ be the \textsc{amosum} constraint of the form \eqref{eq:amosum}, 
$\sigma'$ be the \textsc{sum} constraint of the form \eqref{eq:sum} obtained from $\sigma$ by keeping the same literals, weights and bound, and  $I$ be an assignment.
Let $L$ be the set of literals inferred by rule (RT) associated with $\sigma$, and let $L'$ be the set of literals inferred by $\sigma'$.
Then $L \supseteq L'$ holds, while $L' \supseteq L$ is not guaranteed.
\end{theorem}
\begin{proof}
Let $\ell_i$ be inferred by $\sigma'$, and let $J$ be the first prefix of $I$ that triggers this inference. 
By definition, \eqref{eq:sum:propagation} holds.
Hence
\begin{align*}
   \textcolor{blue}{\sum_{s \in \parts}{\mathit{mwh}_\sigma(J \cupdot \{\overline{\ell_i}\},s)}} \le {} &
    \sum_{s \in \parts}{\max_{(w : \ell\ [s]) \in \sigma,\, \ell \not= \ell_i} w \cdot J^\uparrow(\ell)} \\    
    {}\leq{} & \sum_{j \in [1..n], j \neq i}{w_j \cdot J^\uparrow(\ell_j)} < b
\end{align*}
and \eqref{eq:amosum:propagation:true} holds as well, proving $L \supseteq L'$.
A  witness of $L' \not\supseteq L$ is given in Example~\ref{ex:amosum:propagation:more-2}.
\end{proof}

\textcolor{blue}{Theorem~\ref{thm:inferences} allows us to answer the first part of question \emph{Q1}, showing that inference rule (RT) is able to derive strictly more literals than the classical inference rule~\eqref{eq:sum:propagation}.
Moreover, it is straightforward to observe that the corresponding reason~\eqref{eq:amosum:reason:rt} is a subset of the classical reason~\eqref{eq:sum:reason}, since~\eqref{eq:amosum:reason:rt} contains a subset of the false literals of $\sigma$. Since a smaller reason usually speeds up answer set computation, this answers the second part of question \emph{Q1}.}
Furthermore Rule (RF) identifies deterministic consequences that are not computed by \textsc{amo} and \textsc{sum} constraints considered separately. 
\textcolor{blue}{In particular, \textsc{amo} constraints provide only the usual mutual-exclusion inferences, while standard \textsc{sum} propagation infers literals that are required to reach the lower bound.
By contrast, rule (RF) can infer that a literal must be false when selecting it would make the bound unreachable, even though neither the \textsc{amo} nor the \textsc{sum} constraint alone would derive this consequence.
Example~\ref{ex:amosum:propagation:3} illustrates this behavior.
%
Although domain consistency is traditionally formulated in the constraint-programming literature, we use the following constraint-level notion to characterize the propagation strength of \textsc{amosum}.
Intuitively, every currently undefined literal must admit both possible truth assignments in some satisfying extension, unless one of these assignments is detected to be impossible by propagation.}

{\color{blue}
\begin{definition}[Constraint-level domain consistency]
Let $\sigma$ be an \textsc{amosum} constraint and let $I$ be an assignment consistent with the \textsc{amo} constraints induced by the parts of $\sigma$, and such that $\mps_\sigma(I) \geq \bnd_\sigma$.
We say that $I$ is domain consistent with respect to $\sigma$ if, for every literal $\ell \in \lits_\sigma$ that is undefined in $I$, there exist total interpretations $K_t$ and $K_f$ such that
$$
K_t \supseteq I \cup \{\ell\},
\qquad
K_f \supseteq I \cup \{\overline{\ell}\},
$$
and
$$
K_t \models \sigma,
\qquad
K_f \models \sigma.
$$
Thus, every currently undefined literal has support for both truth values in an extension satisfying the \textsc{amosum} constraint.
\end{definition}


% \begin{proposition}[Constraint-level domain consistency]
% Let $\sigma$ be an \textsc{amosum} constraint of the form~\eqref{eq:amosum}, and let $I$ be an assignment that is consistent with the \textsc{amo} constraints induced by the parts of $\sigma$. Suppose that $\mps_\sigma(I) \geq \bnd_\sigma$ and that $I$ is closed under the exhaustive application of inference rules $(RA)$, $(RT)$, and $(RF)$. Then, $I$ is domain consistent with respect to $\sigma$. More precisely, for every literal $\ell \in \lits_\sigma$ that is undefined in $I$, there exist total interpretations $K_t$ and $K_f$ such that $K_t \supseteq I \cup \{\ell\}$, $K_f \supseteq I \cup \{\overline{\ell}\}$, $K_t \models \sigma$, and $K_f \models \sigma$.
% \end{proposition}
\begin{proposition}[Constraint-level domain consistency]
Let $\sigma$ be an \textsc{amosum} constraint of the form~\eqref{eq:amosum}, and let $I$ be an assignment consistent with the \textsc{amo} constraints induced by the parts of $\sigma$, and such that $\mps_\sigma(I) \geq \bnd_\sigma$.
If $I$ is closed under exhaustive application of inference rules $(RA)$, $(RT)$, and $(RF)$, then $I$ is domain consistent with respect to $\sigma$.
\end{proposition}

\begin{proof}
Let $\ell \in \lits_\sigma$ be undefined in $I$. Suppose first, towards a contradiction, that there is no total interpretation $K_f$ such that $K_f \supseteq I \cup \{\overline{\ell}\}$ and $K_f \models \sigma$.

By the definition of the maximum possible sum, this implies $\mps_\sigma(I,\overline{\ell}) < \bnd_\sigma$.
Therefore, there exists a first prefix of $I$ satisfying condition~\eqref{eq:amosum:propagation:true}, and inference rule $(RT)$ infers $\ell$. This contradicts the assumptions that $I$ is closed under $(RT)$ and that $\ell$ is undefined in $I$. Hence, $I \cup \{\overline{\ell}\}$ admits a total extension satisfying $\sigma$.

Suppose now, towards a contradiction, that there is no total interpretation $K_t$ such that $K_t \supseteq I \cup \{\ell\}$ and $K_t \models \sigma$.
Let $s = \mathit{part}_\sigma(\ell)$. If there is a literal $\ell' \in \lits_\sigma|_s$ such that $\ell' \in I$, then inference rule $(RA)$ infers $\overline{\ell}$. This contradicts the assumptions that $I$ is closed under $(RA)$ and that $\ell$ is undefined in $I$.

Otherwise, $\lits_\sigma|_s \cap I = \emptyset$. Since $I \cup \{\ell\}$ admits no total extension satisfying $\sigma$, the definition of the maximum possible sum yields $\mps_\sigma(I,\ell) < \bnd_\sigma$.

Therefore, there exists a first prefix of $I$ satisfying condition~\eqref{eq:amosum:propagation:false}, and inference rule $(RF)$ infers $\overline{\ell}$. This contradicts the assumptions that $I$ is closed under $(RF)$ and that $\ell$ is undefined in $I$.

Consequently, both $I \cup \{\ell\}$ and $I \cup \{\overline{\ell}\}$ admit a total extension satisfying $\sigma$. Thus, every value in the domain of every undefined literal has a support with respect to $\sigma$, and $I$ is domain consistent with respect to $\sigma$.
\end{proof}
n
\noindent
This proposition establishes a local property of $\sigma$. The supporting interpretations $K_t$ and $K_f$ are not required to satisfy the remaining rules and constraints of the program. Moreover, the result concerns the exhaustive application of the eager propagation rules; the lazy variants introduced in Section~\ref{sub:sec:laziness} intentionally trade propagation completeness for reduced propagation overhead.

}


\subsection{Correctness}

\label{sec:correctness}
We now present a series of theorems that demonstrate the correctness of the reasons defined above.
The first theorem shows that, given an assignment $I$ that triggers inference rule \textcolor{blue}{$(RT)$} for a literal $\ell$, the maximal achievable sum realized by $I$ is captured by the assignment $\mathit{imp}(R,\ell) \subseteq I$, which assigns all literals in $\mathit{imp}(R,\ell) $ to true, where $R$ is the associated reason of the form \textcolor{blue}{\eqref{eq:amosum:reason:rt}.}

\begin{theorem}
    \label{thm:mps:reason:eq:rt}
    Let $\sigma$ be an \textsc{amosum} constraint of the form~\eqref{eq:amosum}, and let $I$ be an assignment such that $\sum_{\ell \in \mathit{lits}_\sigma|_s}{I(\ell)} \leq 1$ for all $s \in \mathit{parts}_{\sigma}$. Let $\ell$ be the literal inferred by  \textcolor{blue}{\eqref{eq:amosum:propagation:true}} of the inference rule (RT), $R$ be the associated reason of the form \eqref{eq:amosum:reason:rt},  and $J$ be the first prefix of $I$ that triggered \textcolor{blue}{(RT).} 
    Then, $\mps_\sigma(J,\overline{\ell})=\mps_\sigma(\mathit{imp}(R,\ell),\overline{\ell})$. 
\end{theorem}
\begin{proof}
    Let $s \in \parts_\sigma$.
    On the one hand, if $\lits_\sigma|_s \cap J = \{\ell'\}$, then by construction $\ell' \in \mathit{imp}(R,\ell)$ and $\lits_\sigma|_s \cap \mathit{imp}(R,\ell) = \{\ell'\}$.
    That is, $\mwh_\sigma(J,s)=\mwh_\sigma(\mathit{imp}(R,\ell),s)$.
    \textcolor{blue}{On the other hand, suppose that $\lits_\sigma|_s \cap J = \emptyset$, and let 
    $$
     q=\mwh_\sigma(J\cupdot\{\overline{\ell}\},s)
     \quad\text{and}\quad
     F=\{\overline{\ell'}\mid \ell'\in\lits_\sigma|_s\cap(\overline{J} \cup \{\ell\}),\ \wh(\ell')>q\}.
    $$
    By construction, $(\mathit{imp}(R,\ell)\cupdot \{\overline{\ell}\})|_s=F$. Hence, every literal of part $s$ whose weight is greater than $q$ is false in $\mathit{imp}(R,\ell)\cupdot\{\overline{\ell}\}$, and therefore
    $$
    \mwh_\sigma(\mathit{imp}(R,\ell)\cupdot\{\overline{\ell}\},s)\le q.
    $$
    If $q=0$, equality follows immediately because weights are nonnegative. 
    If $q>0$, by the definition of $q$ there is a literal $\ml\in\lits_\sigma|_s$ of weight $q$ that is not false in $J\cupdot\{\overline{\ell}\}$. Since the reason contains only false literals of weight strictly greater than $q$, $\ml$ remains possible in $\mathit{imp}(R,\ell)\cupdot\{\overline{\ell}\}$. 
    Thus the reverse inequality holds as well.}
    Hence, $\mwh_\sigma(\mathit{imp}(R,\ell) \cupdot \{\overline{\ell}\},s) = \mwh_\sigma(J \cupdot \{ {\overline{\ell}} \},s)$, and $\mps_\sigma(\mathit{imp}(R,\ell),\overline{\ell}) = \mps_\sigma(J, \overline{\ell})$.
\end{proof}

Given this theorem, it is easy to show that the same reasoning can apply with the 
inference rule \textcolor{blue}{$(RF)$} and the associated reason of the form \textcolor{blue}{\eqref{eq:amosum:reason:rf}.}
\begin{theorem}
    \label{thm:mps:reason:eq:rf}
    Let $\sigma$ be an \textsc{amosum} constraint of the form \eqref{eq:amosum}, and let $I$ be an assignment such that $\sum_{\ell' \in \mathit{lits}_\sigma|_s}{I(\ell')} \leq 1$ for all $s \in \mathit{parts}_{\sigma} $. Let $\overline{\ell}$ be the literal inferred by \textcolor{blue}{\eqref{eq:amosum:propagation:false}} of inference rule \textcolor{blue}{$(RF)$}, $R$ be the associated reason of the form \eqref{eq:amosum:reason:rf}, and $J$ be the first prefix of $I$ that triggered the inference rule.
    Then, $\mps_\sigma(J,\ell)=\mps_\sigma(\mathit{imp}(R,\overline{\ell}),\ell)$. 
\end{theorem}
\begin{proof}
    A reason of the form \textcolor{blue}{\eqref{eq:amosum:reason:rt}} differs from a reason of the form \textcolor{blue}{\eqref{eq:amosum:reason:rf}} just for the application of $\rsn_\sigma$ on part $\mathit{part}(\ell)$.
    By Theorem~\ref{thm:mps:reason:eq:rt} if $s \in \parts$ is different from $\mathit{part}(\ell)$  it follows that
    $\mwh(J \cupdot \{\ell\}, s) = \mwh(\mathit{imp}(R,\overline{\ell}) \cupdot \{\ell\}, s)$.
    Thus, it is sufficient to prove that $\mwh(\mathit{imp}(R,\overline{\ell})\cupdot \{\ell\},\mathit{part}(\ell)) = \mwh(J \cupdot \{ {\ell} \},\mathit{part}(\ell))$ holds.
   \textcolor{blue}{In (RF), $\ell \in \lits_\sigma$ and $\lits_s \cap J = \emptyset$, where $s=\mathit{part}(\ell)$. Moreover, the syntax of~\eqref{eq:amosum} excludes complementary pairs, so $\overline{\ell}\notin\lits_\sigma$, and the (RF) reason contains no literal from part $s$. Therefore, after setting $\ell$ true, it is the unique true literal of that part, and
    $\mwh(J \cupdot \{ {\ell} \},\mathit{part}(\ell)) = \wh(\ell)$, and  $\mwh(\mathit{imp}(R,\overline{\ell}) \cupdot \{\ell\}, \mathit{part}(\ell))= \wh(\ell)$.}
    Hence, $\mwh(J \cupdot \{ {\ell} \},\mathit{part}(\ell)) = \mwh(\mathit{imp}(R,\overline{\ell})\cupdot \{\ell\}, \mathit{part}(\ell))= \wh(\ell)$
\end{proof}

An important property that is instrumental for our discussion is the following.
{\color{blue}
\begin{lemma}
\label{lemma:maximal_subset:antimonote}
    Suppose $I$ and $I'$ are assignments such that $I' \subseteq I$, and let $\ell$ be a literal. Let $K' = I' \cupdot \{\ell\}$ and $K = I \cupdot \{\ell\}$. Suppose that $\sum_{\ell' \in \lits_s} K'(\ell') \le 1$ and $\sum_{\ell' \in \lits_s} K(\ell') \le 1$ for all $s \in \parts$. The $\mps$ function is antimonotone with respect to subset inclusion in its first argument, that is,
    $$
    \mps(I',\ell) \ge \mps(I,\ell).
    $$
    In other words, a smaller assignment (that is, one with more undefined literals) yields a greater $\mps$ value.
\end{lemma}
\begin{proof}
As $I' \subseteq I$, we have $K' \subseteq K$. By assumption, $K$ satisfies all AMO constraints; consequently, so does $K'$.
We prove that $\mwh(K',p) \geq \mwh(K,p)$ for every part $p$.

First, suppose that $K' \cap \lits_p = \{t\}$. Since $K' \subseteq K$, the literal $t$ is also true in $K$. Moreover, AMO consistency ensures that $t$ is the unique true literal of $p$ in both assignments. Therefore, $\mwh(K',p)=\mwh(K,p)=\wh(t)$.

Suppose now that $K'$ contains no true literal from $p$, while $K \cap \lits_p=\{t\}$. The literal $t$ cannot be false in $K'$: otherwise, $\overline{t}\in K'\subseteq K$, contradicting $t\in K$. 
Hence, $K'^\uparrow(t)=1$, and thus
$$
\mwh(K',p) = \max_{t'\in\lits_p} K'^\uparrow(t')\wh(t') \geq \wh(t) = \mwh(K,p).
$$

Finally, suppose that neither $K'$ nor $K$ contains a true literal from $p$. Since $K'\subseteq K$, for every $t'\in\lits_p$ we have $K'^\uparrow(t')\geq K^\uparrow(t')$.
Taking the maximum over the literals of $p$ therefore gives $\mwh(K',p)\geq\mwh(K,p)$.

The inequality holds for every part. Summing over all parts yields 
$$\mps(I',\ell) = \sum_{p\in\parts}\mwh(K',p) \geq \sum_{p\in\parts}\mwh(K,p) = \mps(I,\ell),$$ 
which proves the claim.
\end{proof}
}


What remains to prove is that $R$ is a valid reason, that is $\Pi \models R$.
\begin{theorem}
 Suppose $\Pi$ is a program containing an \textsc{amosum} constraint $\sigma$ of the form \eqref{eq:amosum}, and $I$ is an assignment such that $\sum_{\ell \in \mathit{lits}_\sigma|_s}{I(\ell)} \leq 1$ for every $s \in \mathit{parts}_{\sigma} $. Let $\ell$ be the literal inferred by \textcolor{blue}{\eqref{eq:amosum:propagation:true} of (RT)}, $R$ be the associated reason  of the form \textcolor{blue}{\eqref{eq:amosum:reason:rt}},  and $J$ be the first prefix of $I$ that triggered \textcolor{blue}{(RT)}. Then, $R$ is a valid reason for $\Pi$, that is, $\Pi \models R$.
\end{theorem}
\begin{proof}
    To prove this statement, we have to show that for each interpretation $I'$ fulfilling $I' \models \Pi$, it holds that $I' \models R$, i.e., $\Pi \models R$.
    To show the latter we prove that $\Pi \cup \{\leftarrow \overline{\ell'} \mid \ell' \in \mathit{imp}(R,\ell)\} \cup \{\leftarrow \ell \}$ is unsatisfiable.
    By assumption, $\mps_\sigma(J,\overline{\ell}) < \bnd_\sigma$. \textcolor{blue}{By the definition of $\mps$~\eqref{eq:mps}, if $J$ was extended with $\overline{\ell}$ then the constraint could not be satisfied. 
    More formally, there is no total interpretation $K \supseteq J \mathbin{\dot\cup} \{\overline{\ell}\}$ such that $K \models \sigma$.}
    By Theorem~\ref{thm:mps:reason:eq:rt}, $\mps(\mathit{imp}(R,\ell),\overline{\ell}) = \mps(J,\overline{\ell}) < \bnd_\sigma.$
    Towards a contradiction, suppose $\Pi \cup \{\leftarrow \overline{\ell'} \mid \ell' \in \mathit{imp}(R,\ell)\} \cup \{\leftarrow \ell \}$ is satisfiable, then there exists an interpretation $I'$ such that $I' \models \Pi \cup \{\leftarrow \overline{\ell'} \mid \ell' \in \mathit{imp}(R,\ell)\} \cup \{\leftarrow \ell \}$. Thus $\mathit{imp}(R,\ell) \cup \{\overline{\ell}\} \subseteq I'$, \textcolor{blue}{thanks to Lemma~\ref{lemma:maximal_subset:antimonote}} $\mps_\sigma(I',\overline{\ell}) \le \mps_\sigma(\mathit{imp}(R,\ell),\overline{\ell}) < \bnd$, i.e., $I' \cup \{\overline{\ell}\}= I' \not\models \sigma$.
    Hence $I' \not\models \Pi$, contradicting our initial assumption.
\end{proof}

It is straightforward to check that in the proof of the previous theorem, we can employ instead of Theorem~\ref{thm:mps:reason:eq:rt} the corresponding Theorem~\ref{thm:mps:reason:eq:rf} to show that the reason of the form~\eqref{eq:amosum:reason:rf} generated by (RF) is indeed a valid reason. 

\begin{theorem}
  Suppose $\Pi$ is a program containing an \textsc{amosum} constraint $\sigma$ of the form \eqref{eq:amosum}
    and $I$ is an assignment such that $\sum_{\ell \in \mathit{lits}_\sigma|_s}{I(\ell)} \leq 1$ for all $s \in \mathit{parts}_{\sigma} $. Let $\overline{\ell}$ be the literal inferred by \eqref{eq:amosum:propagation:false} of (RF), $R$ be the associated reason  of the form \eqref{eq:amosum:reason:rf}, and $J$ be the first prefix of $I$ that triggered (RF). 
    Then, $R$ is a valid reason for $\Pi$, that is $\Pi \models R$.
\end{theorem}

{\color{blue}
\begin{proof}
    To show $\Pi \models R$, we prove that $\Pi \cup \{\leftarrow \overline{\ell'} \mid \ell' \in \mathit{imp}(R,\overline{\ell})\} \cup \{\leftarrow \overline{\ell}\}$ is unsatisfiable.
    By assumption, $\mps_\sigma(J,\ell) < \bnd_\sigma$, i.e., by definition of $\mps$~\eqref{eq:mps}, if $J$ was extended with $\ell$, then the constraint could not be satisfied. More formally, there is no total interpretation $K\supseteq J\mathbin{\dot\cup}\{\ell\}$ such that $K\models\sigma$.
    By Theorem~\ref{thm:mps:reason:eq:rf}, $\mps(\mathit{imp}(R,\overline{\ell}),\ell) = \mps(J,\ell) < \bnd_\sigma$.
    Towards a contradiction, suppose $\Pi \cup \{\leftarrow \overline{\ell'} \mid \ell' \in \mathit{imp}(R,\overline{\ell})\} \cup \{\leftarrow \naf \ell\}$ is satisfiable; then there exists an interpretation $I'$ such that $I' \models \Pi \cup \{\leftarrow \overline{\ell'} \mid \ell' \in \mathit{imp}(R,\overline{\ell})\} \cup \{\leftarrow \naf\ell\}$.
    Thus $\mathit{imp}(R,\overline{\ell}) \cup \{\ell\} \subseteq I'$, and by Lemma~\ref{lemma:maximal_subset:antimonote},
    $$\mps_\sigma(I',\ell) \le \mps_\sigma(\mathit{imp}(R,\overline{\ell}),\ell) < \bnd_\sigma,$$
    i.e., $I' \cup \{\ell\} = I' \not\models \sigma$.
    Hence $I' \not\models \Pi$, contradicting our initial assumption.
\end{proof}
}




\section{Reason Minimization}
\label{sec:minimization}
In this section, we  discuss some important properties of a reason. Subsequently, we introduce the concept of \textit{unessential literal} and explain the notion of \textit{increment}.
Using these concepts, we define an algorithm for constructing a subset-minimal reason, prove its correctness, and establish polynomial-time complexity.
\textcolor{blue}{During the following discussion, since the minimization process is applied to a single \textsc{amosum} constraint ($\sigma$), all associated notation in this subsection, including ($\bnd$) and ($\mps$), is understood with respect to ($\sigma$); the corresponding subscripts are therefore omitted whenever they are clear from the context. }
Furthermore, without loss of generality, we assume that for every literal $\ell$ occurring in the reason, either $\ell \in \lits$ or its complement $\overline{\ell} \in \lits$, since, whenever this is not the case, $\ell$ can be trivially removed from the reason, as it is trivially unessential.

\subsubsection*{Unessential Literals}
\label{sec:properties_reason}
Every reason $R$ is entailed by the program, i.e., $\Pi \models R$.
We use this property to remove literals while preserving entailment. 
Specifically, let $R_1,R_2$ be two reasons of $\ell$ of the form:
\begin{align}
    \label{eq:reasons_1}
    &R_1 = \{\ell\} \cup \{\ell_1, \hdots, \ell_n, y\} \\ 
    \label{eq:reasons_2}
    &R_2 = \{\ell\} \cup \{\ell_1, \hdots, \ell_n, \overline{y}\}, 
\end{align}
where $n \ge 0$.
In this case $R_1 \setminus \{y\} = R_2 \setminus \{\overline{y}\}$ and $\ell,y \in R_1$ and $\ell,\overline{y} \in R_2$.
Let $\Pi$ be a program such that $\Pi \models R_1$ and $\Pi \models R_2$.
If $I \models \Pi$ then $I \models R_1$ and $I \models R_2$.
We have seen in Section~\ref{sec:prel} that if $I \models R_1$ and $I \models R_2$ then $I \models  R_1 \otimes_{y} R_2$, hence $\Pi \models R_1 \otimes_{y} R_2 = \{\ell\} \cup \{\ell_1, \hdots, \ell_n\} = R$.
Since $\ell \in R$ and $\Pi \models R$ then $R$ is a reason of $\ell$ in $\Pi$.
\begin{example}[Continuing Example~\ref{ex:amosum:propagation:more-2}]
    In case $I=[\overline{y},\, \overline{w}]$, we have that $\mathit{reason}(z)=\{w, y, z\} = R_1$.
    Let us consider $R_2 =\{w, \overline{y}, z\}$.
    The related max possible sum is $\mps(\mathit{imp}(R_2,z), \overline{z}) = 2 < \bnd_{\sigma_4} = 3$ and we have seen that this is a sufficient condition to confirm that also $R_2$ is a reason.
    Assume then $\Pi$ is a program containing 
    the constraint $\sigma_4$. When we define 
    $R = R_1 \otimes_y R_2 = \{w, z\}$, then $ \Pi \models R$ holds.
\end{example}
We can see in the above example that $y$ can be removed from the reason, getting $R$, and $R$ continues to be a reason, so $y$ is \textit{unessential}.
Formally, a literal $\ell'$ is \textit{unessential} in a reason $R$ (with $\ell' \in R$) of a literal $\ell$ if $R \cupdot \{\overline{\ell'}\}$ is a reason of $\ell$.
Indeed, $R$ and $R \cupdot \{\overline{\ell'}\}$ contain complementary occurrences of $\ell'$, and resolving them on $\ell'$ gives \textcolor{blue}{$R \otimes_{\ell'} (R \cupdot \{\overline{\ell'}\}) = R \setminus \{\ell'\}$} which is still a reason for $\ell$.
Therefore, all unessential literals of a reason $R$ can be removed from it.

\subsubsection*{Increment}
\label{sec:increment}
{\color{blue}
The \textit{increment} measures the increase in the \textit{mps} when moving from $\mathit{imp}(R, \ell)$ to $\mathit{imp}(R \setminus \{\ell'\}, \ell)$ for some $\ell' \in R$.
\begin{definition}
    \label{prop:inc}
        Let $\mathit{reason}(\ell) = R$ be a reason for a literal $\ell$ of the form \eqref{eq:amosum:reason:rt} or \eqref{eq:amosum:reason:rf}.
    The \emph{increment} of a literal $\ell' \in R$ represents:
    %it is computed as:
    $$\mathit{inc}(R, \ell, \ell') = \mps(\mathit{imp}({R \setminus \{ \ell' \}},\ell),\overline{\ell}) - \mps(\mathit{imp}(R,\ell),\overline{\ell}).$$
\end{definition}
}

The increment can also be expressed in terms of two reasons of the form~\eqref{eq:reasons_1} and~\eqref{eq:reasons_2}, namely $R$ and $R \cupdot \{\overline{\ell'}\}$.
\begin{proposition}
    Let $R$ be a reason for a literal $\ell$ of the form \eqref{eq:amosum:reason:rt} or \eqref{eq:amosum:reason:rf}.
    The \emph{increment} of a literal $\ell' \in R$ is defined as:
    $$\mathit{inc}(R, \ell, \ell') = \max\{mps(\mathit{imp}({R \cupdot \{ {\overline{\ell'}} \}},\ell),\overline{\ell}) - mps(\mathit{imp}({R},\ell),\overline{\ell}),0\}.$$
\end{proposition}
\textcolor{blue}{The increment is lower bounded by $0$. Indeed, a negative increment would imply that removing $\ell'$ from the reason decreases the \textit{maximum possible sum} (\textit{mps}), whereas the updated $\mps$ is intended to represent the maximum value reached when $\ell'$ is undefined.
To see this, observe that removing $\ell'$ from $R$ yields a reason candidate $R'$ in which $\ell'$ is undefined. Consequently, by construction, the new $\mps$ is the maximum value obtained by assigning $\ell'$ either \textit{true} or \textit{false}. Therefore, the new $\mps$ cannot be smaller than the original one. In particular, if
$$
\mps(\mathit{imp}(R \cupdot \{\overline{\ell'}\},\ell), \overline{\ell})
<
\mps(\mathit{imp}(R,\ell),\overline{\ell}),
$$
then the original $\mps$ is already the larger value, and the increment is equal to $0$.}

% Let us now talk more in details about how the increment is efficiently computed; to do so, let's introduce this new function: $\mathit{mwh}_\sigma(s) := \max\{w \mid (w : \ell\ [s]) \in \sigma\}$. $\mathit{mwh}_\sigma(s)$ represents the maximum possible weight without considering any interpretation.
Now, in order to prove the correctness of our increment function, the following theorem is shown.
\begin{theorem}
    \label{th:increment}
    Let $R$ be a reason of a literal $\ell \in R$ of the form \textcolor{blue}{ \eqref{eq:amosum:reason:rt} or \eqref{eq:amosum:reason:rf}}, a literal $\ell' \in R$ such that $\ell' \neq \ell$, two parts $p, s$ equal to $\mathit{part}(\ell)$ and $ \mathit{part}(\ell')$, respectively. Then, the following statements hold:
    \begin{enumerate}
        \item $\ell' \not\in \lits|_s $ implies that $\mps(\mathit{imp}(R \cupdot \{ {\overline{\ell'}} \}, \ell),\textcolor{blue}{\overline{\ell}}) - \mps(\mathit{imp}(R,\ell), \textcolor{blue}{\overline{\ell}}) $ is equal to $ \textcolor{blue}{\mwh(\mathit{imp}(R \cupdot \{ {\overline{\ell'}} \},\ell), s)} - \wh(\overline{\ell'})$.
        \item $\ell' \in \lits|_s $ implies $ \mps(\mathit{imp}(R \cupdot \{ {\overline{\ell'}} \}, \ell),\textcolor{blue}{\overline{\ell}}) - \mps(\mathit{imp}(R,\ell), \textcolor{blue}{\overline{\ell}}) $ is equal to $ \wh(\ell') - \mwh(\mathit{imp}(R,\ell) \cupdot \{\overline{\ell}\}, s)$.
    \end{enumerate}
\end{theorem}

{\color{blue}
\begin{proof}
   \emph{Flipping} a literal $\ell'$ from $\mathit{imp}(R,\ell)$ to $\mathit{imp}(R \cupdot \{ {\overline{\ell'}} \},\ell)$ affects just the part of $\ell'$,
    that is, for each $s'$ such that $s'\neq s$ then $\mwh(\mathit{imp}(R \cupdot \{ {\overline{\ell'}} \},\ell), s') $ is equal to $ \mwh(\mathit{imp}(R,\ell), s')$.
    Hence, to compute the increment is enough to locally view the increment in part $s$.
    \begin{enumerate}
        \item If $\ell' \not\in \lits|_s$ (i.e., $\overline{\ell'} \in \lits|_s$) then by construction $\mathit{imp}(R,\ell) \cap \lits|_s = \{\overline{\ell'}\}$, that is, $\mwh(\mathit{imp}(R,\ell), s) = \wh(\overline{\ell'})$, and $s \neq p$.
    By construction $\ell' \in \mathit{imp}(R\cupdot\{\overline{\ell'}\})$ (i.e., $\overline{\ell'} \in \overline{\mathit{imp}(R\cupdot\{\overline{\ell'}\})}$).
    Thus, $\mathit{imp}(R\cupdot\{\overline{\ell'}\})(\overline{\ell'}) = 0$, and $\mwh(\mathit{imp}(R \cupdot \{ {\overline{\ell'}} \},\ell),s)$ is equal to the maximum literal in the remaining $\lits|_s \setminus \{\overline{\ell'}\}$.
    Hence,
    \begin{align*}
        & \mwh(\mathit{imp}(R \cupdot \{ {\overline{\ell'}} \},\ell), s) - \mwh(\mathit{imp}(R,\ell), s) = \\
        & \mwh(\mathit{imp}(R \cupdot \{ {\overline{\ell'}} \},\ell), s) - \wh(\overline{\ell'}).
    \end{align*}
    \item If $\ell' \in \lits|_s$ then $\mathit{imp}(R \cupdot \{ {\overline{\ell'}} \},\ell) \cap \lits|_s = \{\ell'\}$.
    Thus, $$\mwh(\mathit{imp}(R \cupdot \{ {\overline{\ell'}} \},\ell),s) = \wh(\ell').$$
    %
    On the one hand, if $s \neq p$ then setting $\ell$ to false has no effect on part $s$, that is $\mwh(\mathit{imp}(R \cupdot \{\overline{\ell'}\}, \ell), s) = \mwh(\mathit{imp}(R \cupdot \{\overline{\ell'}\}, \ell) \cupdot \{\overline{\ell}\}, s) = \wh(\ell')$.
    %
    On the other hand, if $s=p$, then $R$ cannot be a reason of form \eqref{eq:amosum:reason:rf}, since (RF) excludes $part(\ell)$. Hence $R$ is of form \eqref{eq:amosum:reason:rt}, associated with (RT), and $\ell\in lits|_p$.
    Given that $\ell' \neq \ell$, setting $\ell$ to false does not change the maximum weight of $p$, that is, $$\mwh(\mathit{imp}(R \cupdot \{\overline{\ell'}\}, \ell), s) = \mwh(\mathit{imp}(R \cupdot \{\overline{\ell'}\}, \ell) \cupdot \{\overline{\ell}\}, s) = \wh(\ell').$$
    %
    Hence, let $J = \mathit{imp}(R,\ell)$ and $J' = \mathit{imp}(R \cupdot \{ {\overline{\ell'}} \},\ell)$, $$\mwh(J' \cupdot \{\overline{\ell}\}, s) - \mwh(J \cupdot \{\overline{\ell}\}, s) = \wh(\ell') - \mwh( J \cupdot \{\overline{\ell}\}, s).$$
    \end{enumerate}
    
\end{proof}
}

By construction, we have that $\inc(R, \ell, \ell') \geq 0$ for all $\ell' \in R \setminus \{\ell\}$, since removing  a literal from a reason in the worst case does not affect the $\mps$.
Algorithm \ref{alg:increment} computes the increment as previously described.
\begin{algorithm}[ht]\small
    \caption{inc}
        \label{alg:increment}
        \SetKwInOut{Input}{Input}
        \Input{A reason $R$ of a literal $\ell$, a literal $\ell' \in R$}.
        \SetKwBlock{Loop}{Loop}{end}
        \Begin{
            $J \gets \mathit{imp}(R,\ell)$\;
            $s \gets \mathit{part}(\ell')$ \;
            \lIf {$\ell' \not\in \lits|_s$}{
                \Return{\textcolor{blue}{$\max\{\mwh(\mathit{imp}(R \cupdot \{ {\overline{\ell'}} \},\ell), s) - \mathit{wh}(\overline{\ell'}),0\}$}}\label{alg:increment:line:truecase}}
            \Return{$\max\{\wh(\ell') - \mathit{mwh}(J \cupdot \{ {\overline{\ell}} \},s),0\}$\;}\label{alg:increment:line:falsecase}            
        }
\end{algorithm}

\begin{example}[Continuing Example \ref{ex:amosum:propagation:more-2}]
    \label{ex:increment}    
    Consider again $R=\mathit{reason}(z)=\{w, y, z\}$.
    In this case, $J = \mathit{imp}(R,z) = [\overline{w}, \overline{y}]$ ,
    $s = \mathit{part}(y) = 1$, $\mwh(J \cupdot \{\overline{z}\}, 1) = 1$ and $\wh(y) = 2$.
    Thus, $\inc(R, z, y) = 2 - 1 = 1$, given that
    $y \in \lits|_1$ and $\max\{2 - 1, 0\} = 1$.
\end{example}

\textcolor{blue}{For the purpose of the results presented in this section, let $R' \subseteq R$ be a clause containing $\ell$. We define $R'$ to be a $\sigma$-reason for $\ell$ iff
$$
\mps(\mathit{imp}(R', \ell), \overline{\ell}) < \bnd.
$$
Every $\sigma$-reason considered here is also a valid reason for $\ell$ with respect to $\Pi$. However, the converse does not necessarily hold as it might happen that a clause is entailed by other rules of $\Pi$ without satisfying the above condition.}
\paragraph{Maximal Subset Sum}
At this point, an important parallel must be drawn: minimizing a reason $R$ equates to maximizing the number of literals removed from $R$. 
\begin{algorithm}[ht]\small
    \caption{Maximal Subset Sum (MSS)}
    \label{alg:maximal_subset}
    \SetKwInOut{Input}{Input}
    \SetKwInOut{Parameters}{Parameters}
    \SetKwInOut{Output}{Output}
    \Input{\textcolor{blue}{A reason $R$ of the form~\eqref{eq:amosum:reason:rt} or~\eqref{eq:amosum:reason:rf} for a literal $\ell$, and a bound $\bnd$}}
    \Parameters{\textit{inc} function}
    \Output{\textcolor{blue}{A subset-maximal set of literals such that $R \setminus S$ is a $\sigma$-reason for $\ell$}} 
    \Begin{
        $\mathit{S} \gets [\,]$ \label{alg:maximal_subset:init:S}  \\
        $\mathit{ci} \gets 0$ \label{alg:maximal_subset:init:ci} \\
        \For{$\ell' \in R \setminus \{\ell\} $}{ \label{alg:maximal_subset:for}
            \If{$\mathit{ci} + \inc(R \setminus \mathit{S}, \ell,\ell') < \textcolor{blue}{\bnd_\sigma - \mps(\mathit{imp}(R,\ell),\overline{\ell})}$}{  \label{alg:maximal_subset:condition:ci}
                $\mathit{ci} \gets ci + \inc(R \setminus \mathit{S}, \ell,\ell')$ \label{alg:maximal_subset:update:ci} \\
                $\mathit{S} \gets \mathit{S} \cup \{\ell'\} $ \label{alg:maximal_subset:update:S} \\
            }
        }
        \Return{\textit{S}}
    }
\end{algorithm}

\textcolor{blue}{Algorithm~\ref{alg:maximal_subset} takes as input a reason $R$ of the form~\eqref{eq:amosum:reason:rt} or~\eqref{eq:amosum:reason:rf} for a literal $\ell$ and a bound $\bnd$, and returns a subset-maximal set $S$ such that $R \setminus S$ is a $\sigma$-reason for $\ell$. 
It starts with an empty set $S$ (line~\ref{alg:maximal_subset:init:S}) and \textit{current increment} (\textit{ci}) equal to $0$.
Variable \textit{ci} represents the increment of the \textit{mps} from $\mps(\mathit{imp}(R, \ell), \overline{\ell})$ to $\mps(\mathit{imp}(R \setminus S, \ell), \overline{\ell})$.
Let $\ell' \in R \setminus \{\ell\}$ be the first candidate literal. At this point, \textit{ci} is $0$ because no literal has yet been removed. If $\inc(R, \ell, \ell') < \bnd - \mps(\mathit{imp}(R,\ell), \overline{\ell})$ then $\ell'$ is unessential: the current increment is increased by $\inc(R, \ell, \ell')$ and $S$ is extended with $\ell'$. Otherwise, $\ell'$ is essential and remains in $R$.
Steps~\eqref{alg:maximal_subset:for}--\eqref{alg:maximal_subset:update:S} are repeated for every literal in $R\setminus \{\ell\}$ and after that the subset-maximal set $S$ is returned. }

{\color{blue}
To prove the correctness of \textit{MSS}, we first establish the following invariant, which is used in Theorem~\ref{thm:mss}.

\begin{proposition}
\label{prop:maximal_subset:ci}
    Let $R$ be a reason of the form~\eqref{eq:amosum:reason:rt} or~\eqref{eq:amosum:reason:rf} for a literal $\ell$, and let $S$ be the current set being evaluated in the for loop of \textit{MSS}. The variable $ci$ represents the quantity
    $$\mps(\mathit{imp}(R \setminus S, \ell), \overline{\ell}) - \mps(\mathit{imp}(R, \ell), \overline{\ell}).$$
\end{proposition}
\begin{proof}
    Let's suppose that $\ell_1, \dots, \ell_n$ are the ordered literals added to $S$ at line~\ref{alg:maximal_subset:update:S} and $ci_1, \dots, ci_n$ are the related increments updates at line~\ref{alg:maximal_subset:update:ci}, with $n = \vert S\vert$. The proof is done by induction:
    \begin{enumerate}
        \item[(i)] base case: $$ci_1 = \mps(\mathit{imp}(R \setminus \{\ell_1\}, \ell), \overline{\ell}) - \mps(\mathit{imp}(R, \ell), \overline{\ell})$$
        \item[(ii)] inductive case, for $2\le k \le n$: $$ci_k = \mps(\mathit{imp}(R \setminus S_k, \ell), \overline{\ell}) - \mps(\mathit{imp}(R, \ell), \overline{\ell})$$ where $S_k = \{\ell_i \in S \mid i \le k\}$.
    \end{enumerate}
    Let's prove (i) and then (ii).
    \begin{enumerate}
        \item[(i)] 
        Until $S$ is first updated, $ci = 0$ (i.e., no literals have been removed from the reason).
        When $ci$ is updated from $0$ to $ci_1$, it means that $\inc(R,\ell, \ell_1) < \bnd - \mps(\mathit{imp}(R, \ell), \overline{\ell})$; thus, $ci_1 = \inc(R, \ell,\ell_1)$ and consequently $S_1 = \{\ell_1\}$.
        By Definition~\ref{prop:inc}, $$ci_1 = \mps(\mathit{imp}({R \setminus \{\ell_1\} \}},\ell),\overline{\ell}) - \mps(\mathit{imp}(R,\ell),\overline{\ell}).$$
        %
        \item[(ii)] By induction $ci_{k-1} = \mps(\mathit{imp}(R \setminus S_{k-1}, \ell), \overline{\ell}) - \mps(\mathit{imp}(R , \ell), \overline{\ell})$.
        By construction $ci_k = ci_{k-1} + \inc(R\setminus S_{k-1}, \ell, \ell_k).$
        By Definition~\ref{prop:inc}, $$\inc(R\setminus S_{k-1}, \ell, \ell_k) = \mps(\mathit{imp}({R \setminus S_k},\ell),\overline{\ell}) - \mps(\mathit{imp}(R \setminus S_{k-1},\ell),\overline{\ell}).$$
        Hence, $$ci_k = \mps(\mathit{imp}({R \setminus S_k},\ell),\overline{\ell}) - \mps(\mathit{imp}(R,\ell),\overline{\ell}).$$
    \end{enumerate}
    Thus, the 
    %thesis 
    result follows.
\end{proof}

We can now proceed with the following theorem. 

\begin{theorem}
\label{thm:mss}
    Let $R$ be a reason of the form~\eqref{eq:amosum:reason:rt} or~\eqref{eq:amosum:reason:rf} for a literal $\ell$ derived applying (RT)~\eqref{eq:amosum:propagation:true} or (RF)~\eqref{eq:amosum:propagation:false}. Algorithm~\ref{alg:maximal_subset} (Maximal Subset Sum) returns a subset-maximal $S \subseteq R \setminus \{\ell\} $ such that $R \setminus S$ is a $\sigma$-reason for $\ell$; equivalently, $R\setminus S$ is subset-minimal among the $\sigma$-reasons contained in $R$.
\end{theorem}


\begin{proof}
    First, observe that \(R\setminus S\) is a $\sigma$-reason for \(\ell\). Initially, \(S=\emptyset\) and \(R\) is a reason for \(\ell\).
     Since $R$ is derived applying (RT) or (RF), by Theorems~\ref{thm:mps:reason:eq:rt} and~\ref{thm:mps:reason:eq:rf},
    $$
    \mps(\mathit{imp}(R, \ell), \overline{\ell}) < \bnd,
    $$
    and thus $R$ is a $\sigma$-reason for $\ell$.
Moreover, a literal is added to \(S\) only when the condition at line~\ref{alg:maximal_subset:condition:ci} guarantees, by Propositions~\ref{prop:maximal_subset:ci} and~\ref{prop:inc}, that
\[
\mps(\mathit{imp}(R\setminus S,\ell),\overline{\ell})<\bnd.
\]
Hence, this property is preserved throughout the execution of the algorithm.

We now prove that \(S\) is subset-maximal. Arguing by contradiction, suppose that $S$ is not subset-maximal with respect to the $\sigma$-reason condition. Then, there exists a set $X$ such that
$$
S\subset X\subseteq R
$$
and \(R\setminus X\) is a $\sigma$-reason for \(\ell\). Let \(x\in X\setminus S\), and let \(S_x\) and \(\mathit{ci}_x\) denote, respectively, the values of \(S\) and \(\mathit{ci}\) when \(x\) was examined by the algorithm. Since the algorithm only adds literals to \(S\), we have \(S_x\subseteq S\).

Since \(x\notin S\), the condition at line~\ref{alg:maximal_subset:condition:ci} was not satisfied when \(x\) was examined. Therefore,
\[
\mathit{ci}_x+\inc(R\setminus S_x,\ell,x)
\geq
\bnd_\sigma-\mps(\mathit{imp}(R,\ell),\overline{\ell}).
\]
By Proposition~\ref{prop:maximal_subset:ci},
\[
\mathit{ci}_x=
\mps(\mathit{imp}(R\setminus S_x,\ell),\overline{\ell})
-\mps(\mathit{imp}(R,\ell),\overline{\ell}),
\]
and thus
\[
\mps(\mathit{imp}(R\setminus S_x,\ell),\overline{\ell})
+\inc(R\setminus S_x,\ell,x)
\geq \bnd.
\]

Since \(R\setminus S_x\) is a $\sigma$-reason for \(\ell\), we have
\[
\mps(\mathit{imp}(R\setminus S_x,\ell),\overline{\ell})<\bnd.
\]
Consequently, \(\inc(R\setminus S_x,\ell,x)>0\). By Proposition~\ref{prop:inc}, it follows that
\[
\inc(R\setminus S_x,\ell,x)
=
\mps\bigl(\mathit{imp}((R\setminus S_x)\cupdot\{\overline{x}\},\ell),
          \overline{\ell}\bigr)
-\mps(\mathit{imp}(R\setminus S_x,\ell),\overline{\ell}).
\]
Therefore,
\[
\mps\bigl(\mathit{imp}((R\setminus S_x)\cupdot\{\overline{x}\},\ell),
          \overline{\ell}\bigr)
\geq\bnd.
\]

Since \(S_x\subseteq S\subset X\) and \(x\in X\), we have
\[
R\setminus X
\subseteq
(R\setminus S_x)\cupdot\{\overline{x}\}.
\]
Consequently,
\[
\mathit{imp}(R\setminus X,\ell)
\subseteq
\mathit{imp}((R\setminus S_x)\cupdot\{\overline{x}\},\ell).
\]
By Lemma~\ref{lemma:maximal_subset:antimonote},
\[
\mps(\mathit{imp}(R\setminus X,\ell),\overline{\ell})
\geq
\mps\bigl(\mathit{imp}((R\setminus S_x)\cupdot\{\overline{x}\},\ell),
          \overline{\ell}\bigr)
\geq\bnd.
\]
Hence, \(R\setminus X\) is not a $\sigma$-reason for \(\ell\), contradicting our assumption. Therefore, \(S\) is subset-maximal with respect to the $\sigma$-reason condition, and \(R\setminus S\) is subset-minimal among the $\sigma$-reasons contained in $R$.
\end{proof}

Note that this result does not establish subset-minimality with respect to all consequences of $\Pi$, since additional literals may be redundant because of other rules not considered by Algorithm~\ref{alg:maximal_subset}.
}

\begin{example}[Continuing Example \ref{ex:amosum:propagation:more-2}]
    Consider $R = \mathit{reason}(z) = \{w, y, z\}$ as reason to minimize. In this case $\bnd_{\sigma_4} - \mps(\mathit{imp}(R, z), \overline{z})= 3 - (1 + 0) = 2$. Let us iterate over $\mathit{reason}(z) \setminus \{z\} = \{w, y\}$.
    The first literal to be considered is $w$. Given that $w \in \lits_{\sigma_4}|_2$, $\wh(w)=3$ and $\mwh(\mathit{imp}(R, z) \cupdot \{\overline{z}\}, 2) = 0$ then $\inc(R,z,w) = 3 - 0 = 3$. Given that $3 \ge \bnd_{\sigma_4} - \mps(\mathit{imp}(R, z), \overline{z})$ then $w$ is \textit{not} removed from $R$ and $\mathit{ci}$ remains equal to $0$. Now, let us move on to $y$. We have seen in Example~\ref{ex:increment} that $\mathit{inc}(R,z,y) = 1 < \bnd_{\sigma_4} - \mps(\mathit{imp}(R, z), \overline{z})$. Thus, $y$ can be removed from $R$, i.e., $y$ is added to the set of unessential literals $S$.
\end{example}

\paragraph{Complexity for \textit{MSS}}
\textcolor{blue}{Let's analyze complexity for \textit{MSS} (Algorithm \ref{alg:maximal_subset}).
The loop at line~\ref{alg:maximal_subset:for} performs at most $n=|R|$ iterations. Each increment is computed in $O(n)$ time because finding the maximum weight in a part is linear in the number of its literals. Thus, \textit{MSS} runs in $O(n^2)$ time and computes a subset-minimal reason in polynomial time.}


\subsection{Minimization on the fly}
In general, minimizing a reason improves search-space pruning because shorter clauses can become unit earlier. However, applying a separate quadratic procedure to compute a subset-minimal reason can be counterproductive, especially when the initial and minimized reasons have similar sizes.
For this reason, we compute the subset-minimal reason
\textit{on the fly}, that is, during its construction.
In these cases, we assume without loss of generality that we are just handling literals within the aggregate.
To extend to no-aggregate literals, it is enough to ``ignore'' them, i.e., to not put them in the 
reason by default.
{\color{blue}
To construct a subset-minimal reason directly, we use Algorithm~\ref{alg:minfly} (\textit{MRF}).
\begin{algorithm}[ht]\small
    \caption{Minimal Reason on the Fly (MRF)}
    \label{alg:minfly}
    \SetKwInOut{Input}{Input}
    \SetKwInOut{Parameters}{Parameters}
    \SetKwInOut{Output}{Output}
    \Input{An assignment $J$, a literal $\ell$, and the inference rule $r \in \{\mathrm{(RT)},\mathrm{(RF)}\}$ that derived $\ell$.}
    \Output{A subset-minimal $\sigma$-reason $R$ for literal $\ell$} 
    \Begin{
        $\mathit{ci}\, \gets 0$ \label{alg:minfly:line:initSum}\;
        $R\, \gets \{\ell\}$ \label{alg:minfly:line:reason}\;
        \For{$p \in \parts$}{
            \lIf{$p = \mathit{part}(\ell)$ \textbf{ and } $r=\mathrm{(RF)}$}{\textbf{continue}\label{alg:minfly:line:continue}}
            \If{$\lits_p \cap J = \{\ell'\}$}{ \label{alg:minfly:line:trueCase}
                $s\, \gets \mathit{part}(\ell')$\;
                $\mathit{inc} = \max\{\mwh(J \cupdot \{ \overline{\ell'} \}, s) - \wh(\ell'), 0\}$\;
                \lIf{$\mathit{ci} + \mathit{inc} < \bnd - \mps(J,\overline{\ell})$} {\label{alg:minfly:line:trueCase:checkRed}
                    $\mathit{ci} \gets \mathit{ci} + \mathit{inc}$\label{alg:minfly:line:TrueCase:redundant}}\lElse{
                    $R \gets R \cup \{\overline{\ell'}\}$\label{alg:minfly:line:TrueCase:noredundant}}
            }\Else{ \label{alg:line:falseCase}
                \tcp{literals are ordered by decreasing weight}
                \For{$\ell' \in \lits_p$}{
                    $s\, \gets \mathit{part}(\ell')$\;
                   \lIf{$\wh(\ell') \le \mwh(J \cupdot \{\overline{\ell}\},p)$}{  \label{alg:minfly:line:belowMwh}\textbf{break}\label{alg:minfly:line:belowMwh:break}}
                   $\mathit{inc} = \max\{\wh(\ell') - \mathit{mwh}(J \cupdot \{ {\overline{\ell}} \},s),0\}$\;
                   \If{$\ell' \in \overline{J}$}{\label{alg:minfly:line:falseCase:isFalse}
                        \lIf{$\mathit{ci} + \mathit{inc} < \bnd - \mps(J,\overline{\ell})$}{\label{alg:minfly:line:falseCase:checkRed}
                            $\mathit{ci} \gets \mathit{ci} + \mathit{inc}$\label{alg:minfly:line:falseCase:redundant};
                            \textbf{break}\label{alg:minfly:line:falseCase:break}}
                        \lElse{
                            $R \gets R \cup \{\ell'\}$\label{alg:minfly:line:falseCase:noredundant}}
                   }
                }   
            }
        }
        \Return{\textit{R}}
    }
\end{algorithm}
Let us 
%break down 
go through this algorithm step by step.

Initially, the variables $\mathit{ci}$ and $R$ are initialized to $0$ and $\{\ell\}$, respectively. Analogously to Algorithm~\ref{alg:maximal_subset}, $\mathit{ci}$ plays the role of the current increment.

For each part $p$, we first check whether the input inference rule $r = \mathit{(RF)}$ and $p = \mathit{part}(\ell)$, in this case, $\mathit{part}(\ell)$ does not appear in the reason, and it is skipped (line~\ref{alg:minfly:line:continue}). 
Otherwise, we check whether there is a true literal $\ell' \in \lits|_p$ w.r.t. $J$. If so (line~\ref{alg:minfly:line:trueCase}), the cost of omitting its complement from the reason is $\max\{\mwh(J \cupdot \{ \overline{\ell'} \}, s) - \wh(\ell'), 0\}$, which is the increment computed at line~\ref{alg:increment:line:truecase} of Algorithm~\ref{alg:increment}. If the current increment plus this contribution remains below $\bnd - \mps(J,\overline{\ell})$, the reason literal is unessential (line~\ref{alg:minfly:line:TrueCase:redundant}); otherwise, $\overline{\ell'}$ is added to the reason (line~\ref{alg:minfly:line:TrueCase:noredundant}).


If part $p$ contains no true literal, its literals are processed in decreasing order of weight. If the weight of $\ell'$ does not exceed $\mwh(J \cupdot \{\overline{\ell}\},p)$ (line~\ref{alg:minfly:line:belowMwh}), then $\ell'$ cannot increase the maximum possible sum. All remaining literals have no greater weight and are skipped as well (line~\ref{alg:minfly:line:belowMwh:break}).

Otherwise, only false literals, i.e., literals in $\overline{J}$, are candidates for the reason (line~\ref{alg:minfly:line:falseCase:isFalse}). Their increment is $\max\{\wh(\ell')-\mwh(J\cupdot\{\overline{\ell}\},s),0\}$. If adding it to $\mathit{ci}$ remains below $\bnd-\mps(J,\overline{\ell})$, then $\ell'$ is unessential. Once it is omitted, it becomes the maximum available literal of the part; consequently, all subsequent literals, which have no greater weight, have zero marginal increment and can be omitted as well. 
Instead, if adding it to $\mathit{ci}$ reaches or exceeds $\bnd-\mps(J,\overline{\ell})$, $\ell'$ is added to the reason and the scan continues.

The following result makes the relationship with Algorithm~\ref{alg:maximal_subset} precise.
\begin{theorem}\label{thm:mrf}
Let $J$ trigger (RT) or (RF) for a literal $\ell$, and let $R_0$ be the corresponding reason of the form~\eqref{eq:amosum:reason:rt} or~\eqref{eq:amosum:reason:rf}, then Algorithm~\ref{alg:minfly} returns a subset-minimal $\sigma$-reason $R\subseteq R_0$ for $\ell$.
\end{theorem}
\begin{proof}
Consider Algorithm~\ref{alg:maximal_subset} applied to $R_0$, with candidates ordered by part and, within parts without a true literal, by decreasing weight. Algorithm~\ref{alg:minfly} performs the same omission test while $R_0$ is being generated. In particular, after every processed candidate, its variable $\mathit{ci}$ satisfies the invariant of Proposition~\ref{prop:maximal_subset:ci}:
$$
\mathit{ci}=\mps(\mathit{imp}(R_0\setminus S,\ell),\overline{\ell})-\mps(\mathit{imp}(R_0,\ell),\overline{\ell}),
$$
where $S$ is the set of candidates omitted so far. Hence, a literal is omitted exactly when the resulting maximum possible sum remains strictly below $\bnd$.
The only compressed step is the \textbf{break} in the false-literal case. Suppose that a false literal of weight $w$ is omitted. It then becomes available and determines a possible contribution of at least $w$ for that part. Every subsequent literal has weight at most $w$, and therefore has zero marginal increment; omitting all of them is equivalent to the corresponding zero-increment iterations of Algorithm~\ref{alg:maximal_subset}. The earlier \textbf{break} at line~\ref{alg:minfly:line:belowMwh:break} is justified in the same way, because all remaining literals already have zero increment.

Thus, Algorithm~\ref{alg:minfly} returns the same result as a valid execution of Algorithm~\ref{alg:maximal_subset}. The claim follows from Theorem~\ref{thm:mss}.
\end{proof}

Algorithm~\ref{alg:minfly} runs in $O(|\lits|^2)$ time. This does not add a second quadratic minimization pass: reason construction and minimization are integrated into the same quadratic procedure.


}


\section{ASP Language Extension and Implementation}
\label{sec:language-extension}
In this section, we detail the concrete implementation of the \textsc{amosum} propagator and its syntax in ASP.
Furthermore, we describe its integration into \textsc{wasp} and \textsc{clingo}~\citep{DBLP:conf/iclp/GebserKKOSW16}.

\textcolor{blue}{We deliberately introduce \textsc{amosum} as a constraint construct because this work focuses on propagation, reason generation, and learning. Its non-ground syntax still retains a central advantage of ASP by using variables to represent families of ground constraints compactly. Supporting arbitrary positive, default-negated, or recursive occurrences as first-class aggregate atoms requires a broader semantic treatment and is left for future work.}

\subsection{ASP construct}
\label{sub:sec:asp:const}

We propose the following syntax for a constraint of the type \eqref{eq:amosum}:
\begin{align}\label{eq:amosum:kr}
    %\#\amosum\{w_1 : \ell_1 \ [s_1];\ \cdots;\ w_n : \ell_n\ [s_n]\} \ge b : \mathit{body}
    \#\amosum\{w_1 : \ell_1, \mathit{conj}_1 \ [s_1];\ \cdots;\ w_n : \ell_n, \mathit{conj}_n\ [s_n]\} \ge b : \mathit{body}
\end{align}
where $n \geq 0$, $\ell_1,\ldots,\ell_n$ are literals, $\mathit{conj}_1,\ldots,\mathit{conj_n},\mathit{body}$ are possibly empty conjunctions of literals,
and $s_i$ is a term identifying the part associated with $\ell_i$ (note that $i \neq j$ does not imply $s_i \neq s_j$; parts are not necessarily singletons).
For an interpretation $I$, \textcolor{blue}{let $A_I:=\{i\in[1..n]\mid I\models\mathit{conj}_i\}$ be the active aggregate elements}. The semantics of \eqref{eq:amosum:kr} is as follows: $I \models \eqref{eq:amosum:kr}$ if either $I \not\models \mathit{body}$, or both of the following conditions hold:
\begin{itemize}
\item 
\textcolor{blue}{$I \models \textsc{sum}\{w_i : \ell_i \mid i\in A_I\} \ge b$}, and

\item
\textcolor{blue}{$I \models \textsc{amo}\{\ell_j \mid j\in A_I,\ s_j=s_i\}$ for each $i\in A_I$.}
\end{itemize}
\textcolor{blue}{Thus, when $\mathit{conj}_i$ is false, element $i$ contributes neither to the sum nor to its AMO group. In what follows, we assume that every $\mathit{conj}_i$ is evaluated deterministically before search; handling dynamically changing memberships would require extending both propagation and reason generation.}

\begin{example}
    \label{ex:kr:1}
    The Graph Coloring (GC) problem is the well-known problem where given a graph and a set of colors, the goal is to assign a color to each node such that connected nodes do not share the same color.
    \textcolor{blue}{In this example, we consider a partial weighted variant of the problem, in which each node can be assigned at most one color, but may also remain uncolored. Colors are associated with weights, and the total weight of the selected colors is required to reach a given threshold.}
    For modeling this problem, the following program $\Pi_\mathit{amo}$ can be used.
    $$\begin{array}{l}
        \{\mathit{col}(X,C) : \mathit{colour}(C)\} <= 1 \text{ :- } \mathit{node}(X). \\ 
        \text{:- } \mathit{col}(X1,C), \mathit{col}(X2,C), \mathit{edge}(X1,X2). \\
        \#\amosum\{ W: \mathit{col}(X,C), \mathit{cw}(C,W) [X]\} >= \mathit{M} : \mathit{min}(M). \\
    \end{array}$$
    This is the proposed formulation of the GC problem using \#\amosum.
    \textcolor{blue}{The first rule states that \textit{at most one} color can be chosen for each node.}
    The second rule models that two connected nodes cannot share the same color.
    Finally, the third rule \textcolor{blue}{requires the sum of all color weights to reach the lower bound. An equivalent ASP-Core-2 formulation~\citep{DBLP:journals/tplp/CalimeriFGIKKLM20} is obtained} by replacing the \#\amosum constraint in the last line with the following rule:
    $$\begin{array}{l}
        \text{:- } \#\mathit{sum}\{ W, X: \mathit{col}(X,C), \mathit{cw}(C,W)\} < \mathit{M} , \mathit{min}(M). \\
    \end{array}$$
\end{example}

\textcolor{blue}{Furthermore, we introduce a complementary construct, called \textsc{eosum}, which exploits exactly-one relations already encoded in the input program. An \textsc{eosum} occurrence is well formed if, for every ground part $s$, the literals in $\mathit{lits}_\sigma|_s$ coincide with a group subject to an explicitly declared EO constraint that is active whenever the \textsc{eosum} constraint is active. The \textsc{eosum} propagator does not enforce these EO constraints itself: it derives additional weighted-sum consequences, while the solver's built-in propagator enforces the EO relations and derives true literals. We delegate this propagation because the built-in EO reasons can be significantly smaller. Variables in part annotations are instantiated during grounding, and equal ground part terms identify the same EO group.}

\begin{example}
\label{ex:kr:2}
Consider again the problem described in Example~\ref{ex:kr:1}, where each node has to be colored with exactly one color.
Then, the following program, $\Pi_\mathit{eo}$, can model this scenario:
$$\begin{array}{l}
        \{\mathit{col}(X,C) : \mathit{colour}(C)\} = 1 \text{ :- } \mathit{node}(X). \\ 
        \text{:- }\mathit{col}(X1,C), \mathit{col}(X2,C), \mathit{edge}(X1,X2). \\
        \#\mathit{EOSUM}\{ W: \mathit{col}(X,C), \mathit{cw}(C,W) [X]\} \ge \mathit{M} : \mathit{min}(M). \\
    \end{array}$$
\textcolor{blue}{The first rule explicitly defines the EO group associated with each ground node $(X)$, whereas the \textsc{eosum} constraint exploits these groups to perform the additional weighted-sum propagation.}
\end{example}

\textcolor{blue}{The formal development in Sections~\ref{sec:amosum}--\ref{sec:minimization} concerns lower-bound ($\ge$) constraints. Our implementation additionally supports the dual upper-bound ($\le$) variants. Their propagation uses the minimum possible sum: for each part, it selects the minimum contribution compatible with the current assignment, rather than the maximum one.
It is worth observing that when an at-most-one constraint is combined with a $\le$ condition, the \textsc{amosum} propagator reduces to a standard \textsc{sum} propagator, since the two notions of minimum possible sum coincide.
By contrast, under an exactly-one constraint, the minimum possible sum is necessarily greater than or equal to the one computed by the classical \textsc{sum} propagator.}

Thus, the implementation provides \textsc{amosum} and \textsc{eosum} constructs with both $\ge$ and $\le$. 
\begin{example}[Continuing Example~\ref{ex:kr:2}]
    If we want to enforce that the sum of all colors must be less than or equal to an upper bound, we can introduce the following constraint:
$$\begin{array}{l}
        \#\mathit{EOSUM}\{ W: \mathit{col}(X,C), \mathit{cw}(C,W) [X]\} \le \mathit{M} : \mathit{max}(M). \\
    \end{array}$$
\end{example}

% eq:amosum:kr

\textcolor{blue}{Summing up, with $\#\amosum$ (and its $\#\textsc{eosum}$ counterpart), we define a new constraint construct that succinctly captures at-most-one (or exactly-one) relations between atoms, allowing us to answer question \emph{Q2} affirmatively.}

\subsection{Meta-Encoding Strategies and Integration}

The propagator outlined in Section~\ref{sec:amosum} has been implemented for \textsc{wasp} through its Python interface~\citep{DBLP:journals/tplp/DodaroR20} and for \textsc{clingo}~\citep{DBLP:conf/iclp/GebserKKOSW16} through both its Python and C interfaces.

\textcolor{blue}{The two \textsc{clingo} implementations allow us to compare the performance of the Python and C interfaces. The \textsc{wasp} C counterpart is missing as it currently offers no corresponding external interface.}

In our implementation, \textsc{amosum} constraints are represented by rules. Those rules are created by a \textit{preprocessor} module that takes as input an \textsc{amosum} program and returns a new program where each \textsc{amosum} constraint is replaced by fresh special atoms that are afterward interpreted by the propagator. Specifically, the representation of an \textsc{amosum} constraint $\sigma$ of the form (\ref{eq:amosum:kr}) is the following:
% \#\amosum\{w_1 : \ell_1 \ [s_1];\ \cdots;\ w_n : \ell_n\ [s_n]\} \circ b : a_1, \hdots, a_k, \mathit{not}\, a_{k+1}, \hdots, \mathit{not}\, a_m.
\begin{align}
    % \mathit{group}(\ell_i, w_i, s_i, \sigma) \leftarrow & \qquad\qquad \mathit{for\ all\ } i \in [1..n]\\
    % \mathit{lb}(b, \sigma) \leftarrow
    \label{eq:metaenc:group}
    \{\mathit{\_group\_}(a_i, p(\ell_i), w_i, s_i, \sigma)\}\text{ :- } \ell_i, \textcolor{blue}{\mathit{conj}_i,} \mathit{body}. & \quad \forall i \in [1..n]\\ 
    \label{eq:metaenc:lb}
    \{\mathit{\_lb\_}(b, \sigma)\}\text{ :- } \mathit{body}. & \\
    \label{eq:metaenc:aux}
    \mathit{\_aux\_}(\sigma)\text{ :- } \mathit{body}. & \\
    \label{eq:metaenc:group-aux}
    \{\mathit{\_group\_}(\mathit{\_aux\_}(\sigma), -, b, \textcolor{blue}{\mathit{aux}_\sigma,} \sigma)\}\text{ :- } \mathit{body}. & 
\end{align}
where we assume that $\mathit{\_group\_}$, $\mathit{\_aux\_}$ and $\mathit{\_lb\_}$ do not appear in the program.
Moreover, if $\ell_i$ is a positive atom then $a_i = \ell_i$ and  $p(\ell_i) = ``+" $, otherwise $a_i = \lnot \ell_i$ and $p(\ell_i) = ``-" $.
\
Rule \eqref{eq:metaenc:group} defines a literal $\ell_i$ being in the part $s_i$ with weight $w_i$, in the aggregate $\sigma$.
Rule \eqref{eq:metaenc:lb} defines the (lower) bound
of the constraint.
Rules \eqref{eq:metaenc:aux} and \eqref{eq:metaenc:group-aux} model
the fact that whenever the body is true, the \textsc{amosum} constraint must be fulfilled as well. Intuitively, if body is true then, thanks to \eqref{eq:metaenc:aux}, also the \textit{aux} atom is true. Thanks to rule \eqref{eq:metaenc:group-aux}, the literal $\lnot \mathit{aux}$ is defined inside the \textsc{amosum} aggregate with weight $b$. Given that $\lnot \mathit{aux}$ is false then it `disappears' from the aggregate set and the original \textsc{amosum} constraint must be satisfied.
On the other hand, if \textit{aux} is false then,
thanks to \eqref{eq:metaenc:aux}, also the body is false
and given that $\lnot \mathit{aux}$ is true with weight $b$ then the propagator would not infer anything. Hence, the original \textsc{amosum} constraint could be either true or false.
\textcolor{blue}{
%Concerning 
As for the rules~\eqref{eq:metaenc:group}, \eqref{eq:metaenc:lb}, and \eqref{eq:metaenc:group-aux},
%: these rules 
they use choice rule heads rather than normal rule heads. The reason is that, when the body is undefined, a normal rule would leave the head atom undefined, allowing the solver to choose it as branching literal, incrementing the search space. Instead, by using choice rules, we allow our amosum propagator to derive these atoms as false during its initial preprocessing step reducing the search space, without forcing the body to be false as well.}

\begin{example}[Continuing Example~\ref{ex:kr:1}]
Assume we have three colors (red, green, and blue) and three nodes (1, 2, and 3), that is, the following facts are introduced:
    \begin{align*}
        \begin{array}{ll}
            \mathit{colour}(c). & \textit{for all } c  \in \{r,g,b\}   \\
            \mathit{node}(n). & \textit{for all } n  \in \{1,2,3\}   \\
            \mathit{cw}(c,w). & \textit{for all } (c, w)  \in \{(r,1),(g,2),(b,4) \}   \\
            \mathit{min}(3).
        \end{array}
    \end{align*}
    The constraint
    \begin{align*}
        \#\amosum\{ W: \mathit{col}(X,C), \mathit{cw}(C,W) [X]\} \ge \mathit{M} : \mathit{min}(M).
    \end{align*}
    {\color{blue}
    would be grounded as follows:
    {
    \small
    \begin{align*}
        \#\amosum\{
            1: \mathit{col}(1,r), \mathit{cw}(r,1) [1];
            2: \mathit{col}(1,g), \mathit{cw}(g,2) [1];
            4: \mathit{col}(1,b), \mathit{cw}(b,4) [1];\\
%
            1: \mathit{col}(2,r), \mathit{cw}(r,1) [2];
            2: \mathit{col}(2,g), \mathit{cw}(g,2) [2];
            4: \mathit{col}(2,b), \mathit{cw}(b,4) [2];\\
%
            1: \mathit{col}(3,r), \mathit{cw}(r,1) [3];
            2: \mathit{col}(3,g), \mathit{cw}(g,2) [3];
            4: \mathit{col}(3,b), \mathit{cw}(b,4) [3]\\
        \} \ge \mathit{3} : \mathit{min}(3).
    \end{align*}
    }
    In our implementation, the above ground form is achieved via the following ground rules:
    \begin{align*}
    \begin{array}{rr}
        % \{\mathit{\_group\_}(a_i, p(\ell_i), w, s_i, \sigma)\}\text{:- } a_1, \hdots, a_k, \mathit{not}\, a_{k+1}, \hdots, \mathit{not}\, a_m. & \forall i \in [1..n]\\
        \{\mathit{\_group\_}(\mathit{col}(1, r), + , 1, 1, \sigma)\}\ \text{:- } \mathit{col}(1, r), cw(r,1), \mathit{min}(3).\\
        \{\mathit{\_group\_}(\mathit{col}(2, r), + , 1, 2, \sigma)\}\ \text{:- } \mathit{col}(2, r), cw(r,1), \mathit{min}(3).\\
        \{\mathit{\_group\_}(\mathit{col}(3, r), + , 1, 3, \sigma)\}\ \text{:- } \mathit{col}(3, r), cw(r,1), \mathit{min}(3).\\
        %
        \{\mathit{\_group\_}(\mathit{col}(1, g), + , 2, 1, \sigma)\}\ \text{:- } \mathit{col}(1, g), cw(g,2), \mathit{min}(3).\\
        \{\mathit{\_group\_}(\mathit{col}(2, g), + , 2, 2, \sigma)\}\ \text{:- } \mathit{col}(2, g), cw(g,2), \mathit{min}(3).\\
        \{\mathit{\_group\_}(\mathit{col}(3, g), + , 2, 3, \sigma)\}\ \text{:- } \mathit{col}(3, g), cw(g,2), \mathit{min}(3).\\
        %
        \{\mathit{\_group\_}(\mathit{col}(1, b), + , 4, 1, \sigma)\}\ \text{:- } \mathit{col}(1, b), cw(b,4), \mathit{min}(3).\\
        \{\mathit{\_group\_}(\mathit{col}(2, b), + , 4, 2, \sigma)\}\ \text{:- } \mathit{col}(2, b), cw(b,4), \mathit{min}(3).\\
        \{\mathit{\_group\_}(\mathit{col}(3, b), + , 4, 3, \sigma)\}\ \text{:- } \mathit{col}(3, b), cw(b,4), \mathit{min}(3).\\
        %
        \{\mathit{\_lb\_}(3, \sigma)\}\ \text{:- } \mathit{min}(3). & \\
        %
        \mathit{\_aux\_}(\sigma)\ \text{:- } \mathit{min}(3).& \\
        %
        \{\mathit{\_group\_}(\mathit{\_aux\_}(\sigma), -, 3, \mathit{aux}_\sigma,\sigma)\}\ \text{:- } \mathit{min}(3). & \\
    \end{array}
    \end{align*}
    }
\end{example}

\textcolor{blue}{Moreover, since \textsc{wasp} and \textsc{clingo} already provide efficient propagators for AMO and EO constraints, these relations are explicitly encoded in the input program and enforced by the corresponding built-in propagators. The proposed \textsc{amosum} and \textsc{eosum} propagators exploit the declared groups and focus exclusively on the additional inferences arising from their interaction with the weighted sum and described in Section~\ref{sec:amo:propagator}.}

\subsection{Propagation Activation}
\label{sub:sec:prop:act}
A technical improvement given in this work concerns the invocation of the propagation phase, which plays a crucial role in enhancing the solver's performance.
When the propagation function is called, the algorithm must verify, for each literal, whether it should be propagated. This process generally involves iterating over a set of literals.
Reducing or avoiding such iterations can have a significant impact on the overall efficiency of the algorithm.

\paragraph{Previous Activation}
In \textcolor{blue}{\cite{DBLP:conf/ijcai/AlvianoDFM24},} the propagation phase is invoked only when either the maximum possible sum changes or when the newly inferred literal causes its partition to have only a single remaining undefined literal.
More formally, given a literal $\ell$ and an \textcolor{blue}{assignment} $I$ where $\ell$ is the \textit{last} literal added to $I$, the propagation phase is invoked when $\mps(I \setminus \{\ell\}) \neq \mps(I)$ or $|\lits|_{\mathit{part}(\ell)} \setminus \overline{I}| = 1$.

\paragraph{Improved Activation}
\textcolor{blue}{
%First of all, the 
The first difference w.r.t.~\citep{DBLP:conf/ijcai/AlvianoDFM24} is that we are no longer restricted to the case where $\vert\,\lits|_{\mathit{part}(\ell)} \setminus \overline{I}\, \vert = 1$, that is, we can now infer in early steps.
More specifically, $\ell$ could trigger propagation when it was the \textit{second} max undefined, that is, $\wh(\ell) = \mwh(I \cup \{\overline{m}\}, s)$, where $m \in \{\ell' \in \lits \setminus \{\ell\} \mid \wh(\ell') = \mwh(I,s)\}.$ 
For explaining the second improvement, we consider the concept of \textit{maximum possible decrease (MPD)}, which is the maximum weight of all the literals occurring in \textit{lits}, i.e., $\mathit{MPD} = max\{\wh(\ell') \mid \ell' \in \lits\}$ for \textit{amo} and $\mathit{MPD} = \max_{s \in \parts}\{\max\{\wh(\ell') \mid \ell' \in \lits|_s\} -  \min\{\wh(\ell') \mid \ell' \in \lits|_s\}\}$ for \textit{eo}.
The main idea is as follows.
Let $\ell'$ be the literal with the highest weight. If $\ell'$ becomes false and is the last undefined literal in its part (resp., for \textit{eo}, is the second-to-last undefined literal, with the remaining undefined literal $\ell''$ being the one of minimum weight), then the resulting maximum possible sum is $\mps(I) - \wh(\ell')$ (resp. $\mps(I) - (\wh(\ell') - \wh(\ell''))$).
If $\mps(I) - \mathit{MPD} \ge \bnd$, then no literals can be inferred.
Intuitively, even if the literal with the maximum weight is removed, no propagation can be done as $\mps(I) - \mathit{MPD}$ would still be greater than or equal to the bound $\bnd$.
Therefore, the propagation phase is invoked when $\mps(I \setminus \{\ell\}) \neq \mps(I)$ or $\wh(\ell) = \mwh(I \cup \{\overline{m}\}, s)$ and in both cases $\mps(I) < \bnd + \mathit{MPD}$ must hold.}


\subsection{Laziness}
\label{sub:sec:laziness}
% Often it is more convenient to do not propagate, but delaying it to a second moment.
% It is true especially when the amosum constraint is mostly subsumed from other constraints, that is, the gain of the propagation is not enough to pay the price of the computational time needed to infer those literals.
% To overcome such of cases we implemented a lazy version of the propagator called \textit{amosum-lazy} that propagates to true only when a specific criteria is true.
% Each criteria differentiate various version of the lazy approach. 
% In particular there are two versions of lazyness: \textit{amosum-lazy-h} and 
% \textit{amosum-lazy-f}.

When a propagator is used, it may sometimes be activated to perform a propagation check that ultimately has no effect, that is, no new literals are actually propagated.
This situation typically arises with \amosum constraints involving a large number of literals, where several decisions must be made before any new inference becomes possible. In such cases, the propagator still executes unessential iterations, which can lead to practical inefficiencies.
A common strategy to address this issue is to introduce a degree of laziness in the propagator, so that propagation is triggered only once certain conditions have been met during the search. This idea has also been explored by \cite{DBLP:journals/tplp/CuteriDRS17} (albeit in a different context from the \amosum propagator) where the authors empirically showed, in a controlled setting, that under specific conditions a fully lazy propagator can outperform one that performs propagation at every step.

To address the inefficiencies discussed above, we implemented a \emph{lazy} version of the propagator, called \emph{amosum-lazy}, which activates propagation only when a specific criterion is satisfied.  
The criterion used to trigger propagation determines the particular \emph{degree of laziness}.  
Specifically, we define two variants: \emph{full} (\emph{amosum-lazy-f}) and \emph{hybrid} (\emph{amosum-lazy-h}).

\paragraph{\textcolor{blue}{Parametrized version} (\emph{amosum-lazy-p})}
Both variants can be seen as specific instances of a broader formulation, denoted as \emph{amosum-lazy-p}, where propagation is triggered according to a parameter $p \in [0,1]$.  
This parameter represents a percentage threshold between the constraint bound and the current maximum possible sum.

Formally, let $I$ be an assignment, $\sigma$ be an \amosum constraint, and $\bnd_\sigma$ be its bound.  
Propagation is activated when the following condition holds:
$$
\frac{\bnd_\sigma}{\mps_\sigma(I)} \geq p.
$$
For instance, if $p = 0.7$, propagation will occur only when the bound $\bnd_\sigma$ reaches at least $70\%$ of the current maximum possible sum.  
In this way, both \emph{amosum-lazy-f} and \emph{amosum-lazy-h} correspond to specific choices of the parameter $p$.

\paragraph{Full Version (\emph{amosum-lazy-f})}
The \emph{full} version activates propagation only when the bound is reached or the constraint $\sigma$ becomes violated, that is, when
$$
\bnd_\sigma \geq \mps_\sigma(I)
$$
holds. This case corresponds to \emph{amosum-lazy-p}, with $p = 1$.  
Such a configuration is particularly useful for under-constrained instances or when the constraint $\sigma$ rarely becomes active during the search.

\paragraph{Hybrid Version (\emph{amosum-lazy-h})}
The \emph{hybrid} version lies between the default (non-lazy) propagator and the fully lazy one.  
Here, the degree of laziness is determined by the potential decrease of the maximum possible sum (\emph{lazy MPD}), computed as the difference between the maximum and minimum weights among the literals of $\sigma$.
Propagation is then activated when:
$$
\bnd_\sigma + \mathit{max}_\wh - \mathit{min}_\wh \geq \mps_\sigma(I)
$$
holds. This corresponds to an instance of \emph{amosum-lazy-p} with
$$
p = \frac{\bnd_\sigma}{\bnd_\sigma + \mathit{max}_\wh - \mathit{min}_\wh}\; .
$$
The laziness of this approach stems from the fact that some propagations may be skipped, since the actual safe value to consider, as discussed for the \textit{MPD} in Section~\ref{sub:sec:prop:act}, is equal to $\mathit{max}_\wh$.  
When $\mathit{min}_\wh \gg 0$, the difference between 
$$
p = \frac{\bnd_\sigma}{\bnd_\sigma + \mathit{max}_\wh - \mathit{min}_\wh} \quad \text{and} \quad p' = \frac{\bnd_\sigma}{\bnd_\sigma + \max_\wh}
$$
becomes more pronounced, leading to a higher degree of laziness.  
Note that when $p = \frac{\bnd_\sigma}{\bnd_\sigma + \max_\wh}$, the behavior coincides with the non-lazy (base) configuration.

\section{Experiments}\label{sec:exp}
In this section, we present an empirical evaluation of the proposed propagator by comparing its performance to the state-of-the-art ASP solvers \clingo and \wasp on a set of benchmark instances.  
We first describe the selected benchmarks and the experimental setup, including the systems and configurations used in the evaluation.  
We then analyze the results obtained, highlighting the impact of each modification introduced in this work with respect to~\citep{DBLP:conf/ijcai/AlvianoDFM24} on the overall performance of the solver.

\subsection{Benchmarks}
The experimental analysis is based on a set of benchmarks partly derived from those used by \cite{DBLP:conf/ijcai/AlvianoDFM24} and partly newly introduced ones in order to extend the scope of the evaluation.

The additional benchmarks, namely \emph{Group Assignment}, \emph{Nurse Scheduling}, \textcolor{blue}{\emph{Combinatorial Auctions}, and \textit{Traveling Salesman Problem},} are either directly based on, or inspired by, real-world problems, allowing us to assess the performance of the proposed propagator in more practical settings.

\paragraph{Synthetic benchmark}
This benchmark is designed to empirically assess the properties discussed in Section~\ref{sec:amo:properties}.  
It consists of a synthetic program including a rule that defines the partition and another rule incorporating a \textsc{sum} aggregate.  
The partition is composed of ten parts of uniform size $\mathit{part\_size} \in \{10, 100, 1000\}$, where the $i$-th literal of each part has weight $i$.  
The bound of the aggregate is set to
$\alpha \cdot C_1$ (achievable)
and
$C_1 + \alpha \cdot (C_2 - C_1)$ (unachievable), where
$\alpha \in \{0.15, 0.45, 0.6, 0.9\}$ and
%$C_1$ denotes the sum of the maximum weights in each part, i.e., $10 \cdot \mathit{part\_size}$, and
%$C_2$ denotes the sum of all weights, i.e., $5 \cdot \mathit{part\_size} \cdot (\mathit{part\_size} + 1)$.
$$
C_1 \;=\;10 \cdot \mathit{part\_size} \quad \text{ and }\quad C_2 \;=\; 5 \cdot \mathit{part\_size} \cdot (\mathit{part\_size} + 1),
$$
i.e., $C_1$ denotes the sum of the maximum weights in each part while $C_2$ denotes the sum of all weights.
\textcolor{blue}{Moreover, whenever the computation of a bound yielded a non-integer value, the result was rounded to the nearest integer using Python's built-in \texttt{round} function, which resolves ties by rounding to the nearest even integer.}
Overall, the benchmark comprises 24 instances.

\paragraph{Weighted Graph Coloring (WGC)}%~\label{sec:hard}
%The second benchmark is a modified version of the well-known Graph Coloring problem. 
This benchmark is based on the graph coloring problem discussed in Example~\ref{ex:kr:1}.
Instances are generated from those employed in the ASP competition~\citep{DBLP:journals/ai/CalimeriGMR16}, with the colors red, green, blue, yellow, and cyan associated with the weights 2, 4, 8, 16, and 64, respectively.
For each instance of $n$ nodes, the threshold is set to $\alpha \cdot 64 \cdot n$, where $\alpha \in \{0.15, 0.45, 0.75\}$.
%
%This experiment aims to assess the performance of \textsc{amowasp} in comparison to two leading solvers, namely \textsc{clingo} and \textsc{wasp}, using real-world benchmarks. 


\paragraph{Knapsack (\textcolor{blue}{KS})}%~\label{sec:knapsack}
This benchmark is based on a variant of the classical knapsack problem.  
A set of item types is provided, each associated with a specific weight and value.
Given a knapsack with a fixed capacity and a \textcolor{blue}{minimum} value threshold, the task is to determine whether there exists a selection of arbitrarily many items of each type such that the total weight does not exceed the knapsack capacity while the total value is at least the specified threshold.

The instances are randomly generated according to the following procedure.  
The number of item types, denoted by $n$, ranges from $10$ to $55$ with an increment of $5$.  
For each type, the maximum number of selectable items is fixed to $k = 20$.  
Let $v$ denote the average value of the items.  
Two reference thresholds are defined as:
$$
C_1 := n \cdot v \cdot k \quad \text{and} \quad C_2 := n \cdot v \cdot \frac{k \cdot (k+1)}{2}.
$$

For each value of $n$, ten instances are generated and grouped into four categories:  
\textit{(T1--T3)} having three, three, and two instances, respectively, with the value threshold sampled from a uniform distribution within the intervals $[0, C_1]$, $[0, C_2]$, and $[0.1 \cdot C_1, 1.1 \cdot C_1]$;
\textit{(T4)} having two instances with the value threshold sampled from a normal distribution with mean $C_1$ and variance $5000$.

\paragraph{Group Assignment (GA)}
This benchmark models a simplified personnel allocation problem involving a set of people and a set of projects.  
Each person must be assigned a certain number of working hours to a project in a given month, without exceeding the individual limit of maximum hours ($\max_h$).  
Every person has a role (\textit{full professor}, \textit{associate professor}, or \textit{researcher}), each having a specific hourly salary.  
For each project, a minimum amount of funds must be spent, where the total expenditure corresponds to the sum of salaries multiplied by the assigned working hours.

We consider four categories of instances: \textit{satisfiable (s)}, \textit{unsatisfiable (us)}, \textit{middle (m)}, and \textit{possibly unsatisfiable (pu)}, generated with relative frequencies of $0.4$, $0.1$, $0.1$, and $0.4$, respectively.  
The classification of an instance depends on the required amount of funds to be spent per project.  
Instances with lower funding requirements are more likely to be satisfiable, whereas those with higher requirements tend to be unsatisfiable.

The satisfiable instances are generated with low funding requirements, assuming that each person is assigned to approximately one quarter of the available projects.  
The \textit{middle} instances correspond to boundary cases between satisfiable and unsatisfiable instances, where the funding requirement is computed by assuming that each person works on every project, i.e., for each month, each person works $\frac{1}{n}{\cdot}\max_h$ hours, with $n$ being the number of projects.  
The \textit{possibly unsatisfiable} instances are generated with funding requirements close to the infeasibility threshold, assuming that each person works $90\%$ of the maximum number of hours on each project, i.e., $0.9 \cdot \max_h$.  
Finally, the \textit{unsatisfiable} instances are defined by setting the funding requirements to a clearly infeasible level, assuming that every person works the full $\max_h$ hours on each project.

\paragraph{Nurse Scheduling (NS)}
In this benchmark, we consider the Nurse Scheduling Problem as described by \cite{DBLP:conf/lpnmr/DodaroM17}.
Specifically, it requires assigning nurses to shifts over a planning horizon of one year under multiple operational and regulatory constraints.
In particular, each shift (morning, afternoon, night) has an associated time window and must satisfy minimum and maximum staffing levels per day.
For nurses, the schedule must satisfy bounds on total working hours (per year), ensure legal rest periods (e.g.\ at least 24 hours between shift start times), guarantee a minimum number of rest days in any 14-day window, and enforce that after two consecutive night shifts a special rest day is scheduled. 
Further constraints concern the balance in assignment frequencies: each nurse must work each shift type (morning, afternoon, night) within acceptable bounds, and the distribution of shift types should remain fair across the workforce.
%
The experimental evaluation is based on the five benchmark instances proposed by \cite{DBLP:conf/lpnmr/DodaroM17}.

{\color{blue}
\paragraph{Traveling Salesman Problem (TSP)}
This benchmark models the well-known Traveling Salesman Problem (TSP). Given a graph representing a city map, the goal is to find a tour in which each city is visited at most once (or exactly). This can be modeled by enforcing that each city has \textit{at most one} incoming and \textit{at most one} (\textit{or exactly one}) outgoing street (edge) in the selected tour.
Moreover, each street is associated with a revenue, and the total revenue of the selected tour is required to reach a given bound. Therefore, since the total revenue is expressed as a sum over the selected edges, while for each node \textit{at most one} (or \textit{exactly one}) of its outgoing edges can be selected, the resulting constraints naturally give rise to the AMOSUM structure.
We randomly generated 60 instances with an increasing number of cities. For each instance, we considered five different revenue bounds. Specifically, the bound is defined as $\mps \cdot p$, where $p \in \{0.25, 0.5, 0.75, 1, 1.25\}$ and $\mps$ denotes the value of the maximum possible sum with empty assignment.

\paragraph{Combinatorial Auctions (CA)}
This benchmark models the \emph{Combinatorial Auctions} (CA) problem. In a combinatorial auction, a set of bids is given, where each bid requests a subset of items and is associated with a profit. Two bids are incompatible if they request at least one common item, meaning that they cannot be selected simultaneously. The objective is to select a subset of pairwise compatible bids that maximizes the total profit.
%
In our ASP encoding, each bid is represented by an atom whose truth value indicates whether the corresponding bid is selected. Incompatibilities between bids are encoded through AMO constraints, groups are constructed using the following greedy procedure: starting from an empty collection of disjoint groups, each bid is assigned to the first group in which it is incompatible with every bid already contained in the group; if no such group exists, a new group containing only that bid is created.
%
We use the benchmark instances introduced by~\cite{DBLP:conf/cp/BofillPSV14}, which were generated using the Combinatorial Auctions Test Suite by~\cite{leyton2006test}. The benchmark consists of 170 optimization instances, each containing between 70 and 200 bids. Following the methodology of \cite{DBLP:conf/cp/BofillPSV14}, we derive two decision instances from each optimization instance by considering two lower bounds on the objective value: the best-known optimum and the best-known optimum increased by one. This results in a total of 340 decision instances.

}

\subsection{Systems and Experimental Settings}
As reference solvers, we used the plain version of \textsc{wasp}~\citep{DBLP:conf/lpnmr/AlvianoDLR15} (commit \texttt{f3e4c56}) and \textsc{clingo}~\citep{DBLP:conf/iclp/GebserKKOSW16} (version~5.\textcolor{blue}{8}.0).
The experiments were run on a system with 2.30GHz Intel(R) Xeon(R) Gold 5118 CPU and 512GB of RAM with Ubuntu 20.04.2 LTS (GNU/Linux 5.4.0-137-generic x86\_64).
Time and main memory of each solver run were limited to 20 minutes and 8 GB, respectively.

The proposed propagator has been implemented for both solvers by exploiting their external interfaces.  
Throughout the experimental section, we denote these implementations as \textsc{amowasp} and \textsc{amoclingo}, respectively, and as \textsc{eowasp} and \textsc{eoclingo}, respectively,  if the at-most-one constraint is replaced by an \textcolor{blue}{exactly}-one constraint.  
When referring to a specific configuration or option, we use the notation \textsc{amowasp-option} or \textsc{amoclingo-option}.  
Since multiple variants of the proposed propagator are considered, their detailed descriptions are provided in the corresponding subsections where they are employed.

All the tested systems use the grounder \textsc{gringo} (included in the binary of \textsc{clingo}).

\subsection{Motivating Experiment: Synthetic Benchmark}\label{sec:experiments:syntheticbenchmark}

% \begin{figure*}[t!]
%     \centering
% \begin{tikzpicture}
% 		\begin{axis}[
% 			xtick=data,
% 			xticklabels={$15$, $45$, $60$, $90$, $168$, $302$, $370$, $505$},
% 			ylabel={Conflicts},
%             xlabel={},
%             %xlabel={Asse x: bound/somma totale dei gruppi ($\mathit{bnd}_\sigma/\sum_{s \in S}{\max_{(w : \ell\ [s]) \in \sigma}w}$) per i diversi valori di $|\mathit{lits}_\sigma|_s|$},
%             major x tick style=transparent,
% 			ybar,
% 			ymode=log,
% 			ymin=0.1,
% 			ymax=10000000,
% 			ytick={1, 10, 100, 1000, 10000, 100000, 1000000, 10000000},
% 			yticklabels={0, $10^1$, $10^2$, $10^3$, $10^4$, $10^5$, $10^6$, $10^7$},
% 			x axis line style={opacity=0},
%             x tick label style={rotate=90, anchor=center},
% 			xtick=data,
% 			ymajorgrids=true,
% 			bar width=1pt,
% 			grid style=dotted,
% 			nodes near coords,
% 			scale only axis,
% 			point meta=explicit symbolic,
% 			% legend columns=-1,
% 			% legend pos=north east,
% 			height=0.2\textwidth,
% 			width=0.27\textwidth,
% 			name=plot1,
% 			title={$\mathit{part\_size} = 10$}
% 			]
			
% 			\addplot  coordinates {(0,1) (1,1) (2,27) (3, 2453921) (4, 3207378) (5, 112327) (6, 4845) (7, 22)};
%             \addplot  coordinates {(0,1) (1,1506) (2,70594) (3, 194496) (4, 228698) (5, 58368) (6, 2492) (7, 8)};
% 			\legend{}
% 		\end{axis}	
% 		\begin{axis}[
% 			xtick=data,
% 			xticklabels={$150$, $450$, $600$, $900$, $1680$, $3020$, $3700$, $5050$},
% 			ylabel={},
%             xlabel={Bound},
%             %xlabel={Asse x: bound/somma totale dei gruppi ($\mathit{bnd}_\sigma/\sum_{s \in S}{\max_{(w : \ell\ [s]) \in \sigma}w}$) per i diversi valori di $|\mathit{lits}_\sigma|_s|$},
%             major x tick style=transparent,
% 			ybar,
% 			ymode=log,
% 			ymin=0.1,
% 			ymax=10000000,
% 			ytick={1, 10, 100, 1000, 10000, 100000, 1000000, 10000000},
% 			yticklabels={},
% 			x axis line style={opacity=0},
%             x tick label style={rotate=90, anchor=center},			
% 			xtick=data,
% 			ymajorgrids=true,
% 			bar width=1pt,
% 			grid style=dotted,
% 			nodes near coords,
% 			scale only axis,
% 			point meta=explicit symbolic,	
%             legend columns=-1,
% 			% legend pos=north,
%             legend style={at={(0.5,1.03)},anchor=north},
%             height=0.2\textwidth,
% 			width=0.27\textwidth,
% 			name=plot2,
%             at=(plot1.right of south east), anchor=left of south west,
% 			title={$\mathit{part\_size} = 100$}
% 			]
			
% 			\addplot  coordinates {(0,74) (1,124) (2,232) (3, 47209) (4, 263569) (5, 370134) (6, 476627) (7, 64)};
%             \addplot  coordinates {(0,549) (1,100607) (2,95725) (3, 71962) (4, 53629) (5, 76733) (6, 152531) (7, 98)};	
%             \legend{\textsc{clingo},\textsc{wasp}}
% 		\end{axis}
%         \begin{axis}[
% 			xtick=data,
% 			xticklabels={$1500$, $4500$, $6000$, $9000$, $16800$, $30200$, $37000$, $50500$},
% 			ylabel={},
%             xlabel={},
%             %xlabel={Asse x: bound/somma totale dei gruppi ($\mathit{bnd}_\sigma/\sum_{s \in S}{\max_{(w : \ell\ [s]) \in \sigma}w}$) per i diversi valori di $|\mathit{lits}_\sigma|_s|$},
%             major x tick style=transparent,
% 			ybar,
% 			ymode=log,
% 			ymin=0.1,
% 			ymax=10000000,
% 			ytick={1, 10, 100, 1000, 10000, 100000, 1000000, 10000000},
% 			yticklabels={},
% 			x axis line style={opacity=0},
% 			x tick label style={rotate=90, anchor=center},
% 			xtick=data,
%             ymajorgrids=true,
% 			bar width=1pt,
% 			grid style=dotted,
% 			nodes near coords,
% 			scale only axis,
% 			point meta=explicit symbolic,            
% 			height=0.2\textwidth,
% 			width=0.27\textwidth,
% 			name=plot3,
%             at=(plot2.right of south east), anchor=left of south west,
% 			title={$\mathit{part\_size} = 1000$}            
% 			]
			
% 			\addplot  coordinates {(0,6397) (1,8907) (2,8253) (3, 7461) (4, 59024) (5, 22660) (6, 33102) (7, 855)};
%             \addplot  coordinates {(0,6) (1,3004) (2,4112) (3, 44378) (4, 10811) (5, 15519) (6, 25730) (7, 998)};									
% 		\end{axis}        
% 	\end{tikzpicture}
%     \caption{
%         Number of conflicts on synthetic \textsc{amosum} constraints comprising 10 parts of varying size and bound.
%         For each $\mathit{part\_size}$ there are four satisfiable instances (the first four bounds) and four unsatisfiable instances (the last four bounds).
%         \textsc{amowasp} and \textsc{amoclingo} are not reported in the plots since they solve all tested instances in this benchmark without raising any conflict.
%     }
%     \label{fig:resultssynthetic}
% \end{figure*}



% \begin{table*}[t!]
%     \centering
%     \small
%     \setlength{\tabcolsep}{4pt}
%     \begin{tabular}{rrr r rrr r rrr}
%         \toprule
%         \multicolumn{3}{c}{$\mathit{part\_size}=10$} &
%         \multicolumn{3}{c}{$\mathit{part\_size}=100$} &
%         \multicolumn{3}{c}{$\mathit{part\_size}=1000$} \\
%         \cmidrule(lr){1-3}
%         \cmidrule(lr){4-6}
%         \cmidrule(lr){7-9}
%         Bound & \textsc{clingo} & \textsc{wasp}
%         & Bound & \textsc{clingo} & \textsc{wasp}
%         & Bound & \textsc{clingo} & \textsc{wasp} \\
%         \midrule
%         15 & XXX         & XXX
%         & 150 & XXX        & XXX
%         & 1500 & XXX  & XXX \\

%         45 & XXX         & XXX
%         & 450 & XXX       & XXX
%         & 4,500 & XXX  & XXX \\

%         60        & XXX & XXX
%         & 600 & XXX       & XXX
%         & 6,000 & XXX & XXX \\

%         90 & XXX & XXX
%         & 900 & XXX    & XXX
%         & 9,000 & XXX  & XXX \\
%         \midrule
%         168 & XXX & XXX
%         & 8,425 & XXX & XXX
%         & 759,250 & XXX & XXX \\

%         302 & XXX   & XXX
%         & 23,275 & XXX & XXX
%         & 2,257,750 & XXX & XXX \\

%         370 & XXX     & XXX
%         & 30,700 & XXX & XXX
%         & 3,007,000 & XXX & XXX \\

%         505 & XXX        & XXX
%         & 45,550 & XXX      & XXX
%         & 4,505,500 & XXX    & XXX \\
%         \bottomrule
%     \end{tabular}
%     \caption{
%         Number of conflicts on synthetic \textsc{amosum} constraints
%         comprising 10 parts of varying size and bound.
%         For each value of $\mathit{part\_size}$, the first four instances
%         are satisfiable and the last four are unsatisfiable.
%         \textsc{amowasp} and \textsc{amoclingo} are not reported because
%         they solve all tested instances without raising any conflict.
%     }
%     \label{tab:resultssynthetic}
% \end{table*}



\begin{table*}[t!]
    \centering
    \small
    \setlength{\tabcolsep}{4pt}
    
\begin{tabular}{rrr r rrr r rrr}
        \toprule
        \multicolumn{3}{c}{$\mathit{part\_size}=10$} &
        \multicolumn{3}{c}{$\mathit{part\_size}=100$} &
        \multicolumn{3}{c}{$\mathit{part\_size}=1000$} \\
        \cmidrule(lr){1-3}
        \cmidrule(lr){4-6}
        \cmidrule(lr){7-9}
        Bound & \textsc{clingo} & \textsc{wasp}
        & Bound & \textsc{clingo} & \textsc{wasp}
        & Bound & \textsc{clingo} & \textsc{wasp} \\
        \midrule
        15 & 0         & 0
        & 150 & 74        & 549
        & 1500 & 6  & 6,397 \\

        45 & 0         & 1,506
        & 450 & 124       & 102,882
        & 4,500 & 3,004  & 9,182 \\

        60        & 27 & 70,594
        & 600 & 232       & 96,646
        & 6,000 & 4,112 & 8,506 \\

        90 & 5,097,979 & 220,209
        & 900 & 47,209    & 71,775
        & 9,000 & 98,938  & 7,672 \\
        \midrule
        168 & 5,929,256 & 240,305
        & 8,425 & 483,020 & 54,279
        & 759,250 & 113,820 & 11,253 \\

        302 & 112,327   & 58,368
        & 23,275 & 679,894 & 78,014
        & 2,257,750 & 27,197 & 16,347 \\

        370 & 4,845     & 2,492
        & 30,700 & 800,411 & 162,970
        & 3,007,000 & 74,180 & 27,791 \\

        505 & 22        & 8
        & 45,550 & 64      & 98
        & 4,505,500 & 855    & 998 \\
        \bottomrule
    \end{tabular}


    \caption{
        Number of conflicts on synthetic \textsc{amosum} constraints
        comprising 10 parts of varying size and bound.
        For each value of $\mathit{part\_size}$, the first four instances
        are satisfiable and the last four are unsatisfiable.
        \textsc{amowasp} and \textsc{amoclingo} are not reported because
        they solve all tested instances without raising any conflict.
    }
    \label{tab:resultssynthetic}
\end{table*}


We first start with the analysis of the synthetic benchmark.
First of all, we observe that all instances used in this experiment are trivial for \textsc{amowasp}, as they can be solved without raising any conflict.  
In contrast, traditional solvers such as \textsc{wasp} and \textsc{clingo} cannot perform the same inferences (see Theorem~\ref{thm:inferences}).  
We use this benchmark as a motivating example to illustrate the behavior of these two solvers on such cases, and we measure the number of conflicts generated by each system within the first minute of computation.

The experimental results are summarized in Table~\ref{tab:resultssynthetic}.
In particular, we can see that when $\mathit{part\_size}$ is 10 or 100, \textsc{clingo} can handle the first three instances in an efficient way, solving them with only a few conflicts. However, this is not true when $\mathit{part\_size}$ is $1000$. This suggests that the way in which \textsc{clingo} makes decisions (i.e., its branching heuristic) works better for instances with \textsc{sum} aggregates of small to medium size.
Moreover, as we expected, the number of conflicts is related to the bound. Specifically, instances with the smallest and largest bounds have fewer conflicts. This happens because these instances are either not constrained enough or overly constrained.
Note that our propagator can solve 
such instances without raising any conflict.

% \begin{table}[t!]
%   \centering    
%   \small
% \begin{tabular}{llrrrrr}
% \toprule
% &  & \textbf{WGC} & \textbf{GA} & \textbf{KS} & \textbf{NS} & \textbf{Total}\\
% \midrule
% & \#Instances & 300 & 350 & 100 & 5 & 755\\
% \midrule
% \multirow{3}{*}{\rotatebox{90}{\textsc{amo}}}
%  & \amowaspij & 64 & 65 & 71 & 0 & 200\\
%  & \amoclingoijpy & 95 & 141 & 75 & 0 & 311\\
%  & \amoclingoijc & \textbf{102} & \textbf{307} & \textbf{77} & \textbf{2} & \textbf{488}\\
% \midrule
% \multirow{3}{*}{\rotatebox{90}{\textsc{eo}}} 
%  & \eowaspij & 50 & 130 & 73 & 0 & 253\\
%  & \eoclingoijpy & 79 & 153 & 74 & 0 & 306\\
%  & \eoclingoijc & \textbf{100} & \textbf{311} & \textbf{76} & \textbf{1} & \textbf{488}\\
% \bottomrule
% \end{tabular}
% \caption{Comparison between baseline version \amowaspij\ and \clingo\ versions, namely \amoclingoijpy\ and \amoclingoijc. Best results are highlighted in bold.}
% \label{tab:IJCAI:baseline}
% \end{table}

\begin{table}[t!]
  \centering    
  \small




    \begin{tabular}{llrrrrrrr}
    \toprule
    &  & \textbf{WGC} & \textbf{GA} & \textbf{KS} & \textbf{NS} & \textbf{TSP} & \textbf{CA} & \textbf{Total}\\
    \midrule
    & \#Instances & 300 & 350 & 100 & 5 & 300 & 340 & 1395\\
    \midrule
    \multirow{3}{*}{\rotatebox{90}{\textsc{amo}}}
     & \amoclingoijc & \textbf{102} & \textbf{307} & \textbf{77} & \textbf{1} & \textbf{253} & \textbf{307} & \textbf{1047}\\
 & \amoclingoijpy & 95 & 141 & 75 & 0 & 252 & 285 & 848\\
 & \amowaspij & 64 & 65 & 71 & 0 & 131 & 306 & 637\\
    \midrule
    & \#Instances & 300 & 350 & 100 & 5 & 300 & 0 & 1055\\
    \midrule
    \multirow{3}{*}{\rotatebox{90}{\textsc{eo}}} 
     & \eoclingoijc & \textbf{101} & \textbf{350} & \textbf{77} & \textbf{1} & \textbf{256} & - & \textbf{785}\\
 & \eoclingoijpy & 79 & 280 & 74 & 0 & 253 & - & 686\\
 & \eowaspij & 50 & 258 & 73 & 0 & 250 & - & 631\\
    \bottomrule
    \end{tabular}
    

\caption{Comparison between baseline version \amowaspij\ and \clingo\ versions, namely \amoclingoijpy\ and \amoclingoijc. Best results are highlighted in bold.}
\label{tab:IJCAI:baseline}
\end{table}

\subsection{Comparison to \textsc{amowasp} by \cite{DBLP:conf/ijcai/AlvianoDFM24}}

\textcolor{blue}{We first compare the enhanced propagator introduced in this work with respect to the version presented by \cite{DBLP:conf/ijcai/AlvianoDFM24}.  
It is important to note that in the latter work the propagator was integrated into \textsc{wasp} through the  Python interface of the system~\citep{DBLP:journals/tplp/DodaroR20}.  
In the following, we refer to that implementation as \amowaspij.}

Our first improvement concerns the adoption of \textsc{clingo}, where the propagator has been implemented using its external Python (referred to as \amoclingoijpy) and C interfaces (referred to as \amoclingoijc). \textcolor{blue}{It is worth stressing that \amoclingoijpy and \amoclingoijc differ from the propagator proposed by~\cite{DBLP:conf/ijcai/AlvianoDFM24} \emph{only}  in the solver used and the programming language adopted (in case of \amoclingoijc). }

The comparison is reported in Table \ref{tab:IJCAI:baseline}.
First, we observe that \amoclingoijpy solves \textcolor{blue}{211} instances more than \amowaspij, and that \eoclingoijpy solves \textcolor{blue}{55} instances more than \eowaspij.  
These results can be explained by two main observations:  
\textit{(i)} \textsc{clingo} is generally faster than \textsc{wasp} as a baseline solver; and  
\textit{(ii)} as discussed for the synthetic benchmark in Section~\ref{sec:experiments:syntheticbenchmark}, \textsc{clingo} produces fewer conflicts, suggesting that its heuristic strategy is better suited for this class of problems.

Moreover, it can be observed that \amoclingoijc and \eoclingoijc are significantly faster than their Python counterparts, solving \textcolor{blue}{199 and 99} additional instances, respectively.

Overall, compared to the version presented by \cite{DBLP:conf/ijcai/AlvianoDFM24}, we obtain a substantial improvement: \textcolor{blue}{410} additional instances solved on problems involving \amosum\ constraints (\amoclingoijc vs.\ \amowaspij) and \textcolor{blue}{154} additional instances solved on problems with \textsc{eosum} constraints (\eoclingoijc vs. \eowaspij), which amounts to  %corresponding 
improvements of \textcolor{blue}{64\% and 24\%}, respectively.

\textcolor{blue}{Since the best-performing configuration in the base version is the one based on \textsc{clingo} with the C interface, the remaining improvements are based on \amoclingoijc.}

% \begin{table}[t!]
%   \centering    
%   \small
% \begin{tabular}{llrrrrr}
% \toprule
% &  & \textbf{WGC} & \textbf{GA} & \textbf{KS} & \textbf{NS} & \textbf{Total}\\
% \midrule
% \multirow{2}{*}{\rotatebox{90}{\textsc{amo}}}
%  & \amoclingoijc & 102 & 307 & 77 & 2 & 488\\
%  & \amoclingoinf & \textbf{120} & 307 & 77 & 2 & \textbf{506}\\
% \midrule
% \multirow{2}{*}{\rotatebox{90}{\textsc{eo}}} 
%  & \eoclingoijc & 100 & 311 & 76 & 1 & 488\\
%  & \eoclingoinf & \textbf{153} & \textbf{314} & \textbf{78} & 1 & \textbf{546}\\
% \bottomrule
% \end{tabular}
% \caption{Comparison between version \amoclingoijc and \amoclingoinf to show the impact of the inference rule \eqref{eq:amosum:propagation:true}. Best results are highlighted in bold.}
\begin{table}[t!]
  \centering    
  \small
    \begin{tabular}{llrrrrrrr}
    \toprule
    &  & \textbf{WGC} & \textbf{GA} & \textbf{KS} & \textbf{NS} & \textbf{TSP} & \textbf{CA} & \textbf{Total}\\
    \midrule
& \#Instances & 300 & 350 & 100 & 5 & 300 & 340 & 1395\\
    \midrule
    \multirow{2}{*}{\rotatebox{90}{\textsc{amo}}}
     & \amoclingoijc  & 102 & \textbf{307} & \textbf{77} & 1 & 253 & 307 & 1047\\
 & \amoclingoinf & \textbf{120} & \textbf{307} & \textbf{77} & \textbf{2} & \textbf{256} & \textbf{325} & \textbf{1087}\\
 \midrule
 & \#Instances & 300 & 350 & 100 & 5 & 300 & 0 & 1055\\
    \midrule
    \multirow{2}{*}{\rotatebox{90}{\textsc{eo}}} 
 & \eoclingoijc & 101 & \textbf{350} & \textbf{77} & \textbf{1} & 256 & - & 785\\
     & \eoclingoinf & \textbf{153} & \textbf{350} & \textbf{77} & \textbf{1} & \textbf{258} & - & \textbf{839}\\
    \bottomrule
    \end{tabular}
\caption{Comparison between version \amoclingoijc and \amoclingoinf to show the impact of the inference rule \eqref{eq:amosum:propagation:true}. Best results are highlighted in bold.}
\label{tab:inference}
\end{table}
\begin{table}[t!]
  \centering    
  \small
    \begin{tabular}{llrrrrrrr}
    \toprule
    &  & \textbf{WGC} & \textbf{GA} & \textbf{KS} & \textbf{NS} & \textbf{TSP} & \textbf{CA} & \textbf{Total}\\
    \midrule
    & \#Instances & 300 & 350 & 100 & 5 & 300 & 340 & 1395\\
    \midrule
    \multirow{2}{*}{\rotatebox{90}{\textsc{amo}}}
     & \amoclingoinf & 120 & 307 & \textbf{77} & \textbf{2} & \textbf{256} & \textbf{325} & 1087\\
 & \amoclingoinfr & \textbf{122} & \textbf{309} & 76 & \textbf{2} & 255 & 324 & \textbf{1088}\\
    \midrule
    & \#Instances & 300 & 350 & 100 & 5 & 300 & 0 & 1055\\
    \midrule
    \multirow{2}{*}{\rotatebox{90}{\textsc{eo}}} 
     & \eoclingoinf & 153 & \textbf{350} & \textbf{77} & \textbf{1} & \textbf{258} & - & 839\\
 & \eoclingoinfr & \textbf{157} & \textbf{350} & \textbf{77} & \textbf{1} & 256 & - & \textbf{841}\\
    \bottomrule
    \end{tabular}
\caption{Comparison between version \amoclingoinf\ and \amoclingoinfr\ to show the impact of the reason \eqref{eq:amosum:reason:rt}. Best results are highlighted in bold.}
\label{tab:inference-reason}
\end{table}



% \begin{table}[t!]
%   \centering    
%   \small
% \begin{tabular}{llrrrrr}
% \toprule
% &  & \textbf{WGC} & \textbf{GA} & \textbf{KS} & \textbf{NS} & \textbf{Total}\\
% \midrule
% \multirow{2}{*}{\rotatebox{90}{\textsc{amo}}}
%  & \amoclingoinfr & 120 & 307 & 77 & 2 & 506\\
%  & \amoclingoinfmr & \textbf{184} & \textbf{328} & \textbf{86} & \textbf{3} & \textbf{601}\\
% \midrule
% \multirow{2}{*}{\rotatebox{90}{\textsc{eo}}} 
%  & \eoclingoinfr & 157 & \textbf{314} & 78 & \textbf{1} & \textbf{550}\\
%  & \eoclingoinfmr & \textbf{160} & 303 & \textbf{82} & 0 & 545\\
% \bottomrule
% \end{tabular}
% \caption{Comparison between version \amoclingoinfr\ and \amoclingoinfmr\ to show the impact of the minimal reason. Best results are highlighted in bold.}
% \label{tab:minimal:reason}
% \end{table}


\begin{table}[t!]
  \centering    
  \small
    \begin{tabular}{llrrrrrrr}
    \toprule
    &  & \textbf{WGC} & \textbf{GA} & \textbf{KS} & \textbf{NS} & \textbf{TSP} & \textbf{CA} & \textbf{Total}\\
    \midrule
& \#Instances & 300 & 350 & 100 & 5 & 300 & 340 & 1395\\
    \midrule
    \multirow{2}{*}{\rotatebox{90}{\textsc{amo}}}
 & \amoclingoinfr & 122 & \textbf{309} & \textbf{76} & \textbf{2} & \textbf{255} & 324 & 1088\\
     & \amoclingoinfmr & \textbf{127} & 305 & 75 & 1 & \textbf{255} & \textbf{339} & \textbf{1102}\\
    \midrule
    & \#Instances & 300 & 350 & 100 & 5 & 300 & 0 & 1055\\
    \midrule
    \multirow{2}{*}{\rotatebox{90}{\textsc{eo}}}     
     & \eoclingoinfr & \textbf{157} & \textbf{350} & \textbf{77} & \textbf{1} & \textbf{256} & - & \textbf{841}\\
 & \eoclingoinfmr & 144 & \textbf{350} & 74 & 0 & 254 & - & 822\\
    \bottomrule
    \end{tabular}
\caption{Comparison between \amoclingoinfr\ and \amoclingoinfmr\, showing the impact of subset-minimal reasons. Best results are highlighted in bold.}
\label{tab:minimal:reason}
\end{table}


\subsection{Impact of Inference, (Minimal) Reason, and Lazy Propagation Improvements}
\textcolor{blue}{In the following, we analyze the individual contribution of the improvements introduced \textcolor{blue}{with respect to~\cite{DBLP:conf/ijcai/AlvianoDFM24}.} In particular, we evaluate the inference rules and the computation of (minimal) reasons, as well as the impact of lazy propagation.  }
The goal is to isolate and quantify the effect of each contribution on solver performance.  
To this end, we consider different configurations of the propagator, each enabling or disabling specific features related to inference and reason minimization.  

\paragraph{Inference Improvement}
First of all, we show the impact of the theoretical and practical improvements discussed in Section~\ref{sub:sec:prop:act}.
These enhancements concern the earlier inference of true literals through the inference rule~\eqref{eq:amosum:propagation:true}, together with an improved activation mechanism for the propagation.  
We refer to this version as \amoclingoinf.

%
Table~\ref{tab:inference} presents a comparison between \amoclingoinf and the previous version \amoclingoijc, together with their \textsc{eo} counterparts (\eoclingoijc).  
It should be noted that for \eoclingoijc, only the improvement concerning the activation of the propagation is applicable, since  as discussed in Section~\ref{sub:sec:asp:const} the \textsc{eo} propagator does not implement inference rules for true literals.
%
The results indicate that improving the inference mechanism alone leads to a consistent gain in performance, with the number of solved instances increasing \textcolor{blue}{from 1047 to 1087 for \textsc{amo}, and from 785 to 839 for \textsc{EO}, corresponding to improvements of approximately $3.8\%$ and $6.9\%$, respectively.}

\paragraph{Reason Improvement}
Continuing our investigation of the proposed enhancements, we introduce a new version, denoted as \amoclingoinfr, which combines the inference improvement discussed above with reason~\eqref{eq:amosum:reason:rt}.
\textcolor{blue}{In particular, let $\ell$ be the literal inferred by (RT)~\eqref{eq:amosum:propagation:true}. The reason proposed by~\cite{DBLP:conf/ijcai/AlvianoDFM24} differs only for the part of $\ell$ where all $\lits_s\setminus \{\ell\}$ were added to the reason.}
The results, reported in Table~\ref{tab:inference-reason}, show that this modification provides only a marginal benefit. Indeed, concerning \textsc{amo}, we observe an increase of one solved instance \textcolor{blue}{overall, due to an improvement} in both the WGC and GA benchmarks, while for \textsc{eo}, we observe an overall improvement of \textcolor{blue}{2} instances, due to a better performance in the benchmark WGC.
This limited improvement can be explained by the fact that the two reasons are in practice very similar, differing at most in the best case for all literals within a single part.

\paragraph{\textcolor{blue}{Subset-}Minimal Reason Improvement}
To further refine the reasoning mechanism, we move from reason~\eqref{eq:amosum:reason:rt} \textcolor{blue}{to reason created by Algorithm~\ref{alg:minfly} (\textit{MRF})}, thus defining a new version, \amoclingoinfmr.
Table~\ref{tab:minimal:reason} summarizes the results obtained for this configuration.  
For the \textsc{amo} case, we observe an overall improvement of \textcolor{blue}{14} additional instances solved compared to \amoclingoinfr, confirming the effectiveness of the proposed approach.  
Conversely, in the \textsc{eo} case, we observed a different behavior, with an overall decrease of 19 solved instances.  
\textcolor{blue}{Since neither approach consistently outperforms the other across all instances, we consider both to be valid and competitive alternatives. Accordingly, in the following experiments, we evaluate the subsequent extensions on both approaches.}

\begin{table}[t!]
  \centering    
  \small
    \begin{tabular}{llrrrrrrr}
    \toprule
    &  & \textbf{WGC} & \textbf{GA} & \textbf{KS} & \textbf{NS} & \textbf{TSP} & \textbf{CA} & \textbf{Total}\\
    \midrule
& \#Instances & 300 & 350 & 100 & 5 & 300 & 340 & 1395\\
    \midrule
    \multirow{6}{*}{\rotatebox{90}{\textsc{amo}}}
 &  \amoclingoinfr & 122 & 309 & 76 & 2 & 255 & 324 & 1088\\
 & \amoclingoinfrhybrid & 122 & 308 & 76 & 2 & 255 & 323 & 1086\\
 & \amoclingoinfrlazy & 165 & 313 & \textbf{77} & \textbf{3} & 274 & 330 & 1162\\
     & \amoclingoinfmr & 127 & 305 & 75 & 1 & 255 & \textbf{339} & 1102\\
 & \amoclingoinfmrhybrid & 126 & 309 & 75 & 1 & 255 & 338 & 1104\\
 & \amoclingoinfmrlazy & \textbf{183} & \textbf{314} & 75 & 1 & \textbf{275} & \textbf{339} & \textbf{1187}\\
    \midrule
    & \#Instances & 300 & 350 & 100 & 5 & 300 & 0 & 1055\\
    \midrule
    \multirow{6}{*}{\rotatebox{90}{\textsc{eo}}} 
     & \eoclingoinfr& 157 & \textbf{350} & \textbf{77} & 1 & 256 & - & 841\\
 & \eoclingoinfrhybrid & 160 & \textbf{350} & \textbf{77} & 1 & 257 & - & 845\\
 & \eoclingoinfrlazy & \textbf{172} & \textbf{350} & \textbf{77} & \textbf{3} & \textbf{274} & - & \textbf{876}\\
 & \eoclingoinfmr& 144 & \textbf{350} & 74 & 0 & 254 & - & 822\\
 & \eoclingoinfmrhybrid & 145 & \textbf{350} & 74 & 0 & 255 & - & 824\\
 & \eoclingoinfmrlazy & 167 & \textbf{350} & 75 & 2 & 256 & - & 850\\
    \bottomrule
    \end{tabular}
\caption{Comparison between lazy versions against no-lazy ones\ to show the impact of the lazy approach. Best results are highlighted in bold.}
\label{tab:laziness}
\end{table}


\paragraph{Lazy Propagation}
{\color{blue}
In the following, we evaluate the impact of lazy propagation on \amoclingoinfr, \amoclingoinfmr, \eoclingoinfr, and \eoclingoinfmr. For each configuration, the suffixes \textsc{-h} and \textsc{-l} denote the hybrid and fully lazy variants, respectively. 
Table~\ref{tab:laziness} reports the results of their comparison.
%
The results show that the fully lazy approach consistently improves performance across all configurations. In particular, \amoclingoinfrlazy and \amoclingoinfmrlazy solve 74 and 85 more instances than \amoclingoinfr and \amoclingoinfmr, respectively. Similarly, \eoclingoinfrlazy and \eoclingoinfmrlazy solve 35 and 28 additional instances compared to their non-lazy counterparts.

By contrast, the hybrid approach has only a marginal and inconsistent impact: relative to the corresponding non-lazy configurations, the number of solved instances ranges from two fewer (\amoclingoinfrhybrid versus \amoclingoinfr) to four additional instances (\eoclingoinfrhybrid versus \eoclingoinfr).
}

A possible explanation lies in the fact that, in practice, the value of \textit{MPD} used in the non-lazy approach (see Section~\ref{sub:sec:prop:act}) and that is employed in the hybrid version (see Section~\ref{sub:sec:laziness}) tend to coincide, thus limiting the benefit of introducing partial laziness.

\subsection{Comparison to \textsc{clingo} and \textsc{wasp}}

% \begin{table}[t!]
%   \centering    
%   \small
% \begin{tabular}{llrrrrr}
% \toprule
% &  & \textbf{WGC} & \textbf{GA} & \textbf{KS} & \textbf{NS} & \textbf{Total}\\
% \midrule
% \multirow{6}{*}{\rotatebox{90}{\textsc{amo}}}
%     & \textbf{\#Instances} & 300 & 350 & 100 & 5 & 755\\
% \midrule
%     & \wasp & 60 & 57 & 14 & 0 & 131\\ 
%     & \clingo & 60 & 238 & 48 & \textbf{4}& 350\\
%     & \amowaspij & 64 & 65 & 71 & 0 & 200\\
%     & \amoclingoinfmr & 184 & \textbf{328} & 86 & 3 & 601\\
%     & \amoclingolazy & \textbf{228} & 323 & \textbf{88} & \textbf{4} & \textbf{643}\\

% \midrule
% \multirow{6}{*}{\rotatebox{90}{\textsc{eo}}} 
% & \wasp & 32 & 39 & 12 & 0 & 83\\
% & \clingo & 95 & 280 & 48 & \textbf{4} & 427\\
% & \eowaspij & 50 & 130 & 73 & 0 & 253\\
% & \eoclingoinfmr & 160 & 303 & 82 & 0 & 545\\
% & \eoclingolazy & \textbf{188} & \textbf{316} & \textbf{87}   & 1 & \textbf{592}\\
% \bottomrule
% \end{tabular}
% \caption{Comparison with state-of-the-art solvers \clingo\ and \wasp\ with baseline \amowaspij\ and best no-lazy and lazy versions, \amoclingoinfmr\ and \amoclingolazy, respectively. Best results are highlighted in bold.}
% \label{tab:overall}
% \end{table}


\begin{table}[t!]
  \centering    
  \small



    \begin{tabular}{llrrrrrrr}
    \toprule
    &  & \textbf{WGC} & \textbf{GA} & \textbf{KS} & \textbf{NS} & \textbf{TSP} & \textbf{CA} & \textbf{Total}\\
    \midrule
& \#Instances & 300 & 350 & 100 & 5 & 300 & 340 & 1395\\
    \midrule
    \multirow{5}{*}{\rotatebox{90}{\textsc{amo}}}
 & \amoclingoinfrlazy & 165 & 313 & \textbf{77} & 3 & 274 & 330 & 1162\\
     & \amoclingoinfmrlazy & \textbf{183} & \textbf{314} & 75 & 1 & \textbf{275} & \textbf{339} & \textbf{1187}\\
 & \amowaspij & 64 & 65 & 71 & 0 & 131 & 306 & 637\\
 & \clingo & 60 & 238 & 48 & \textbf{4} & 128 & 252 & 730\\
 & \wasp & 60 & 57 & 14 & 0 & 125 & 173 & 429\\
    \midrule
    & \#Instances & 300 & 350 & 100 & 5 & 300 & 0 & 1055\\
    \midrule
    \multirow{5}{*}{\rotatebox{90}{\textsc{eo}}} 
     & \eoclingoinfrlazy& \textbf{172} & \textbf{350} & \textbf{77} & 3 & \textbf{274} & - & \textbf{876}\\
 & \eoclingoinfmrlazy & 167 & \textbf{350} & 75 & 2 & 256 & - & 850\\
 & \eowaspij & 50 & 258 & 73 & 0 & 250 & - & 631\\
 & \clingo & 95 & \textbf{350} & 47 & \textbf{4} & 129 & - & 625\\
 & \wasp & 32 & 139 & 16 & 0 & 124 & - & 311\\
    \bottomrule
    \end{tabular}
\caption{Comparison with state-of-the-art solvers \clingo\ and \wasp\ with baseline \amowaspij\ and best lazy versions. Best results are highlighted in bold.}
\label{tab:overall}
\end{table}

In this section, we compare our best configuration against the state-of-the-art solvers \textsc{clingo} and \textsc{wasp}, as well as the baseline proposed by \cite{DBLP:conf/ijcai/AlvianoDFM24}; the results are reported in Table~\ref{tab:overall}.

\textcolor{blue}{For both \textsc{amo} and \textsc{eo} constraints, our approaches obtain the best overall results in terms of solved instances. In the \textsc{amo} setting, \amoclingoinfmrlazy \textcolor{blue}{solves 1187 out of 1395 instances}. This corresponds to \textcolor{blue}{457} additional instances compared to plain \textsc{clingo}, representing a 62.6\% improvement. Similarly, in the \textsc{eo} setting, \eoclingoinfrlazy obtains the best result by solving 876 instances—251 more than plain \textsc{clingo}, corresponding to a 40.2\% improvement.}

\textcolor{blue}{Looking at the benchmark-level results, the baseline \amowaspij configuration remains competitive with plain \textsc{wasp} and \textsc{clingo} on WGC, KS, TSP, and CA, solving 64, 71, 131, and 306 instances, respectively. However, it performs substantially worse than \textsc{clingo} on GA and NS, where it solves only 65 and 0 instances, compared to 238 and 4 solved by \textsc{clingo}.
The improved configurations outperform or match \textsc{clingo} on every benchmark except NS. On WGC, \amoclingoinfmrlazy solves 183 instances, 123 more than \textsc{clingo}, whereas \eoclingoinfrlazy solves 172 instances, 77 instances more than \clingo.
On GA, \amoclingoinfmrlazy obtains the best \textsc{amo} result with 314 solved instances (76 more than \textsc{clingo}). In the \textsc{eo} setting, \eoclingoinfrlazy and \textsc{clingo} both solve the complete set of 350 instances.
On KS, \amoclingoinfrlazy and \eoclingoinfrlazy both solve 77 instances, exceeding \textsc{clingo} by 29 and 30 instances in the \textsc{amo} and \textsc{eo} settings, respectively. 
Substantial gains are also observed on TSP: the best lazy configurations solve 275 \textsc{amo} and 274 \textsc{eo} instances, compared to 128 and 129 solved by \textsc{clingo}. 
On CA, which is evaluated only in the \textsc{amo} setting, \amoclingoinfmrlazy solves 339 instances, 87 more than \textsc{clingo}. 
Finally, NS remains the only benchmark on which plain \textsc{clingo} performs best, solving 4 out of 5 instances in both settings, compared to at most 3 solved by the best proposed lazy configurations.}

\begin{figure}[t!]
\centering
\begin{tikzpicture}[scale=1]
    \pgfkeys{/pgf/number format/set thousands separator = {}}
    \begin{axis}[
    scale only axis
    , xlabel={Solved instances}
    , ylabel={Time (s)}    
    , xmin=0, 
    , ymin=0, ymax=1220
    , ytick distance=200,
    , legend style={ at={(0.5,-0.25)}, anchor=north, fill=none}
    , legend columns=2
    % , legend columns=1
    , width=0.80\textwidth
    , height=0.40\textwidth
    , major tick length=2pt
    , title= {}
    ]        
    \addplot [mark size=2pt, color=amoclingo, mark=o] [unbounded coords=jump] table[col sep=comma, y index=3] {./plots/catcus_amo.csv};
    \addlegendentry{\amoclingoinfmrlazy}
    \addplot [mark size=2pt, color=amowasp, mark=o] [unbounded coords=jump] table[col sep=comma, y index=13] {./plots/catcus_amo.csv};
    \addlegendentry{\amowaspij} 
    \addplot [mark size=2pt, color=clingo, mark=+] [unbounded coords=jump] table[col sep=comma, y index=14] {./plots/catcus_amo.csv};
    \addlegendentry{\clingo}    
    \addplot [mark size=2pt, color=wasp, mark=+] [unbounded coords=jump] table[col sep=comma, y index=15] {./plots/catcus_amo.csv};
    \addlegendentry{\textsc{wasp}}            
        \end{axis}
    \end{tikzpicture}%    
    \hfill    
    \caption{
        Number of solved instances ($x$-axis) within a time limit ($y$-axis) in the \textsc{amo} setting.
    }\label{fig:catcusamo}          
    \centering
    \begin{tikzpicture}[scale=1]
        \pgfkeys{/pgf/number format/set thousands separator = {}}
    \begin{axis}[
    scale only axis
    , xlabel={Solved instances}
    , ylabel={Time (s)}    
    , xmin=0, 
    , ymin=0, ymax=1220
    , ytick distance=200,
    , legend style={ at={(0.5,-0.25)}, anchor=north, fill=none}
    , legend columns=2
    % , legend columns=1
    , width=0.80\textwidth
    , height=0.40\textwidth
    , major tick length=2pt
    , title= {}
    ]      
            \addplot [mark size=2pt, color=amoclingo, mark=o] [unbounded coords=jump] table[col sep=comma, y index=8] {./plots/catcus_eo.csv};
            \addlegendentry{\eoclingoinfrlazy}            
            \addplot [mark size=2pt, color=amowasp, mark=o] [unbounded coords=jump] table[col sep=comma, y index=15] {./plots/catcus_eo.csv};
            \addlegendentry{\eowaspij}
            \addplot [mark size=2pt, color=clingo, mark=+] [unbounded coords=jump] table[col sep=comma, y index=12] {./plots/catcus_eo.csv};
            \addlegendentry{\textsc{clingo}}                    
            \addplot [mark size=2pt, color=wasp, mark=+] [unbounded coords=jump] table[col sep=comma, y index=16] {./plots/catcus_eo.csv};
            \addlegendentry{\textsc{wasp}}
        \end{axis}
    \end{tikzpicture}%
    \caption{
        Number of solved instances ($x$-axis) within a time limit ($y$-axis) in the \textsc{eo} setting.
    }\label{fig:catcuseo}
\end{figure}

Finally, we also show the evolution of solving times across the benchmark set
%Specifically, 
in Figures~\ref{fig:catcusamo} and~\ref{fig:catcuseo}.
They report the cactus plots for the \textsc{amo} and \textsc{eo} benchmarks, respectively.  
We recall that a cactus plot is a standard way to visualize solver performance: each point $(i, j)$ on a solver's curve indicates that the $i$-th instance was solved in $j$ seconds, where instances are ordered by increasing solution time.  
Hence, solvers that appear lower and further to the right in the plot are both faster and able to solve a larger number of instances.

It is evident that \textcolor{blue}{\amoclingoinfmrlazy and \eoclingoinfrlazy} scale significantly better than the plain \textsc{clingo} and \textsc{wasp} solvers.  
In particular, their curves remain consistently below and to the right of the ones of \textsc{clingo} and \textsc{wasp}, indicating faster solving times and improved robustness across a wider range of instances.

{\color{blue}
\subsection{Impact of \textsc{amosum} and \textsc{eosum} propagator with $\le$ constraint}
Table~\ref{tab:le} shows the impact of enabling dedicated propagation for constraints with a $\le$ bound. We restrict this analysis to the GA, KS, and NS benchmarks as they are the ones with at least one suitable $\le$ constraint, which comprise 455 instances overall. For each \textsc{amo} and \textsc{eo} configuration, we compare the standard lazy version with its counterpart extended to handle $\le$ constraints.
It is worth noting that, in the \textsc{amo} case, the propagator does not propagate more than the standard one, given that the minimal possible sum coincides with the sum of true literals.
As expected, in the \textsc{amo} setting, enabling propagation for $\le$ constraints does not improve performance. 
However, the resulting performance degradation is modest.
Indeed, for \amoclingoinfrlazy, the number of solved instances decreases from 393 to 390, with two fewer instances solved on KS and one fewer on NS. Similarly, \amoclingoinfmrlazy decreases from 390 to 388 solved instances, due to a reduction of two instances on GA. This result is consistent with the theoretical observation that combining an at-most-one constraint with an upper-bound sum does not strengthen propagation beyond that already provided by the standard \textsc{sum} propagator.

The results are more heterogeneous in the \textsc{eo} setting. Enabling $\le$ propagation for \eoclingoinfrlazy has an almost neutral effect, decreasing the total from 430 to 429 solved instances. By contrast, the $\le$ variant of \eoclingoinfmrlazy improves the total number of solved instances from 427 to 436, obtaining the best overall result in this comparison. This improvement is entirely driven by KS, where the number of solved instances increases from 75 to 86. On GA and NS, the same configuration solves one fewer instance in each benchmark.

Overall, dedicated propagation for $\le$ constraints provides a limited overhead in the \textsc{amo} setting, while its effectiveness in the \textsc{eo} setting depends on the underlying reason-generation strategy and the benchmark. In particular, its combination with on-the-fly minimal reasons yields a substantial improvement on KS, although this benefit does not generalize to GA and NS.
}

\begin{table}[t!]
  \centering    
  \small
    \begin{tabular}{llrrrr}
    \toprule
    &  & \textbf{GA} & \textbf{KS} & \textbf{NS} & \textbf{Total}\\
    \midrule
    & \#Instances & 350 & 100 & 5 & 455\\
    \midrule
    \multirow{4}{*}{\rotatebox{90}{\textsc{amo}}}
& \amoclingoinfrlazy & 313 & \textbf{77} & \textbf{3} & \textbf{393}\\
& \amoclingoinfrlazy $\le$ & 313 & 75 & 2 & 390\\
& \amoclingoinfmrlazy & \textbf{314} & 75 & 1 & 390\\
 & \amoclingoinfmrlazy $\le$& 312 & 75 & 1 & 388\\
\midrule
\multirow{4}{*}{\rotatebox{90}{\textsc{eo}}} 
& \eoclingoinfrlazy & \textbf{350} & 77 & \textbf{3} & 430\\
& \eoclingoinfrlazy $\le$ & \textbf{350} & 77 & 2 & 429\\
& \eoclingoinfmrlazy & \textbf{350} & 75 & 2 & 427\\
& \eoclingoinfmrlazy $\le$ & 349 & \textbf{86} & 1 & \textbf{436}\\
    \bottomrule
    \end{tabular}
\caption{Impact of enabling propagation for $\le$ constraints on the lazy \textsc{amo} and \textsc{eo} configurations over the GA, KS, and NS benchmarks. The best result for each benchmark and setting is highlighted in bold.}
\label{tab:le}
\end{table}

\subsection{Summary of the Experimental Analysis}
In this section, we summarize the answers to Research Questions \textit{\textcolor{blue}{Q1, Q3,} Q4--Q6} and provide an overall assessment based on the results presented throughout Section~\ref{sec:exp}.

{\color{blue}
\begin{itemize}
\item[\textit{Q1}] \emph{How can inference and learning be improved in the solver?}
Section~\ref{sec:amosum} addresses this question from a theoretical perspective by introducing enhanced inference and learning mechanisms. The experimental evaluation further demonstrates that these theoretical advances translate into tangible performance improvements in practice.

\item[\textit{Q3}] \emph{Which combinations of \textsc{sum} and \textsc{amo} constraints can be effectively propagated jointly?}
Table~\ref{tab:le} helps us answer this question, showing that the only combination that fails to yield an improvement is the AMO constraint paired with the $\le$ relation. Indeed, as discussed, since the minimum possible sum coincides with the classical one in this case, it does not infer more literals than~\eqref{eq:sum:propagation}. All other combinations, namely AMO with $\ge$, and EO with both $\le$ and $\ge$, proved effective in practice.

\item[\textit{Q4}] \emph{Can reasons be minimized to improve learning and efficiency?}\\
To answer this question, we investigated on-the-fly reason minimization through Algorithm~\ref{alg:minfly} (\textit{MRF}), leading to the \amoclingoinfmr\ and \eoclingoinfmr\ configurations. The results in Table~\ref{tab:minimal:reason} show that the impact of subset-minimal reasons depends on the constraint type.

In the \textsc{amo} setting, \amoclingoinfmr\ solves 1102 instances, 14 more than \amoclingoinfr, corresponding to an improvement of approximately 1.3\%. At the benchmark level, minimization improves performance on WGC and CA, with 5 and 15 additional instances solved, respectively. These gains compensate for small decreases on GA, KS, and NS, while performance on TSP remains unchanged. The positive impact is also preserved in the fully lazy setting, where \amoclingoinfmrlazy\ solves 1187 instances, 25 more than \amoclingoinfrlazy.

In the \textsc{eo} setting, reason minimization is less effective. The number of solved instances decreases from 841 with \eoclingoinfr\ to 822 with \eoclingoinfmr, corresponding to a loss of 19 instances. A similar trend is observed for the fully lazy variants, where \eoclingoinfmrlazy\ solves 850 instances compared to 876 solved by \eoclingoinfrlazy. These results indicate that smaller reasons do not always translate into better performance, as their interaction with the solver's branching and learning heuristics differs between the \textsc{amo} and \textsc{eo} settings. Overall, on-the-fly minimization is beneficial for \textsc{amo}, whereas the non-minimized reason remains preferable for \textsc{eo}.

\item[\textit{Q5}] \emph{How does the approach generalize to exactly-one constraints and upper bounds?}\\
The experimental results show that the proposed approach generalizes effectively to exactly-one constraints. On the benchmarks shared by the two settings, the best \textsc{eo} configuration, \eoclingoinfrlazy, substantially outperforms plain \textsc{clingo} on WGC, KS, and TSP, solving 77, 30, and 145 additional instances, respectively. It matches \textsc{clingo} on all 350 GA instances and solves one fewer instance on NS. Overall, \eoclingoinfrlazy\ solves 876 out of the 1055 instances evaluated in the \textsc{eo} setting, compared to 625 solved by plain \textsc{clingo}.

The results for upper-bound constraints are more configuration-dependent. As reported in Table~\ref{tab:le}, enabling dedicated propagation for $\le$ constraints does not improve the \textsc{amo} configurations. This is consistent with the theoretical observation that combining an at-most-one constraint with an upper-bound sum does not strengthen propagation beyond the standard SUM propagator. In the \textsc{eo} setting, enabling $\le$ propagation has an almost neutral effect on \eoclingoinfrlazy, which decreases from 430 to 429 solved instances on GA, KS, and NS. By contrast, the $\le$ variant of \eoclingoinfmrlazy\ improves from 427 to 436 solved instances, primarily because of the increase from 75 to 86 solved KS instances.

These results indicate that the approach extends successfully to exactly-one constraints, while the benefit of dedicated upper-bound propagation depends on both the reason-generation strategy and the benchmark.

\item[\textit{Q6}] \emph{How effective is the approach in practice?}\\
The results across all benchmarks confirm that the proposed propagators are highly effective in practice. Both non-lazy and fully lazy configurations substantially outperform the corresponding plain solvers, with the fully lazy variants achieving the best overall results.

In the \textsc{amo} setting, \amoclingoinfmrlazy\ obtains the strongest performance, solving 1187 out of 1395 instances. This corresponds to 457 additional instances compared to plain \textsc{clingo}, which solves 730 instances, representing a 62.6\% improvement. The best non-lazy configuration, \amoclingoinfmr, also provides a substantial gain, solving 1102 instances.

In the \textsc{eo} setting, the best result is 876 out of 1055 applicable instances, obtained by \eoclingoinfrlazy. This is 251 more than plain \textsc{clingo}, corresponding to a 40.2\% improvement. The best non-lazy configuration is \eoclingoinfr, with 841 solved instances, and it also significantly outperforms the plain solver.

The cactus plots in Figures~\ref{fig:catcusamo} and~\ref{fig:catcuseo} further show that the proposed configurations not only solve more instances but also exhibit better runtime scalability. The improvements are observed across heterogeneous benchmarks, with NS being the only benchmark on which plain \textsc{clingo} obtains the best performance.

\end{itemize}
}

\paragraph{Overall Summary}
{\color{blue}
The experimental evaluation demonstrates that the proposed propagators substantially improve the practical handling of \amosum\ and \eosum constraints. Enhanced inference and propagation activation provide consistent gains over the baseline implementations, while the impact of reason minimization depends on the constraint type: it improves the \textsc{amo} configurations but is less effective for \textsc{eo}. Fully lazy propagation provides the strongest overall improvement, increasing the number of solved instances from 1102 to 1187 for the best \textsc{amo} configurations and from 841 to 876 for the best \textsc{eo} configurations.

Propagation for $\le$ constraints has a limited impact on \textsc{amo}, as expected from the theoretical analysis, but can improve \textsc{eo} performance when combined with on-the-fly minimal reasons, particularly on KS. Compared to the previous baseline and the plain \textsc{clingo} and \textsc{wasp} solvers, the best configurations solve hundreds of additional instances and exhibit improved runtime scalability. Overall, the results show that the proposed approach is effective across different constraint types, solver configurations, and benchmark domains, although the most effective reason-generation and propagation strategies remain dependent on the specific setting.
}

\section{Related Work}\label{sec:rw}

\textsc{amosum} constraints, as introduced in this work, extend the (propositional) language of Answer Set Programming (ASP). However, their applicability is not limited to ASP:
they can, in principle, be integrated into any logic-based formalism that extends propositional logic, including \emph{Satisfiability Modulo Theories (SMT)}~\citep{DBLP:journals/jacm/NieuwenhuisOT06} and \emph{Constraint Satisfaction Problems (CSP)}~\citep{DBLP:journals/eor/BrailsfordPS99,Rossi2006Handbook}.
This generality opens the door to a broader adoption of \textsc{amosum} constraints across a variety of automated reasoning paradigms.

From an implementation standpoint, two main strategies in ASP have emerged for enhancing the reasoning capabilities of logic-based languages: 
\emph{propagator-based} and \emph{translation-based} approaches. 
We discuss these in turn.


\subsection{Propagator-Based Approaches}

Propagator-based approaches extend the inference engine of a solver by introducing custom mechanisms that infer additional literals from partial assignments.
These are implemented as procedural modules, which the solver invokes during unit propagation. 
The \textsc{amosum} constraints introduced in this paper are supported via such a propagator, described in Section~\ref{sec:amo:propagator}.

Several ASP solvers, including \textsc{clingo}~\citep{DBLP:journals/ai/GebserKS12,DBLP:journals/tplp/GebserKKS19}, \linebreak \textsc{wasp}~\citep{DBLP:journals/tplp/AlvianoDM18}, and \textsc{idp}~\citep{Denecker2010DPLLAggAE,DBLP:journals/ngc/0001JCJBD16}, adopt propagator-based mechanisms to handle aggregates and other advanced constructs.
\textcolor{blue}{In particular, \textsc{clingo} adopts a hybrid approach in which small aggregates are compiled into rules, whereas more complex ones, as can occur for complex \textsc{amo} constraints, are handled by dedicated internal C-level propagators~\citep{DBLP:conf/iclp/GebserKKS09}.}
The threshold controlling this behavior is adjustable via command-line parameters.
Additionally, both \textsc{clingo} and \textsc{wasp} expose Python APIs that allow users to define external propagators, making the integration of new constraints modular and accessible~\citep{DBLP:journals/algorithms/CabalarFSW23,DBLP:journals/tplp/DodaroR20}.

The propagator for \textsc{amosum} constraints proposed in this work builds upon our previous version~\citep{DBLP:conf/ijcai/AlvianoDFM24}.
That earlier version focused primarily on the inference rule and its application within \textsc{wasp}. 
In contrast, the present paper introduces several significant improvements: 
a refined inference strategy that minimizes reasons both statically and on-the-fly, an ASP language construct to express \textsc{amosum} constraints declaratively, and a more flexible integration mechanism. 
Furthermore, while the previous version was implemented as a Python external propagator for \textsc{wasp}, the current version has also been ported to \textsc{clingo}, as C and Python external propagators.
These implementations enable a more direct comparison with other propagator-based approaches and demonstrate the portability and generality of the proposed design.


\subsection{Translation-Based Approaches}
\textcolor{blue}{Translation-based approaches instead compile high-level constructs into equivalent solver-native representations, such as normal ASP rules or CNF formulas. 
This strategy is often applied to aggregates, cardinality constraints, and pseudo-Boolean (PB) expressions.}

In the context of ASP, several encodings have been proposed that translate PB constraints into SAT or ASP-friendly forms.
These include adder circuits and binary decision diagrams (BDDs)~\citep{DBLP:journals/jair/AbioNORM12,DBLP:journals/jsat/EenS06}, as well as sorting networks and watchdog encodings~\citep{DBLP:conf/sat/BailleuxBR09}. 
The ASP tools \textsc{lp2sat} and \textsc{lp2normal} incorporate many such translations~\citep{DBLP:conf/jelia/BomansonGJ14,DBLP:conf/lpnmr/BomansonJ13}, producing CNF and normal logic programs respectively. 

Similarly, \emph{asp-fzn}~\citep{DBLP:journals/tplp/EiterGMOK25} translates ground (constraint) answer set programs into the intermediate format \emph{FlatZinc}, which can in turn be given to Constraint or Mixed Integer Programming solvers. Aggregates are translated into linear constraints and the respective solvers generally use similar propagation techniques as described in this work. However, to the best of our knowledge, none provide a propagator for \textsc{AMOSUM}.
Another example is \textsc{cmodels}~\citep{DBLP:journals/jar/GiunchigliaLM06}, which compiles aggregates into nested logic programs~\citep{DBLP:journals/tplp/FerrarisL05}.

%% Added
{\color{blue}
In the SAT literature, several approaches exploit AMO relations when compiling PB constraints into CNF. \cite{DBLP:conf/ijcai/BofillCSV17} proposed a Multi-Decision Diagram (MDD)-based encoding of PB constraints subject to an upper bound. When AMO relations are known over disjoint subsets of variables, combinations that violate these relations can be pruned from the underlying decision diagram. This allows the constraint to be represented using a MDD, rather than a Binary Decision Diagram, resulting in a more compact encoding and improved solving performance.

Building on this idea, \cite{DBLP:journals/ai/BofillCNSUV22} introduced AMO-aware generalizations of several established PB encodings for the \textit{less-than-or-equal} case, including the Sequential Weight Counter, Generalized Totalizer, \(n\)-Level Modulo Totalizer, Global Polynomial Watchdog, and Local Polynomial Watchdog. Their results demonstrate that exploiting AMO relations can substantially reduce encoding size and improve solving time. A related contribution is due to \cite{DBLP:conf/cp/AnsoteguiBCDEMN19}, who proposed the automatic detection of AMO and Exactly-One (EO) relations and exploited the detected structure to improve MDD-based SAT encodings of PB constraints.

Our contribution is complementary to these translation-based approaches. Whereas the aforementioned methods exploit AMO and EO relations to produce smaller CNF encodings and thereby reduce propagation overhead, we retain the structured constraint and use these relations directly during propagation. In particular, the proposed propagators compute a tighter maximum attainable sum under a partial interpretation, allowing them to derive consequences that may be missed when the PB and AMO constraints are propagated independently.

Although our implementation targets ASP solvers, the underlying propagation principles are not inherently ASP-specific. A native PB solver that represents PB constraints directly, rather than translating them into SAT, could associate known or automatically detected AMO and EO groups with a weighted constraint and incorporate the same maximum-possible-sum tests and reason-minimization techniques into its propagation engine. The ASP implementations presented in this work can thus be seen as a concrete instantiation of a more general approach to native PB propagation.

For AMO constraints specifically, several SAT encodings are available, including pairwise (binomial), binary (bitwise), commander, product, sequential-counter, and bimander encodings. These encodings offer different trade-offs among representation size, propagation strength, and ease of integration. A recent comparison is provided by \cite{DBLP:conf/sma2/NguyenNKB20}.

Finally, the proposed \textsc{amosum} propagator is also compatible with translation-based encodings. For example, the consequences of the initial inference rule described in Section~\ref{sec:amo:propagator} could be precompiled into clauses, while the remaining inferences are handled by the dedicated propagator. Such a hybrid configuration would combine the low-level efficiency of established encodings with the stronger propagation enabled by preserving the \textsc{amosum} structure.
}

\subsection{Differences with~\cite{DBLP:conf/ijcai/AlvianoDFM24}}

\textcolor{blue}{As discussed in Sections~\ref{sub:sec:prop:act}--\ref{sec:exp}, this work extends the propagator proposed by~\cite{DBLP:conf/ijcai/AlvianoDFM24} with several improvements.
The first improvement concerns the programming language and underlying solver used. Since \textsc{wasp} only provides a Python interface for implementing custom propagators, we decided to switch to \textsc{clingo}, both to exploit a more efficient solver and to use a compiled language that avoids unnecessary overhead.
We also improved propagation activation, as discussed in Section~\ref{sub:sec:prop:act}.
In addition, we introduced a new reason. More specifically, the reason proposed by~\cite{DBLP:conf/ijcai/AlvianoDFM24} was larger, leading to unnecessary literals within it.
Given a literal $\ell \in \lits_\sigma$ propagated by inference rule (RT)~\eqref{eq:amosum:propagation:true}, the old reason was
$$\mathit{reason_{old}}(\ell) = \lits_\sigma|_{\mathit{part}_\sigma(\ell)} \cup \bigcup_{s \in \mathit{parts}_\sigma \setminus \{\mathit{part}_\sigma(\ell)\}} \mathit{rsn}_\sigma(J,\ell,s),$$
whereas our new reason satisfies $\mathit{reason}(\ell) \subseteq \mathit{reason_{old}}(\ell)$.
Regarding minimization and laziness: the previous version did not define either of these techniques, so they are presented here as entirely novel.}

\section{Conclusion}\label{sec:conclusion}

%\textsc{amosum} constraints extend the language of Answer Set Programming (ASP) with a construct that is designed not to increase the expressive power of the language, but rather to enhance the \emph{search-space pruning capabilities} of solvers.
%The construct provides a compact representation of groups of constraints that are already expressible in ASP, such as a single weighted sum constraint accompanied by multiple at-most-one constraints.
%By combining the semantic information of these components, \textsc{amosum} constraints enable more effective propagation and allow solvers to deduce deterministic consequences that would not be apparent when reasoning over the individual constraints in isolation.

%This work builds upon 
%%our previous formulation 
%a preliminary version presented %at the IJCAI 2024 conference 
%in IJCAI 2024 \citep{DBLP:conf/ijcai/AlvianoDFM24} and introduces several key extensions.
%Notably, we present improved propagation algorithms, including a novel on-the-fly minimization strategy, and an ASP language construct that allows \textsc{amosum} constraints to be expressed declaratively.
%Furthermore, we provide two separate implementations: one as an external Python propagator for \textsc{wasp}, and another as a C extension for \textsc{clingo}. These implementations enable broader adoption and facilitate direct comparison with existing solver architectures.

In this paper, we have presented \textsc{amosum} constraints, which are a novel construct for the Answer Set Programming (ASP) language that allows one to compactly express groups of constraints in a declarative manner, such as a single weighted sum constraint accompanied by multiple at-most-one constraints. While the construct does not increase the expressive power of the language, it provides the user with further means to guide \emph{search-space pruning} of
solvers. By combining the semantic information of their constituents, \textsc{amosum} constraints allow for more effective propagation and enable solvers to deduce deterministic consequences that would not be apparent when reasoning over the individual constraints in isolation.
Our work significantly enhances and further develops a preliminary work presented at IJCAI \citep{DBLP:conf/ijcai/AlvianoDFM24} by introducing several key extensions, among them improved propagation algorithms, a novel  strategy for on-the-fly minimization, and the \textsc{amosum} construct for the ASP user language. Furthermore, we provide three implementations: an external Python propagator for \textsc{wasp}, and both Python and C integrations for \textsc{clingo}.
The implementations enable broader adoption and facilitate direct comparison with existing solver architectures.

The empirical evaluation confirms the potential of \textsc{amosum} constraints to deliver substantial performance gains.
Compared to both the baseline and the previously introduced version, the new implementations yield faster solving times and reduced search efforts on a range of benchmark instances.
In particular, the integration into \textsc{clingo} demonstrates the portability and efficiency of the proposed techniques, and the declarative ASP construct provides an intuitive interface for users.
\smallskip

\noindent{\em Outlook.} There are different avenues for future work. 
%An avenue for future work 
One concerns the generalization of \textsc{amosum} constraints. While the current construct is tailored to capture the interaction between a weighted sum and multiple at-most-one constraints, other combinations of constraints may similarly benefit from integrated propagation.
Identifying such patterns and designing corresponding propagators could further enhance solver performance across a broader range of problems.

\textcolor{blue}{A further direction is to extend \textsc{amosum} from a constraint construct to a first-class aggregate atom that can occur in arbitrary rule bodies, including the study of recursive occurrences under stable-model semantics.}

Another promising direction for further research is the application of \textsc{amosum} constraints to practical encodings. 
In this work, we have already demonstrated their effectiveness in the context of nurse scheduling, a domain where weighted preferences and mutual exclusions frequently interact~\citep{DBLP:conf/lpnmr/DodaroM17}. 
Further investigation may reveal additional domains where \textsc{amosum}-based reasoning improves both expressiveness and performance.

All code, encodings, and data used in our experiments are publicly available at \url{https://github.com/instafiore/AMOSUM/tree/amosum}, ensuring full reproducibility\footnote{The IJCAI benchmarks are available at
\url{https://doi.org/10.5281/zenodo.11115982}~\citep{alviano_2024_11115982}. Upon acceptance, this record will be updated to include the extended benchmark set.}.




%{\color{red}SIAMO A PAGINA \thepage: limite 7+2}

%\clearpage
% \section*{Ethical Statement}
% There are no ethical issues.

\section*{Acknowledgments}
This work was supported by the Italian Ministry of University and Research (MUR) 
    under PRIN project PRODE ``Probabilistic declarative process mining'', CUP H53D23003420006,
    under PNRR project FAIR ``Future AI Research'', CUP H23C22000860006,
    under PNRR project Tech4You ``Technologies for climate change adaptation and quality of life improvement'', CUP H23C22000370006, and~
    under PNRR project SERICS ``SEcurity and RIghts in the CyberSpace'', CUP H73C22000880001;
by the Italian Ministry of Health (MSAL)
    under POS projects CAL.HUB.RIA (CUP H53C22000800006) and RADIOAMICA (CUP H53C22000650006);
by the Italian Ministry of Enterprises and Made in Italy
    under project STROKE 5.0 (CUP B29J23000430005);
    under PN RIC project ASVIN ``Assistente Virtuale Intelligente di Negozio'' (CUP B29J24000200005);
by Regione Calabria
    under PR CALABRIA FESR FSE 2021–2027 projects
    IOS4OS ``Intelligent Optimization System for Optimization Solutions'' (CUP J89I24002920005);
    NextGenGuides ``Innovazione Tecnologica per le Guide Turistiche del Futuro tramite l'ausilio di tecnologie innovative AI e sistemi di geolocalizzazione avanzata'' (CUP J39I24002040005);
    and MOZART ``Modelli e tecniche di mOdernizzazione di applicaZioni e data silos a supporto di clienti esterni, field force e impiegati di backoffice basati su AI GeneRativa e apprendimento auTomatico'' (CUP J49I24001740005);
and by the LAIA lab (part of the SILA labs). 
Mario Alviano is a member of Gruppo Nazionale Calcolo Scientifico-Istituto Nazionale di Alta Matematica (GNCS-INdAM).

    \section*{Declaration of generative AI and AI-assisted technologies}
    During the preparation of this work the author(s) used ChatGPT to check for grammatical errors and improve the readability of some sentences. After using this tool/service, the author(s) reviewed and edited the content as needed and take(s) full responsibility for the content of the published article.

\bibliographystyle{elsarticle-harv}
\bibliography{bibl}
\end{document}
